\input{preamble}

% OK, commencez ici.
%
\begin{document}

\title{Catégories dérivées des variétés}


\maketitle

\phantomsection
\label{section-phantom}

\tableofcontents

\section{Introduction}
\label{section-introduction}

\noindent
Dans ce chapitre, nous poursuivons la discussion entamée dans
Catégories dérivées des schémas, section \ref{perfect-section-introduction}.
Nous étudierons les transformations de Fourier-Mukai,
introduites par Mukai dans \cite{Mukai}.
Nous démontrerons le théorème d'Orlov sur les équivalences dérivées (\cite{Orlov-K3}).
Nous étudierons également la dénombrabilité des classes d'équivalence dérivée,
établie par Anel et Toën dans \cite{AT}.

\medskip\noindent
Le livre de Daniel Huybrechts \cite{Huybrechts} constitue une bonne
introduction à ce sujet. Parmi les autres articles qui ont contribué
à populariser ce domaine, citons
\begin{enumerate}
\item l'article de Bondal et Kapranov, voir \cite{Bondal-Kapranov},
\item l'article de Bondal et Orlov, voir \cite{Bondal-Orlov},
\item l'article de Bondal et Van den Bergh, voir \cite{BvdB},
\item les articles de Beilinson, voir
\cite{Beilinson} et \cite{Beilinson-derived},
\item l'article d'Orlov, voir \cite{Orlov-AV},
\item l'article d'Orlov, voir \cite{Orlov-motives},
\item l'article de Rouquier, voir \cite{Rouquier-dimensions},
\item et bien d'autres encore pourraient être cités ici.
\end{enumerate}




\section{Conventions et notations}
\label{section-conventions}

\noindent
Soit $k$ un corps. Une catégorie triangulée $k$-linéaire $\mathcal{T}$
est une catégorie triangulée (Catégories dérivées, section
\ref{derived-section-triangulated-definitions})
munie d'une structure $k$-linéaire
(Algèbre différentielle graduée, section \ref{dga-section-linear})
telle que les foncteurs de translation $[n] : \mathcal{T} \to \mathcal{T}$
soient $k$-linéaires pour tout $n \in \mathbf{Z}$.

\medskip\noindent
Soit $k$ un corps. Nous notons $\text{Vect}_k$ la catégorie des
espaces vectoriels sur $k$. Pour un espace vectoriel sur $k$, noté $V$, nous notons
$V^\vee$ le dual $k$-linéaire de $V$, c'est-à-dire $V^\vee = \Hom_k(V, k)$.

\medskip\noindent
Soit $X$ un schéma. Nous notons $D_{perf}(\mathcal{O}_X)$ la
sous-catégorie pleine de $D(\mathcal{O}_X)$ formée des complexes parfaits
(Cohomologie, section \ref{cohomology-section-perfect}).
Si $X$ est noethérien, alors
$D_{perf}(\mathcal{O}_X) \subset D^b_{\textit{Coh}}(\mathcal{O}_X)$, voir
Catégories dérivées des schémas, lemme \ref{perfect-lemma-perfect-on-noetherian}.
Si $X$ est noethérien et régulier, alors
$D_{perf}(\mathcal{O}_X) = D^b_{\textit{Coh}}(\mathcal{O}_X)$, voir
Catégories dérivées des schémas, lemme \ref{perfect-lemma-perfect-on-regular}.

\medskip\noindent
Soit $k$ un corps. Soient $X$ et $Y$ des schémas sur $k$. Dans cette
situation, nous écrirons $X \times Y$ au lieu de $X \times_{\Spec(k)} Y$.


\medskip\noindent
Soit $S$ un schéma. Soient $X$ et $Y$ des schémas sur $S$.
Soit $\mathcal{F}$ un $\mathcal{O}_X$-module et soit
$\mathcal{G}$ un $\mathcal{O}_Y$-module. Nous posons
$$
\mathcal{F} \boxtimes \mathcal{G} =
\text{pr}_1^*\mathcal{F} \otimes_{\mathcal{O}_{X \times_S Y}}
\text{pr}_2^*\mathcal{G}
$$
comme $\mathcal{O}_{X \times_S Y}$-module.
Si $K \in D(\mathcal{O}_X)$ et $M \in D(\mathcal{O}_Y)$, nous posons
$$
K \boxtimes M =
L\text{pr}_1^*K \otimes_{\mathcal{O}_{X \times_S Y}}^\mathbf{L} L\text{pr}_2^*M
$$
comme objet de $D(\mathcal{O}_{X \times_S Y})$.
Notre notation est donc potentiellement ambiguë, mais le contexte devrait
indiquer clairement laquelle des deux significations est visée.





\section{Foncteurs de Serre}
\label{section-Serre-functors}

\noindent
Le contenu de cette section est repris de \cite{Bondal-Kapranov}.

\begin{lemma}
\label{lemma-Serre-functor-exists}
Soit $k$ un corps. Soit $\mathcal{T}$ une catégorie triangulée
$k$-linéaire telle que $\dim_k \Hom_\mathcal{T}(X, Y) < \infty$
pour tous $X, Y \in \Ob(\mathcal{T})$. Les assertions suivantes sont équivalentes :
\begin{enumerate}
\item il existe une équivalence $k$-linéaire
$S : \mathcal{T} \to \mathcal{T}$ et des isomorphismes $k$-linéaires
$c_{X, Y} : \Hom_\mathcal{T}(X, Y) \to \Hom_\mathcal{T}(Y, S(X))^\vee$
fonctoriels en $X, Y \in \Ob(\mathcal{T})$ ;
\item pour tout $X \in \Ob(\mathcal{T})$,
le foncteur $Y \mapsto \Hom_\mathcal{T}(X, Y)^\vee$
est représentable et le foncteur $Y \mapsto \Hom_\mathcal{T}(Y, X)^\vee$
est coreprésentable.
\end{enumerate}
\end{lemma}

\begin{proof}
La condition (1) entraîne (2) car, étant donnés $(S, c)$ et $X \in \Ob(\mathcal{T})$,
l'objet $S(X)$ représente le foncteur
$Y \mapsto \Hom_\mathcal{T}(X, Y)^\vee$ et l'objet $S^{-1}(X)$ coreprésente
le foncteur $Y \mapsto \Hom_\mathcal{T}(Y, X)^\vee$.

\medskip\noindent
Supposons (2). Nous utiliserons à plusieurs reprises le lemme de Yoneda, voir
Catégories, lemme \ref{categories-lemma-yoneda}.
Pour tout $X$, notons $S(X)$ l'objet qui représente le
foncteur $Y \mapsto \Hom_\mathcal{T}(X, Y)^\vee$. Étant donné
$\varphi : X \to X'$, on obtient un unique morphisme $S(\varphi) : S(X) \to S(X')$
déterminé par la transformation de foncteurs correspondante
$\Hom_\mathcal{T}(X, -)^\vee \to \Hom_\mathcal{T}(X', -)^\vee$.
Ainsi, $S$ est un foncteur et l'on obtient les isomorphismes $c_{X, Y}$
par construction. Il reste à montrer que $S$ est une équivalence.
Pour tout $X$, notons $S'(X)$ l'objet qui coreprésente le
foncteur $Y \mapsto \Hom_\mathcal{T}(Y, X)^\vee$. En raisonnant comme
ci-dessus, on voit que $S'$ est un foncteur. Nous affirmons que $S'$
est un quasi-inverse de $S$. En effet, observons que
$$
\Hom_\mathcal{T}(X, Y) = \Hom_\mathcal{T}(Y, S(X))^\vee =
\Hom_\mathcal{T}(S'(S(X)), Y)
$$
bifonctoriellement, c'est-à-dire que $S' \circ S \cong \text{id}_\mathcal{T}$.
De même, on a
$$
\Hom_\mathcal{T}(Y, X) = \Hom_\mathcal{T}(S'(X), Y)^\vee =
\Hom_\mathcal{T}(Y, S(S'(X)))
$$
et l'on obtient $S \circ S' \cong \text{id}_\mathcal{T}$.
\end{proof}

\begin{definition}
\label{definition-Serre-functor}
Soit $k$ un corps. Soit $\mathcal{T}$ une catégorie triangulée
$k$-linéaire telle que $\dim_k \Hom_\mathcal{T}(X, Y) < \infty$
pour tous $X, Y \in \Ob(\mathcal{T})$. On dit qu'{\it il existe un foncteur
de Serre} si les conditions équivalentes du lemme \ref{lemma-Serre-functor-exists}
sont satisfaites. Dans ce cas, un {\it foncteur de Serre} est une équivalence $k$-linéaire
$S : \mathcal{T} \to \mathcal{T}$ munie d'isomorphismes $k$-linéaires
$c_{X, Y} : \Hom_\mathcal{T}(X, Y) \to \Hom_\mathcal{T}(Y, S(X))^\vee$
fonctoriels en $X, Y \in \Ob(\mathcal{T})$.
\end{definition}

\begin{remark}
\label{remark-matress}
Soient $X^0 \to X^1 \to X^2 \to X^0[1]$ et $Y^0 \to Y^1 \to Y^2 \to Y^0[1]$
des triangles distingués dans une catégorie triangulée. Pour
$p \in \mathbf{Z}$, écrivons $p = 3n + i$ avec $i \in \{0, 1, 2\}$ et
posons $X^p = X^i[n]$. De même pour $Y^q$. Considérons le complexe double
dont les termes sont
$$
K^{p, q} = \Hom(X^{-p}, Y^q)
$$
La différentielle $d_1 : K^{p, q} \to K^{p + 1, q}$ est induite par le morphisme
$X^{-p - 1} \to X^{-p}$ (égal au morphisme correspondant
du premier triangle distingué, à un décalage près),
et la différentielle $d_2 : K^{p, q} \to K^{p, q + 1}$ est de même induite
par le morphisme $Y^q \to Y^{q + 1}$. D'après
Catégories dérivées, lemme \ref{derived-lemma-representable-homological},
les lignes et les colonnes de ce complexe double sont exactes.
De plus, $K^{p, q} = K^{p + 3, q - 3}$, et ces
égalités sont compatibles avec les différentielles. Enfin, l'axiome TR3
entraîne une propriété supplémentaire : étant donnés
$\alpha \in K^{p, q}$ et $\beta \in K^{p - 1, q + 1}$ tels
que $d_2 \alpha = d_1 \beta$, il existe $\gamma \in K^{p - 2, q + 2}$
tel que
$d_1 \gamma = d_2 \beta $ dans $K^{p - 1, q + 2}$ et
$d_2 \gamma = d_1 \alpha $ dans $K^{p - 2, q + 3} = K^{p + 1, q}$.
(Indication : pour $p = q = 0$, c'est exactement l'énoncé de TR3 ;
pour les autres indices, on le démontre par décalage.)
Un complexe double ayant ces propriétés est appelé {\it matress}
(voir \cite{Bondal-Kapranov}).
\end{remark}

\begin{remark}
\label{remark-dual-matress}
Soit $k$ un corps et soit $K^{\bullet, \bullet}$ un complexe double
d'espaces vectoriels de dimension finie sur $k$, qui est aussi une {\it matress} au sens de la
remarque \ref{remark-matress}. Nous affirmons que le complexe double dual
$L^{p, q} = \Hom_k(K^{-q, -p}, k)$ est lui aussi une {\it matress}. L'exactitude des
lignes et des colonnes ainsi que la 3-périodicité sont immédiates. Pour vérifier
la condition supplémentaire, considérons l'application linéaire
$$
\partial :
K^{p, q} \oplus K^{p - 1, q + 1} \oplus K^{p - 2, q + 2}
\longrightarrow
K^{p, q + 1} \oplus K^{p - 1, q + 2} \oplus K^{p - 2, q + 3}
$$
qui envoie $(\alpha, \beta, \gamma)$ sur
$(d_2 \alpha - d_1 \beta, d_2 \gamma - d_1 \alpha, d_1 \gamma - d_2 \beta)$,
avec les identifications de la remarque \ref{remark-matress}.
La condition d'être une {\it matress} s'écrit
$$
\Im(\partial) \cap
\left( 0 \oplus K^{p - 1, q + 2} \oplus K^{p - 2, q + 3} \right)
=
\Im(\partial|_{0 \oplus 0 \oplus K^{p - 2, q + 2}})
$$
En notant ${}^\wedge$ le dual d'un espace vectoriel ou d'une application, la dualisation
donne
$$
\Ker(\partial^\wedge) +
\left(L^{-p, -q - 1} \oplus 0 \oplus 0 \right)
=
\Ker(\text{pr}_{L^{-p + 2, -q - 2}} \circ \partial^\wedge)
$$
où le terme $\text{pr}$ désigne la projection. Cela implique
que
$\Im(\partial^\wedge) \cap (L^{-p, -q} \oplus L^{-p + 1, -q - 1} \oplus 0)$
est égal à
$\Im(\partial^\wedge|_{L^{-p, -q - 1} \oplus 0 \oplus 0})$, ce qui
se traduit par la condition de {\it matress} pour $L^{\bullet, \bullet}$.
Nous omettons certains détails.
\end{remark}

\begin{lemma}
\label{lemma-Serre-functor}
Dans la situation de la définition \ref{definition-Serre-functor},
s'il existe un foncteur de Serre, celui-ci est unique à isomorphisme unique près
et c'est un foncteur exact de catégories triangulées.
\end{lemma}

\begin{proof}
Étant donné un foncteur de Serre $S$, l'objet $S(X)$ représente
le foncteur $Y \mapsto \Hom_\mathcal{T}(X, Y)^\vee$.
Ainsi, l'objet $S(X)$, muni de l'identification fonctorielle
$\Hom_\mathcal{T}(X, Y)^\vee = \Hom_\mathcal{T}(Y, S(X))$,
est déterminé à isomorphisme unique près par le lemme de Yoneda
(Catégories, lemme \ref{categories-lemma-yoneda}).
De plus, pour $\varphi : X \to X'$, le morphisme $S(\varphi) : S(X) \to S(X')$
est uniquement déterminé par la transformation de foncteurs correspondante
$\Hom_\mathcal{T}(X, -)^\vee \to \Hom_\mathcal{T}(X', -)^\vee$.

\medskip\noindent
Pour des objets $X, Y$ de $\mathcal{T}$, on a
\begin{align*}
\Hom(Y, S(X)[1])^\vee
& =
\Hom(Y[-1], S(X))^\vee \\
& =
\Hom(X, Y[-1]) \\
& =
\Hom(X[1], Y) \\
& =
\Hom(Y, S(X[1]))^\vee
\end{align*}
Le lemme de Yoneda fournit donc un unique isomorphisme
$S(X[1]) \to S(X)[1]$ qui induit l'isomorphisme obtenu de la première à la dernière expression.
Comme chacun des isomorphismes ci-dessus est fonctoriel à la fois en $X$ et en $Y$,
on obtient ainsi un isomorphisme de foncteurs
$S \circ [1] \to [1] \circ S$.

\medskip\noindent
Soit $(A, B, C, f, g, h)$ un triangle distingué de $\mathcal{T}$.
Il faut montrer que le triangle $(S(A), S(B), S(C), S(f), S(g), S(h))$
est distingué. On utilise ici l'isomorphisme canonique $S(A[1]) \to S(A)[1]$
construit ci-dessus pour identifier le but $S(A[1])$ de $S(h)$ à $S(A)[1]$.
Observons d'abord que, pour tout $X$ dans $\mathcal{T}$,
le triangle $(S(A), S(B), S(C), S(f), S(g), S(h))$ induit
une suite exacte longue
$$
\ldots \to
\Hom(X, S(A)) \to
\Hom(X, S(B)) \to
\Hom(X, S(C)) \to
\Hom(X, S(A)[1]) \to \ldots
$$
d'espaces vectoriels de dimension finie sur $k$. En effet, cette suite est
le dual $k$-linéaire de la suite
$$
\ldots \leftarrow
\Hom(A, X) \leftarrow
\Hom(B, X) \leftarrow
\Hom(C, X) \leftarrow
\Hom(A[1], X) \leftarrow
\ldots
$$
qui est exacte d'après Catégories dérivées, lemme
\ref{derived-lemma-representable-homological}.
Choisissons ensuite un triangle distingué $(S(A), E, S(C), i, p, S(h))$ ;
c'est possible d'après les axiomes TR1 et TR2. Nous voulons construire la flèche
en pointillés qui rende commutatif le diagramme suivant :
$$
\xymatrix{
S(C)[-1] \ar[r]_-{S(h[-1])} &
S(A) \ar[r]_{S(f)} &
S(B) \ar[r]_{S(g)} &
S(C) \ar[r]_{S(h)} &
S(A)[1] \\
S(C)[-1] \ar[r]^-{S(h[-1])} \ar@{=}[u] &
S(A) \ar[r]^i \ar@{=}[u] &
E \ar[r]^p \ar@{..>}[u]^\varphi &
S(C) \ar[r]^{S(h)} \ar@{=}[u] &
S(A)[1] \ar@{=}[u]
}
$$
En effet, si l'on dispose de $\varphi$, alors, pour tout $X$, l'application obtenue
$\Hom(X, E) \to \Hom(X, S(B))$ est un isomorphisme d'espaces vectoriels
sur $k$. On obtient en effet un diagramme commutatif
$$
\xymatrix{
\Hom(X, S(C)[-1]) \ar[r] &
\Hom(X, S(A)) \ar[r] &
\Hom(X, S(B)) \ar[r] &
\Hom(X, S(C)) \ar[r] &
\Hom(X, S(A)[1]) \\
\Hom(X, S(C)[-1]) \ar[r] \ar@{=}[u] &
\Hom(X, S(A)) \ar[r] \ar@{=}[u] &
\Hom(X, E) \ar[r] \ar[u]^\varphi &
\Hom(X, S(C)) \ar[r] \ar@{=}[u] &
\Hom(X, S(A)[1]) \ar@{=}[u]
}
$$
dont les lignes sont exactes (voir ci-dessus), et le lemme des cinq
(Homologie, lemme \ref{homology-lemma-five-lemma}) montre
que le morphisme médian est un isomorphisme. Le lemme de Yoneda
permet alors de conclure que $\varphi$ est un isomorphisme.

\medskip\noindent
Pour construire $\varphi$, formons la {\it matress} de la remarque \ref{remark-matress}
à partir des triangles distingués
$A \to B \to C \to A[1]$ et $S(A) \to E \to S(C) \to S(A)[1]$.
D'après la remarque \ref{remark-dual-matress}, son dual $k$-linéaire
est également une {\it matress}. Comme $S$ est un foncteur de Serre, on
voit, comme ci-dessus, que cette {\it matress} duale est le complexe double formé par
les triangles (dont le second n'est pas encore connu pour être distingué)
$S(A) \to E \to S(C) \to S(A)[1]$ et $S(A) \to S(B) \to S(C) \to S(A)[1]$.
Or la condition de {\it matress} dit exactement que les morphismes
$S(A) \to S(A)$ et $S(C) \to S(C)$ s'insèrent dans un morphisme de triangles,
c'est-à-dire qu'il existe un $\varphi$ qui rend commutatif le diagramme ci-dessus.
\end{proof}






\section{Exemples de foncteurs de Serre}
\label{section-examples-Serre-functors}

\noindent
Le lemme ci-dessous est l'exemple classique.

\begin{lemma}
\label{lemma-Serre-functor-Gorenstein-proper}
Soit $k$ un corps. Soit $X$ un schéma propre sur $k$ et de Gorenstein.
Considérons le complexe $\omega_X^\bullet$ du
lemme \ref{duality-lemma-duality-proper-over-field} de Dualité pour les schémas.
Alors le foncteur
$$
S : D_{perf}(\mathcal{O}_X) \longrightarrow D_{perf}(\mathcal{O}_X),\quad
K \longmapsto S(K) = \omega_X^\bullet \otimes_{\mathcal{O}_X}^\mathbf{L} K
$$
est un foncteur de Serre.
\end{lemma}

\begin{proof}
L'énoncé a un sens puisque $\dim \Hom_X(K, L) < \infty$
pour $K, L \in D_{perf}(\mathcal{O}_X)$ d'après
Catégories dérivées des schémas, lemme \ref{perfect-lemma-ext-finite}.
Comme $X$ est de Gorenstein, le complexe dualisant $\omega_X^\bullet$
est un objet inversible de $D(\mathcal{O}_X)$, voir
Dualité pour les schémas, lemme \ref{duality-lemma-gorenstein}.
En particulier, localement sur $X$, le complexe $\omega_X^\bullet$
possède un unique faisceau de cohomologie non nul, lequel est un module inversible, voir
Cohomologie, lemme \ref{cohomology-lemma-invertible-derived}.
Ainsi, $S(K)$ appartient à $D_{perf}(\mathcal{O}_X)$.
D'autre part, l'inversibilité de $\omega_X^\bullet$
implique clairement que $S$ est une auto-équivalence de $D_{perf}(\mathcal{O}_X)$.
Enfin, il faut construire un isomorphisme
$$
c_{K, L} : \Hom_X(K, L) \longrightarrow
\Hom_X(L, \omega_X^\bullet \otimes_{\mathcal{O}_X}^\mathbf{L} K)^\vee
$$
bifonctoriel en $K, L$. Pour cela, utilisons les isomorphismes canoniques
$$
\Hom_X(K, L) = H^0(X, L \otimes_{\mathcal{O}_X}^\mathbf{L} K^\vee)
$$
et
$$
\Hom_X(L, \omega_X^\bullet \otimes_{\mathcal{O}_X}^\mathbf{L} K) =
H^0(X, 
\omega_X^\bullet \otimes_{\mathcal{O}_X}^\mathbf{L} K
\otimes_{\mathcal{O}_X}^\mathbf{L} L^\vee)
$$
donnés dans Cohomologie, lemme \ref{cohomology-lemma-dual-perfect-complex}.
Puisque $(L \otimes_{\mathcal{O}_X}^\mathbf{L} K^\vee)^\vee =
(K^\vee)^\vee \otimes_{\mathcal{O}_X}^\mathbf{L} L^\vee$
et qu'il existe un isomorphisme canonique $K \to (K^\vee)^\vee$,
ces espaces vectoriels sur $k$ sont canoniquement duaux d'après
Dualité pour les schémas, lemme
\ref{duality-lemma-duality-proper-over-field-perfect}.
On obtient ainsi les isomorphismes $c_{K, L}$.
Nous omettons la démonstration de leur fonctorialité.
\end{proof}





\section{Caractérisation des modules cohérents}
\label{section-coherent}

\noindent
Cette section prolonge, en un certain sens, la discussion menée
dans Catégories dérivées des schémas, section \ref{perfect-section-pseudo-coherent}
et Compléments sur les morphismes, section
\ref{more-morphisms-section-characterize-pseudo-coherent}.

\medskip\noindent
Avant de pouvoir énoncer le résultat, nous avons besoin de quelques notations.
Soit $k$ un corps. Soit $n \geq 0$ un entier.
Soit $S = k[X_0, \ldots, X_n]$. Pour un entier $e$, notons
$S_e \subset S$ l'espace des polynômes homogènes de degré $e$.
Considérons la $k$-algèbre (non commutative)
$$
R =
\left(
\begin{matrix}
S_0 & S_1 & S_2 & \ldots & \ldots \\
0 & S_0 & S_1 & \ldots & \ldots\\
0 & 0 & S_0 & \ldots & \ldots \\
\ldots & \ldots & \ldots & \ldots & \ldots \\
0 & \ldots & \ldots & \ldots & S_0
\end{matrix}
\right)
$$
(à $n + 1$ lignes et colonnes), munie de la multiplication et de l'addition évidentes.

\begin{lemma}
\label{lemma-perfect-for-R}
Avec $k$, $n$ et $R$ comme ci-dessus, pour un objet $K$ de $D(R)$,
les assertions suivantes sont équivalentes :
\begin{enumerate}
\item $\sum_{i \in \mathbf{Z}} \dim_k H^i(K) < \infty$, et
\item $K$ est un objet compact.
\end{enumerate}
\end{lemma}

\begin{proof}
Si $K$ est un objet compact, alors $K$ peut être représenté par un complexe
$M^\bullet$ qui, comme $R$-module gradué, est projectif de type fini, voir
Algèbre différentielle graduée, lemme \ref{dga-lemma-compact}.
Puisque $\dim_k R < \infty$, nous concluons que $\sum \dim_k M^i < \infty$
et, a fortiori, que $\sum \dim_k H^i(M^\bullet) < \infty$.
(On peut aussi déduire facilement cette implication de la proposition plus élémentaire
Algèbre différentielle graduée, proposition \ref{dga-proposition-compact}.)

\medskip\noindent
Supposons que $K$ satisfasse (1). Considérons le triangle distingué
de troncation $\tau_{\leq m}K \to K \to \tau_{\geq m + 1}K$, voir
Catégories dérivées, remarque
\ref{derived-remark-truncation-distinguished-triangle}.
Il est clair que $\tau_{\leq m}K$ et $\tau_{\geq m + 1} K$
satisfont (1). Si l'on peut montrer qu'ils sont tous deux compacts, alors $K$ l'est aussi, voir
Catégories dérivées, lemme \ref{derived-lemma-compact-objects-subcategory}.
Ainsi, par récurrence sur le nombre de modules de cohomologie non nuls de $K$,
nous pouvons supposer que $H^i(K)$ n'est non nul que pour une seule valeur de $i$.
En décalant, nous pouvons supposer que $K$ est donné par le complexe
constitué d'un seul $R$-module $M$ de dimension finie, placé
en degré $0$.

\medskip\noindent
Puisque $\dim_k(M) < \infty$, nous voyons que $M$ est artinien comme $R$-module.
Il suffit donc de montrer que tout $R$-module simple représente un
objet compact de $D(R)$. Observons que
$$
I =
\left(
\begin{matrix}
0 & S_1 & S_2 & \ldots & \ldots \\
0 & 0 & S_1 & \ldots & \ldots\\
0 & 0 & 0 & \ldots & \ldots \\
\ldots & \ldots & \ldots & \ldots & \ldots \\
0 & \ldots & \ldots & \ldots & 0
\end{matrix}
\right)
$$
est un idéal bilatère nilpotent de $R$ et que $R/I$
est une $k$-algèbre commutative isomorphe à un produit de $n + 1$ copies de
$k$ (placées sur la diagonale de la matrice, c'est-à-dire que $R/I$ peut être relevée
en une sous-$k$-algèbre de $R$). Il s'ensuit que $R$ possède exactement $n + 1$
classes d'isomorphisme de modules simples $M_0, \ldots, M_n$ (placés le long
de la diagonale). Considérons le $R$-module à droite $P_i$ de vecteurs lignes
$$
P_i =
\left(
\begin{matrix}
0 &
\ldots &
0 &
S_0 &
\ldots &
S_{i - 1} &
S_i
\end{matrix}
\right)
$$
muni de la multiplication évidente $P_i \times R \to P_i$. Nous voyons alors que
$R \cong P_0 \oplus \ldots \oplus P_n$ comme $R$-module à droite. Comme il est clair que
$R$ est un objet compact de $D(R)$, nous concluons que chaque $P_i$ est un objet compact
de $D(R)$. (Nous concluons bien sûr aussi que chaque $P_i$ est projectif
comme $R$-module, mais ce n'est pas ce que nous avons à montrer dans cette démonstration.)
Il est clair que $P_0 = M_0$ est le premier de nos modules simples sur $R$.
Pour $P_1$, nous avons une suite exacte courte
$$
0 \to P_0^{\oplus n + 1} \to P_1 \to M_1 \to 0
$$
qui montre que $M_1$ s'insère dans un triangle distingué dont
les autres membres sont des objets compacts ; ainsi, $M_1$ est un objet compact
de $D(R)$. Plus généralement, il existe une suite exacte courte
$$
0 \to C_i \to P_i \to M_i \to 0
$$
où $C_i$ est un $R$-module de dimension finie dont les facteurs de composition
sont isomorphes à $M_j$ pour $j < i$. Par récurrence, nous concluons d'abord que
$C_i$ détermine un objet compact de $D(R)$, puis que $M_i$
en détermine un également, comme souhaité.
\end{proof}

\begin{lemma}
\label{lemma-coherent-on-projective-space}
Soit $k$ un corps. Soit $n \geq 0$. Soit
$K \in D_\QCoh(\mathcal{O}_{\mathbf{P}^n_k})$.
Les assertions suivantes sont équivalentes :
\begin{enumerate}
\item $K$ appartient à $D^b_{\textit{Coh}}(\mathcal{O}_{\mathbf{P}^n_k})$,
\item $\sum_{i \in \mathbf{Z}}
\dim_k H^i(\mathbf{P}^n_k, E \otimes^\mathbf{L} K) < \infty$
pour tout objet parfait $E$ de
$D(\mathcal{O}_{\mathbf{P}^n_k})$,
\item $\sum_{i \in \mathbf{Z}}
\dim_k \Ext^i_{\mathbf{P}^n_k}(E, K) < \infty$
pour tout objet parfait $E$ de $D(\mathcal{O}_{\mathbf{P}^n_k})$,
\item $\sum_{i \in \mathbf{Z}} \dim_k H^i(\mathbf{P}^n_k,
K \otimes^\mathbf{L} \mathcal{O}_{\mathbf{P}^n_k}(d)) < \infty$
pour $d = 0, 1, \ldots, n$.
\end{enumerate}
\end{lemma}

\begin{proof}
Les assertions (2) et (3) sont équivalentes d'après
Cohomologie, lemme \ref{cohomology-lemma-dual-perfect-complex}.
Si (1) est satisfaite, alors, pour $E$ parfait, le produit tensoriel dérivé
$E \otimes^\mathbf{L} K$ appartient à
$D^b_{\textit{Coh}}(\mathcal{O}_{\mathbf{P}^n_k})$
et l'on voit que (2) est satisfaite d'après
Catégories dérivées des schémas, lemme \ref{perfect-lemma-direct-image-coherent}.
Il est clair que (2) implique (4), puisque $\mathcal{O}_{\mathbf{P}^n_k}(d)$
peut être considéré
comme un objet parfait de la catégorie dérivée de $\mathbf{P}^n_k$.
Ainsi, il suffit de montrer que (4) implique (1).

\medskip\noindent
Supposons (4). Soit $R$ comme dans le lemme \ref{lemma-perfect-for-R}.
Soit $P = \bigoplus_{d = 0, \ldots, n} \mathcal{O}_{\mathbf{P}^n_k}(-d)$.
Rappelons que $R = \text{End}_{\mathbf{P}^n_k}(P)$, tandis que tous les autres groupes
d'extensions de $P$ par lui-même sont nuls, et que $P$ détermine une équivalence
$- \otimes^\mathbf{L} P : D(R) \to D_\QCoh(\mathcal{O}_{\mathbf{P}^n_k})$
d'après Catégories dérivées des schémas, lemme \ref{perfect-lemma-Pn-module-category}.
Disons que $K$ correspond à $L$ dans $D(R)$. Alors
\begin{align*}
H^i(L)
& =
\Ext^i_{D(R)}(R, L) \\
& =
\Ext^i_{\mathbf{P}^n_k}(P, K) \\
& =
H^i(\mathbf{P}^n_k, K \otimes P^\vee) \\
& =
\bigoplus\nolimits_{d = 0, \ldots, n}
H^i(\mathbf{P}^n_k, K \otimes \mathcal{O}(d))
\end{align*}
d'après Algèbre différentielle graduée, lemme
\ref{dga-lemma-upgrade-tensor-with-complex-derived}
(et le fait que $- \otimes^\mathbf{L} P$ est une équivalence)
et Cohomologie, lemme \ref{cohomology-lemma-dual-perfect-complex}.
Ainsi, notre hypothèse (4) implique que $L$ satisfait la condition (2) du
lemme \ref{lemma-perfect-for-R} et est donc un objet compact de $D(R)$.
Par conséquent, $K$ est un objet compact de
$D_\QCoh(\mathcal{O}_{\mathbf{P}^n_k})$.
Ainsi, $K$ est parfait d'après
Catégories dérivées des schémas, proposition
\ref{perfect-proposition-compact-is-perfect}.
Puisque $D_{perf}(\mathcal{O}_{\mathbf{P}^n_k}) =
D^b_{\textit{Coh}}(\mathcal{O}_{\mathbf{P}^n_k})$
d'après
Catégories dérivées des schémas, lemme \ref{perfect-lemma-perfect-on-regular},
nous concluons que (1) est satisfaite.
\end{proof}

\begin{lemma}
\label{lemma-finiteness}
Soit $X$ un schéma propre sur un corps $k$. Soit
$K \in D^b_{\textit{Coh}}(\mathcal{O}_X)$ et soit $E$ dans $D(\mathcal{O}_X)$
un objet parfait. Alors
$\sum_{i \in \mathbf{Z}} \dim_k \Ext^i_X(E, K) < \infty$.
\end{lemma}

\begin{proof}
Cela résulte, par exemple, de la combinaison des
lemmes \ref{perfect-lemma-ext-finite} et
\ref{perfect-lemma-ext-from-perfect-into-bounded-QCoh} de Catégories dérivées des schémas.
Autre démonstration : combiner
les lemmes de Catégories dérivées des schémas
\ref{perfect-lemma-perfect-on-noetherian} et
\ref{perfect-lemma-direct-image-coherent}.
\end{proof}

\begin{lemma}
\label{lemma-characterize-dbcoh-projective}
\begin{reference}
Dans le cas projectif, il s'agit de \cite[Lemme 7.46]{Rouquier-dimensions}
et ce résultat est implicite dans \cite[Théorème A.1]{BvdB}.
\end{reference}
Soit $X$ un schéma propre sur un corps $k$. Soit
$K \in \Ob(D_\QCoh(\mathcal{O}_X))$. Les assertions suivantes sont équivalentes :
\begin{enumerate}
\item $K \in D^b_{\textit{Coh}}(\mathcal{O}_X)$, et
\item $\sum_{i \in \mathbf{Z}} \dim_k \Ext^i_X(E, K) < \infty$
pour tout objet parfait $E$ de $D(\mathcal{O}_X)$.
\end{enumerate}
\end{lemma}

\begin{proof}
L'implication (1) $\Rightarrow$ (2) découle du
lemme \ref{lemma-finiteness}.
L'implication (2) $\Rightarrow$ (1) découle de
Compléments sur les morphismes, lemme
\ref{more-morphisms-lemma-characterize-relatively-perfect}
(voir Catégories dérivées des schémas, exemple
\ref{perfect-example-relatively-perfect-field}
pour la signification d'un objet relativement parfait sur un corps) ;
la démonstration plus simple du cas projectif figure au paragraphe suivant.

\medskip\noindent
Supposons (2) et que $X$ soit projectif sur $k$.
Choisissons une immersion fermée $i : X \to \mathbf{P}^n_k$. Il suffit de montrer
que $Ri_*K$ appartient à $D^b_{\textit{Coh}}(\mathbf{P}^n_k)$, car un module quasi-cohérent
$\mathcal{F}$ sur $X$ est cohérent, resp.\ nul, si et seulement si
$i_*\mathcal{F}$ est cohérent, resp.\ nul. Pour un objet parfait $E$
de $D(\mathcal{O}_{\mathbf{P}^n_k})$, $Li^*E$ est un objet parfait de
$D(\mathcal{O}_X)$ et
$$
\Ext^q_{\mathbf{P}^n_k}(E, Ri_*K) = \Ext^q_X(Li^*E, K)
$$
D'après notre hypothèse, nous voyons donc que
$\sum_{q \in \mathbf{Z}} \dim_k \Ext^q_{\mathbf{P}^n_k}(E, Ri_*K) < \infty$.
Nous concluons par le lemme \ref{lemma-coherent-on-projective-space}.
\end{proof}





\section{Un théorème de représentabilité}
\label{section-bondal-van-den-bergh}

\noindent
Le contenu de cette section est tiré de \cite{BvdB}.

\medskip\noindent
Soit $\mathcal{T}$ une catégorie triangulée $k$-linéaire.
Dans cette section, nous considérons les foncteurs cohomologiques $k$-linéaires
$H$ de $\mathcal{T}$ vers la catégorie des espaces vectoriels sur $k$.
Cela signifiera que $H$ est un foncteur
$$
H : \mathcal{T}^{opp} \longrightarrow \text{Vect}_k
$$
qui est $k$-linéaire et tel que, pour tout triangle distingué
$X \to Y \to Z$ de $\mathcal{T}$, la suite $H(Z) \to H(Y) \to H(X)$
est une suite exacte d'espaces vectoriels sur $k$. Voir
Catégories dérivées, définition \ref{derived-definition-homological}
et Algèbre différentielle graduée, section \ref{dga-section-linear}.

\begin{lemma}
\label{lemma-maps-from-compact-filtered}
Soit $\mathcal{D}$ une catégorie triangulée. Soit
$\mathcal{D}' \subset \mathcal{D}$ une sous-catégorie triangulée pleine. Soit
$X \in \Ob(\mathcal{D})$. La catégorie des flèches $E \to X$ avec
$E \in \Ob(\mathcal{D}')$ est filtrante.
\end{lemma}

\begin{proof}
Vérifions les conditions de
Catégories, définition \ref{categories-definition-directed}.
La catégorie est non vide, car elle contient $0 \to X$.
Si $E_i \to X$, $i = 1, 2$ sont des objets, alors $E_1 \oplus E_2 \to X$
est un objet et il existe des morphismes $(E_i \to X) \to (E_1 \oplus E_2 \to X)$.
Enfin, supposons que $a, b : (E \to X) \to (E' \to X)$ soient des morphismes.
Choisissons un triangle distingué $E \xrightarrow{a - b} E' \to E''$
dans $\mathcal{D}'$. Par l'axiome TR3, on obtient un morphisme de triangles
$$
\xymatrix{
E \ar[r]_{a - b} \ar[d] &
E' \ar[d] \ar[r] & E'' \ar[d] \\
0 \ar[r] &
X \ar[r] &
X
}
$$
et l'on constate que la flèche obtenue $(E' \to X) \to (E'' \to X)$
égalise $a$ et $b$.
\end{proof}

\begin{lemma}
\label{lemma-van-den-bergh}
\begin{reference}
\cite[Lemme 2.14]{CKN}
\end{reference}
Soit $k$ un corps. Soit $\mathcal{D}$ une catégorie triangulée $k$-linéaire
qui admet les sommes directes et est compactement engendrée.
Notons $\mathcal{D}_c$ la sous-catégorie pleine
des objets compacts. Soit $H : \mathcal{D}_c^{opp} \to \text{Vect}_k$
un foncteur cohomologique $k$-linéaire tel que
$\dim_k H(X) < \infty$ pour tout $X \in \Ob(\mathcal{D}_c)$.
Alors $H$ est isomorphe au foncteur $X \mapsto \Hom(X, Y)$
pour un certain $Y \in \Ob(\mathcal{D})$.
\end{lemma}

\begin{proof}
Nous utiliserons Catégories dérivées, lemme
\ref{derived-lemma-compact-objects-subcategory} sans autre mention.
Notons $G : \mathcal{D}_c \to \text{Vect}_k$ le foncteur homologique
$k$-linéaire qui envoie $X$ sur $H(X)^\vee$. Pour tout objet $Y$ de $\mathcal{D}$,
posons
$$
G'(Y) = \colim_{X \to Y, X \in \Ob(\mathcal{D}_c)} G(X)
$$
Il s'agit d'une colimite filtrante d'après le lemme \ref{lemma-maps-from-compact-filtered}.
Nous affirmons que $G'$ est un foncteur homologique $k$-linéaire,
que la restriction de $G'$ à $\mathcal{D}_c$ est $G$, et que $G'$
transforme les sommes directes en sommes directes.

\medskip\noindent
En effet, supposons que $Y_1 \to Y_2 \to Y_3$ soit un triangle distingué.
Soit $\xi \in G'(Y_2)$ dont l'image dans $G'(Y_3)$ est nulle. Comme la colimite est
filtrante, $\xi$ est représenté par un certain $X \to Y_2$ avec
$X \in \Ob(\mathcal{D}_c)$ et $g \in G(X)$.
Le fait que l'image de $\xi$ dans $G'(Y_3)$ soit nulle signifie que le composé
$X \to Y_2 \to Y_3$ se factorise sous la forme $X \to X' \to Y_3$ avec $X' \in \mathcal{D}_c$
et que l'image de $g$ dans $G(X')$ est nulle. Choisissons un triangle distingué
$X'' \to X \to X'$. Alors $X'' \in \Ob(\mathcal{D}_c)$.
Comme $G$ est homologique, on voit que $g$ est l'image d'un certain
$g'' \in G'(X'')$. Par l'axiome TR3, les morphismes $X \to Y_2$ et $X' \to Y_3$ s'insèrent dans
un morphisme de triangles distingués
$(X'' \to X \to X') \to (Y_1 \to Y_2 \to Y_3)$,
et l'on voit que $\xi$ est bien l'image de
l'élément de $G'(Y_1)$ représenté par $X'' \to Y_1$ et $g'' \in G(X'')$.

\medskip\noindent
Si $Y \in \Ob(\mathcal{D}_c)$, alors $\text{id} : Y \to Y$ est l'objet final
de la catégorie des flèches $X \to Y$ avec $X \in \Ob(\mathcal{D}_c)$.
On voit donc que $G'(Y) = G(Y)$ dans ce cas, ce qui prouve l'assertion
sur la restriction. Soit $Y = \bigoplus_{i \in I} Y_i$
une somme directe. Soient $a : X \to Y$ avec $X \in \Ob(\mathcal{D}_c)$
et $g \in G(X)$ représentant un élément $\xi$ de $G'(Y)$.
Le morphisme $a : X \to Y$ s'écrit de façon unique comme une somme de morphismes
$a_i : X \to Y_i$ presque tous nuls, puisque $X$ est un objet compact de $\mathcal{D}$.
Soit $I' = \{i \in I \mid a_i \not = 0\}$. On peut alors factoriser
$a$ comme le composé
$$
X \xrightarrow{(1, \ldots, 1)}
\bigoplus\nolimits_{i \in I'} X
\xrightarrow{\bigoplus_{i \in I'} a_i}
\bigoplus\nolimits_{i \in I} Y_i = Y
$$
On en déduit que $\xi = \sum_{i \in I'} \xi_i$
est la somme des images des éléments
$\xi_i \in G'(Y_i)$ correspondant à $a_i : X \to Y_i$
et $g \in G(X)$. Ainsi, $\bigoplus G'(Y_i) \to G'(Y)$
est surjective. Nous omettons la vérification (triviale) de son injectivité.

\medskip\noindent
Il s'ensuit que le foncteur $Y \mapsto G'(Y)^\vee$ est cohomologique
et transforme les sommes directes en produits directs. La représentabilité de Brown,
voir Catégories dérivées, proposition \ref{derived-proposition-brown},
montre donc qu'il existe $Y \in \Ob(\mathcal{D})$
et un isomorphisme
$G'(Z)^\vee = \Hom(Z, Y)$ fonctoriel en $Z$.
Pour $X \in \Ob(\mathcal{D}_c)$, on a
$G'(X)^\vee = G(X)^\vee = (H(X)^\vee)^\vee = H(X)$
car $\dim_k H(X) < \infty$, ce qui achève la démonstration.
\end{proof}

\begin{theorem}
\label{theorem-bondal-van-den-bergh}
\begin{reference}
Dans le cas projectif, il s'agit de \cite[Théorème A.1]{BvdB}
\end{reference}
Soit $X$ un schéma propre sur un corps $k$.
Soit $F : D_{perf}(\mathcal{O}_X)^{opp} \to \text{Vect}_k$
un foncteur cohomologique $k$-linéaire tel que
$$
\sum\nolimits_{n \in \mathbf{Z}} \dim_k F(E[n]) < \infty
$$
pour tout $E \in D_{perf}(\mathcal{O}_X)$. Alors $F$ est isomorphe à un foncteur
de la forme $E \mapsto \Hom_X(E, K)$ pour un certain
$K \in D^b_{\textit{Coh}}(\mathcal{O}_X)$.
\end{theorem}

\begin{proof}
La catégorie dérivée $D_\QCoh(\mathcal{O}_X)$ admet les sommes directes,
est compactement engendrée, et $D_{perf}(\mathcal{O}_X)$ est la sous-catégorie pleine
des objets compacts, voir
Catégories dérivées des schémas, lemme
\ref{perfect-lemma-quasi-coherence-direct-sums},
théorème \ref{perfect-theorem-bondal-van-den-Bergh}, et
proposition \ref{perfect-proposition-compact-is-perfect}.
Par le lemme \ref{lemma-van-den-bergh}, nous pouvons supposer que
$F(E) = \Hom_X(E, K)$ pour un certain $K \in \Ob(D_\QCoh(\mathcal{O}_X))$.
Il s'ensuit alors que $K$ appartient à $D^b_{\textit{Coh}}(\mathcal{O}_X)$
par le lemme \ref{lemma-characterize-dbcoh-projective}.
\end{proof}

\begin{lemma}
\label{lemma-homological-representable}
Soit $X$ un schéma propre et régulier sur un corps $k$.
Soit $G : D_{perf}(\mathcal{O}_X) \to \text{Vect}_k$
un foncteur homologique $k$-linéaire tel que
$$
\sum\nolimits_{n \in \mathbf{Z}} \dim_k G(E[n]) < \infty
$$
pour tout $E \in D_{perf}(\mathcal{O}_X)$. Alors $G$ est isomorphe à un foncteur
de la forme $E \mapsto \Hom_X(K, E)$ pour un certain $K \in D_{perf}(\mathcal{O}_X)$.
\end{lemma}

\begin{proof}
Considérons le foncteur contravariant $E \mapsto E^\vee$
sur $D_{perf}(\mathcal{O}_X)$, voir
Cohomologie, lemme \ref{cohomology-lemma-dual-perfect-complex}.
Ce foncteur est une anti-auto-équivalence exacte de $D_{perf}(\mathcal{O}_X)$.
Nous pouvons donc appliquer le théorème \ref{theorem-bondal-van-den-bergh}
au foncteur $F(E) = G(E^\vee)$ pour trouver
$K \in D_{perf}(\mathcal{O}_X)$ tel que $G(E^\vee) = \Hom_X(E, K)$.
Il s'ensuit que $G(E) = \Hom_X(E^\vee, K) = \Hom_X(K^\vee, E)$,
et l'on conclut que $K^\vee$ convient.
\end{proof}






\section{Existence d'adjoints}
\label{section-adjoints}

\noindent
Les résultats de l'article de Bondal et van den Bergh ont pour conséquence
l'énoncé suivant sur l'existence automatique d'adjoints.

\begin{lemma}
\label{lemma-always-right-adjoints}
Soit $k$ un corps. Soient $X$ et $Y$ des schémas propres sur $k$.
Si $X$ est régulier, tout foncteur exact $k$-linéaire
$F : D_{perf}(\mathcal{O}_X) \to D_{perf}(\mathcal{O}_Y)$
admet un adjoint à droite exact et un adjoint à gauche exact.
\end{lemma}

\begin{proof}
S'il existe, un adjoint est un foncteur exact d'après le lemme très général
\ref{derived-lemma-adjoint-is-exact} de Catégories dérivées.

\medskip\noindent
Démontrons l'existence d'un adjoint à droite.
Pour cela, il suffit de montrer que, pour
$M \in D_{perf}(\mathcal{O}_Y)$, le foncteur contravariant
$K \mapsto \Hom_Y(F(K), M)$ est représentable.
Ce foncteur est contravariant, $k$-linéaire et cohomologique.
Par le théorème \ref{theorem-bondal-van-den-bergh},
il suffit donc de montrer que
$$
\sum\nolimits_{i \in \mathbf{Z}} \dim_k \Ext^i_Y(F(K), M) < \infty
$$
Cela résulte du lemme \ref{lemma-finiteness}.

\medskip\noindent
Pour l'existence de l'adjoint à gauche, on raisonne de même
en utilisant le lemme \ref{lemma-homological-representable}
au lieu du théorème \ref{theorem-bondal-van-den-bergh}.
\end{proof}






\section{Foncteurs de Fourier-Mukai}
\label{section-fourier-mukai}

\noindent
Ces foncteurs ont été introduits pour la première fois dans \cite{Mukai}.

\begin{definition}
\label{definition-fourier-mukai-functor}
Soit $S$ un schéma. Soient $X$ et $Y$ des schémas sur $S$.
Soit $K \in D(\mathcal{O}_{X \times_S Y})$. Le foncteur exact
$$
\Phi_K : D(\mathcal{O}_X) \longrightarrow D(\mathcal{O}_Y),\quad
M \longmapsto R\text{pr}_{2, *}(
L\text{pr}_1^*M \otimes_{\mathcal{O}_{X \times_S Y}}^\mathbf{L} K)
$$
de catégories triangulées est appelé un {\it foncteur de Fourier-Mukai},
et $K$ est appelé un {\it noyau de Fourier-Mukai} pour ce foncteur.
De plus,
\begin{enumerate}
\item si $\Phi_K$ envoie $D_\QCoh(\mathcal{O}_X)$ dans $D_\QCoh(\mathcal{O}_Y)$,
alors le foncteur exact obtenu
$\Phi_K : D_\QCoh(\mathcal{O}_X) \to D_\QCoh(\mathcal{O}_Y)$
est appelé un foncteur de Fourier-Mukai,
\item si $\Phi_K$ envoie $D_{perf}(\mathcal{O}_X)$ dans
$D_{perf}(\mathcal{O}_Y)$, alors le foncteur exact obtenu
$\Phi_K : D_{perf}(\mathcal{O}_X) \to D_{perf}(\mathcal{O}_Y)$
est appelé un foncteur de Fourier-Mukai, et
\item si $X$ et $Y$ sont noethériens et si $\Phi_K$ envoie
$D^b_{\textit{Coh}}(\mathcal{O}_X)$ dans $D^b_{\textit{Coh}}(\mathcal{O}_Y)$,
alors le foncteur exact obtenu
$\Phi_K : D^b_{\textit{Coh}}(\mathcal{O}_X) \to
D^b_{\textit{Coh}}(\mathcal{O}_Y)$
est appelé un foncteur de Fourier-Mukai.
De même pour $D_{\textit{Coh}}$, $D^+_{\textit{Coh}}$, $D^-_{\textit{Coh}}$.
\end{enumerate}
\end{definition}

\begin{lemma}
\label{lemma-fourier-Mukai-QCoh}
Soit $S$ un schéma. Soient $X$ et $Y$ des schémas sur $S$.
Soit $K \in D(\mathcal{O}_{X \times_S Y})$.
Le foncteur de Fourier-Mukai correspondant $\Phi_K$ envoie
$D_\QCoh(\mathcal{O}_X)$ dans $D_\QCoh(\mathcal{O}_Y)$
si $K$ appartient à $D_\QCoh(\mathcal{O}_{X \times_S Y})$ et si $X \to S$ est
quasi-compact et quasi-séparé.
\end{lemma}

\begin{proof}
Cela résulte du fait que l'image inverse dérivée préserve
$D_\QCoh$
(Catégories dérivées des schémas, lemme
\ref{perfect-lemma-quasi-coherence-pullback}),
que les produits tensoriels dérivés préservent $D_\QCoh$
(Catégories dérivées des schémas, lemme
\ref{perfect-lemma-quasi-coherence-tensor-product}),
que la projection $\text{pr}_2 : X \times_S Y \to Y$ est
quasi-compacte et quasi-séparée
(Schémas, lemmes
\ref{schemes-lemma-quasi-compact-preserved-base-change} et
\ref{schemes-lemma-separated-permanence}), et
que l'image directe totale par un morphisme quasi-séparé et quasi-compact
préserve $D_\QCoh$
(Catégories dérivées des schémas, lemme
\ref{perfect-lemma-quasi-coherence-direct-image}).
\end{proof}

\begin{lemma}
\label{lemma-compose-fourier-mukai}
Soit $S$ un schéma. Soient $X, Y, Z$ des schémas sur $S$. Supposons que
$X \to S$, $Y \to S$ et $Z \to S$ soient quasi-compacts et quasi-séparés.
Soit $K \in D_\QCoh(\mathcal{O}_{X \times_S Y})$.
Soit $K' \in D_\QCoh(\mathcal{O}_{Y \times_S Z})$.
Considérons les foncteurs de Fourier-Mukai
$\Phi_K : D_\QCoh(\mathcal{O}_X) \to D_\QCoh(\mathcal{O}_Y)$
et $\Phi_{K'} : D_\QCoh(\mathcal{O}_Y) \to D_\QCoh(\mathcal{O}_Z)$.
Si $X$ et $Z$ sont tor-indépendants sur $S$ et si $Y \to S$ est plat,
alors
$$
\Phi_{K'} \circ \Phi_K = \Phi_{K''} :
D_\QCoh(\mathcal{O}_X)
\longrightarrow
D_\QCoh(\mathcal{O}_Z)
$$
où
$$
K'' = R\text{pr}_{13, *}(
L\text{pr}_{12}^*K
\otimes_{\mathcal{O}_{X \times_S Y \times_S Z}}^\mathbf{L}
L\text{pr}_{23}^*K')
$$
dans $D_\QCoh(\mathcal{O}_{X \times_S Z})$.
\end{lemma}

\begin{proof}
L'énoncé a un sens d'après le lemme \ref{lemma-fourier-Mukai-QCoh}.
Nous allons utiliser
Catégories dérivées des schémas, lemmes
\ref{perfect-lemma-quasi-coherence-pullback},
\ref{perfect-lemma-quasi-coherence-tensor-product}, et
\ref{perfect-lemma-quasi-coherence-direct-image},
ainsi que Schémas, lemmes
\ref{schemes-lemma-quasi-compact-preserved-base-change} et
\ref{schemes-lemma-separated-permanence},
sans autre mention.
D'après Catégories dérivées des schémas, lemme
\ref{perfect-lemma-flat-base-change-tor-independent},
$X \times_S Y$ et $Y \times_S Z$ sont tor-indépendants
sur $Y$. Autrement dit, le changement de base s'applique au diagramme cartésien
$$
\xymatrix{
X \times_S Y \times_S Z \ar[d] \ar[r] &
Y \times_S Z \ar[d]^{p^{YZ}_Y} \\
X \times_S Y \ar[r]^{p^{XY}_Y} & Y
}
$$
pour les complexes à faisceaux de cohomologie quasi-cohérents, voir
Catégories dérivées des schémas, lemme \ref{perfect-lemma-compare-base-change}.
En posant $p^* = Lp^*$, $p_* = Rp_*$ et $\otimes = \otimes^\mathbf{L}$,
on a, pour $M \in D_\QCoh(\mathcal{O}_X)$, la suite d'égalités
\begin{align*}
\Phi_{K'}(\Phi_K(M))
& =
p^{YZ}_{Z, *}(p^{YZ, *}_Y p^{XY}_{Y, *}(p^{XY, *}_X M \otimes K) \otimes K') \\
& =
p^{YZ}_{Z, *}(\text{pr}_{23, *} \text{pr}_{12}^*(p^{XY, *}_X M \otimes K)
\otimes K') \\
& =
p^{YZ}_{Z, *}(\text{pr}_{23, *}(\text{pr}_1^*M \otimes \text{pr}_{12}^*K)
\otimes K') \\
& =
p^{YZ}_{Z, *}(\text{pr}_{23, *}(\text{pr}_1^*M \otimes \text{pr}_{12}^*K
\otimes \text{pr}_{23}^*K')) \\
& =
\text{pr}_{3, *}(\text{pr}_1^*M \otimes \text{pr}_{12}^*K
\otimes \text{pr}_{23}^*K') \\
& =
p^{XZ}_{Z, *}\text{pr}_{13, *}(\text{pr}_1^*M \otimes \text{pr}_{12}^*K
\otimes \text{pr}_{23}^*K') \\
& =
p^{XZ}_{Z, *} (p^{XZ, *}_X M \otimes \text{pr}_{13, *}(\text{pr}_{12}^*K
\otimes \text{pr}_{23}^*K'))
\end{align*}
comme voulu. Nous avons utilisé ici la remarque sur le changement de base dans la
deuxième égalité et Catégories dérivées des schémas, lemme
\ref{perfect-lemma-cohomology-base-change}, dans la $4$e et la
dernière égalité.
\end{proof}

\begin{lemma}
\label{lemma-fourier-mukai}
Soit $S$ un schéma. Soient $X$ et $Y$ des schémas sur $S$.
Soit $K \in D(\mathcal{O}_{X \times_S Y})$.
Le foncteur de Fourier-Mukai correspondant $\Phi_K$ envoie
$D_{perf}(\mathcal{O}_X)$ dans $D_{perf}(\mathcal{O}_Y)$ si au moins
l'une des conditions suivantes est satisfaite :
\begin{enumerate}
\item $S$ est noethérien, $X \to S$ et $Y \to S$ sont de type fini,
$K \in D^b_{\textit{Coh}}(\mathcal{O}_{X \times_S Y})$, le support de $H^i(K)$
est propre sur $Y$ pour tout $i$, et $K$ est de tor-dimension finie
comme objet de $D(\text{pr}_2^{-1}\mathcal{O}_Y)$,
\item $X \to S$ est de présentation finie et $K$ peut être représenté
par un complexe borné $\mathcal{K}^\bullet$ de modules de présentation finie sur
$\mathcal{O}_{X \times_S Y}$, plat sur $Y$, et dont le support
est propre sur $Y$,
\item $X \to S$ est un morphisme propre, plat et de présentation finie,
et $K$ est parfait,
\item $S$ est noethérien, $X \to S$ est plat et propre, et $K$ est parfait,
\item $X \to S$ est un morphisme propre, plat et de présentation finie,
et $K$ est $Y$-parfait,
\item $S$ est noethérien, $X \to S$ est plat et propre, et $K$ est
$Y$-parfait.
\end{enumerate}
\end{lemma}

\begin{proof}
Si $M$ est parfait sur $X$, alors $L\text{pr}_1^*M$
est parfait sur $X \times_S Y$, voir
Cohomologie, lemme \ref{cohomology-lemma-perfect-pullback}.
Nous utiliserons ce fait sans autre mention ci-dessous.
Nous utiliserons également le fait que, si $X \to S$ est de type fini, propre,
plat ou de présentation finie, il en va de même pour
le changement de base $\text{pr}_2 : X \times_S Y \to Y$, voir
Morphismes, lemmes
\ref{morphisms-lemma-base-change-finite-type},
\ref{morphisms-lemma-base-change-proper},
\ref{morphisms-lemma-base-change-flat}, et
\ref{morphisms-lemma-base-change-finite-presentation}.

\medskip\noindent
L'assertion (1) résulte de
Catégories dérivées des schémas, lemme \ref{perfect-lemma-perfect-direct-image},
combiné avec
Catégories dérivées des schémas, lemme \ref{perfect-lemma-perfect-on-noetherian}.

\medskip\noindent
L'assertion (2) résulte de
Catégories dérivées des schémas, lemme
\ref{perfect-lemma-base-change-tensor-perfect}.

\medskip\noindent
L'assertion (3) résulte de
Catégories dérivées des schémas, lemme
\ref{perfect-lemma-flat-proper-perfect-direct-image-general}.

\medskip\noindent
L'assertion (4) résulte de l'assertion (3) et du fait qu'un morphisme de type fini
entre schémas noethériens est de présentation finie d'après Morphismes, lemme
\ref{morphisms-lemma-noetherian-finite-type-finite-presentation}.

\medskip\noindent
L'assertion (5) résulte de
Catégories dérivées des schémas, lemme
\ref{perfect-lemma-derived-pushforward-rel-perfect}, combiné avec
Catégories dérivées des schémas, lemme
\ref{perfect-lemma-perfect-relatively-perfect}.

\medskip\noindent
L'assertion (6) résulte de l'assertion (5) de la même manière que l'assertion (4) résulte de
l'assertion (3).
\end{proof}

\begin{lemma}
\label{lemma-fourier-mukai-Coh}
Soit $S$ un schéma noethérien. Soient $X$ et $Y$ des schémas de type fini
sur $S$. Soit $K \in D^b_{\textit{Coh}}(\mathcal{O}_{X \times_S Y})$.
Le foncteur de Fourier-Mukai correspondant $\Phi_K$ envoie
$D^b_{\textit{Coh}}(\mathcal{O}_X)$ dans $D^b_{\textit{Coh}}(\mathcal{O}_Y)$
si au moins l'une des conditions suivantes est satisfaite :
\begin{enumerate}
\item le support de $H^i(K)$ est propre sur $Y$ pour tout $i$, et $K$
est de tor-dimension finie comme objet de $D(\text{pr}_1^{-1}\mathcal{O}_X)$,
\item $K$ peut être représenté par un complexe borné $\mathcal{K}^\bullet$
de $\mathcal{O}_{X \times_S Y}$-modules cohérents, plat sur $X$, et dont le support
est propre sur $Y$,
\item le support de $H^i(K)$ est propre sur $Y$ pour tout $i$
et $X$ est un schéma régulier,
\item $K$ est parfait, le support de $H^i(K)$ est propre sur $Y$ pour tout $i$,
et $Y \to S$ est plat.
\end{enumerate}
En outre, dans chacun des cas, la condition sur le support est automatique
si $X \to S$ est propre.
\end{lemma}

\begin{proof}
Soit $M$ un objet de $D^b_{\textit{Coh}}(\mathcal{O}_X)$.
Dans chaque cas, nous utiliserons Catégories dérivées des schémas, lemme
\ref{perfect-lemma-direct-image-coherent}, pour montrer que
$$
\Phi_K(M) = R\text{pr}_{2, *}(
L\text{pr}_1^*M
\otimes_{\mathcal{O}_{X \times_S Y}}^\mathbf{L}
K)
$$
appartient à $D^b_{\textit{Coh}}(\mathcal{O}_Y)$. Le produit tensoriel dérivé
$L\text{pr}_1^*M \otimes_{\mathcal{O}_{X \times_S Y}}^\mathbf{L} K$
est un objet pseudo-cohérent de $D(\mathcal{O}_{X \times_S Y})$
(d'après
Cohomologie, lemme \ref{cohomology-lemma-pseudo-coherent-pullback},
Catégories dérivées des schémas, lemme
\ref{perfect-lemma-identify-pseudo-coherent-noetherian}, et
Cohomologie, lemme \ref{cohomology-lemma-tensor-pseudo-coherent}),
et possède donc des faisceaux de cohomologie cohérents (à nouveau d'après
Catégories dérivées des schémas, lemme
\ref{perfect-lemma-identify-pseudo-coherent-noetherian}).
Dans chaque cas, les supports des faisceaux de cohomologie
$H^i(L\text{pr}_1^*M \otimes_{\mathcal{O}_{X \times_S Y}}^\mathbf{L} K)$
sont propres sur $Y$, car ces supports sont contenus dans la
réunion des supports des $H^i(K)$. Il suffit donc, dans chaque cas,
de prouver que ce produit tensoriel est borné inférieurement.

\medskip\noindent
Cas (1). D'après Cohomologie, lemme \ref{cohomology-lemma-variant-derived-pullback},
on a
$$
L\text{pr}_1^*M
\otimes_{\mathcal{O}_{X \times_S Y}}^\mathbf{L}
K
\cong
\text{pr}_1^{-1}M
\otimes_{\text{pr}_1^{-1}\mathcal{O}_X}^\mathbf{L}
K
$$
avec les notations évidentes. L'hypothèse de tor-dimension finie
et le fait que $M$ ne possède qu'un nombre fini de
faisceaux de cohomologie non nuls impliquent donc la borne voulue.

\medskip\noindent
Le cas (2) résulte du fait qu'ici l'hypothèse implique que $K$ est de
tor-dimension finie comme objet de $D(\text{pr}_1^{-1}\mathcal{O}_X)$ ;
l'argument du paragraphe précédent s'applique donc.

\medskip\noindent
Dans le cas (3), $K$ est également de tor-dimension finie
comme objet de $D(\text{pr}_1^{-1}\mathcal{O}_X)$. En effet, choisissons
des ouverts affines $U = \Spec(A)$ et $V = \Spec(B)$, respectivement de $X$ et $Y$, dont les images sont contenues dans
l'ouvert affine $W = \Spec(R)$ de $S$. Alors
$K|_{U \times V}$ est donné par un complexe borné de
$A \otimes_R B$-modules de type fini $M^\bullet$. Comme $A$ est un anneau régulier
de dimension finie, chaque $M^i$ est de dimension projective finie
comme $A$-module (Algèbre, lemme
\ref{algebra-lemma-finite-gl-dim-finite-dim-regular})
et donc de tor-dimension finie comme $A$-module.
Ainsi, $M^\bullet$ est de tor-dimension finie comme complexe de $A$-modules
(Compléments sur l'algèbre, lemme
\ref{more-algebra-lemma-complex-finite-tor-dimension-modules}).
Puisque $X \times Y$ est quasi-compact, il existe $[a, b]$
tel que, pour tout point $z \in X \times Y$, le germe $K_z$
soit d'amplitude de tor contenue dans $[a, b]$ sur $\mathcal{O}_{X, \text{pr}_1(z)}$.
Cela implique que $K$ est de tor-dimension bornée comme objet de
$D(\text{pr}_1^{-1}\mathcal{O}_X)$, voir
Cohomologie, lemme \ref{cohomology-lemma-tor-amplitude-stalk}.
On conclut comme dans les deux paragraphes précédents.

\medskip\noindent
Cas (4). Avec les notations ci-dessus, l'homomorphisme d'anneaux $R \to B$ est plat.
L'homomorphisme d'anneaux $A \to A \otimes_R B$ est donc plat. Tout module projectif
sur $A \otimes_R B$ est donc plat sur $A$. Ainsi, tout complexe parfait de
$A \otimes_R B$-modules est de tor-dimension finie comme complexe
de $A$-modules, et l'on conclut comme précédemment.
\end{proof}

\begin{example}
\label{example-diagonal-fourier-mukai}
Soit $X \to S$ un morphisme séparé de schémas. Alors la diagonale
$\Delta : X \to X \times_S X$ est une immersion fermée et, par conséquent,
$\mathcal{O}_\Delta = \Delta_*\mathcal{O}_X = R\Delta_*\mathcal{O}_X$
est un $\mathcal{O}_{X \times_S X}$-module quasi-cohérent de type fini
qui est plat sur $X$ (pour l'une ou l'autre projection). Le foncteur de Fourier-Mukai
$\Phi_{\mathcal{O}_\Delta}$ est égal à l'identité dans ce cas.
En effet, pour tout $M \in D(\mathcal{O}_X)$, on a
\begin{align*}
L\text{pr}_1^*M \otimes_{\mathcal{O}_{X \times_S X}}^\mathbf{L}
\mathcal{O}_\Delta
& =
L\text{pr}_1^*M \otimes_{\mathcal{O}_{X \times_S X}}^\mathbf{L}
R\Delta_*\mathcal{O}_X \\
& =
R\Delta_*(
L\Delta^*L\text{pr}_1^*M \otimes_{\mathcal{O}_X}^\mathbf{L} \mathcal{O}_X) \\
& =
R\Delta_*(M)
\end{align*}
Nous avons examiné la première égalité ci-dessus.
La deuxième égalité est donnée par Cohomologie, lemme
\ref{cohomology-lemma-projection-formula-closed-immersion}.
La troisième vient de ce que $\text{pr}_1 \circ \Delta = \text{id}_X$ et du
lemme \ref{cohomology-lemma-derived-pullback-composition} de Cohomologie.
Pour obtenir l'image directe vers $X$, appliquons à ce qui précède le foncteur $R\text{pr}_{2, *}$ ;
on obtient $M$ d'après
Cohomologie, lemme \ref{cohomology-lemma-derived-pushforward-composition},
et le fait que $\text{pr}_2 \circ \Delta = \text{id}_X$.
\end{example}

\begin{lemma}
\label{lemma-fourier-mukai-right-adjoint}
\begin{reference}
Comparer avec la discussion de \cite{Rizzardo}.
\end{reference}
Soient $X \to S$ et $Y \to S$ des morphismes de schémas quasi-compacts et quasi-séparés.
Soit $\Phi : D_\QCoh(\mathcal{O}_X) \to D_\QCoh(\mathcal{O}_Y)$
un foncteur de Fourier-Mukai à noyau pseudo-cohérent
$K \in D_\QCoh(\mathcal{O}_{X \times_S Y})$.
Soit $a : D_\QCoh(\mathcal{O}_Y) \to  D_\QCoh(\mathcal{O}_{X \times_S Y})$
l'adjoint à droite de $R\text{pr}_{2, *}$, voir
Dualité pour les schémas, lemme \ref{duality-lemma-twisted-inverse-image}.
Notons
$$
K' = (Y \times_S X \to X \times_S Y)^*
R\SheafHom_{\mathcal{O}_{X \times_S Y}}(K, a(\mathcal{O}_Y)) \in
D_\QCoh(\mathcal{O}_{Y \times_S X})
$$
et notons $\Phi' : D_\QCoh(\mathcal{O}_Y) \to D_\QCoh(\mathcal{O}_X)$
la transformée de Fourier-Mukai correspondante. Il existe une application canonique
$$
\Hom_X(M, \Phi'(N)) \longrightarrow \Hom_Y(\Phi(M), N)
$$
fonctorielle en $M$ dans $D_\QCoh(\mathcal{O}_X)$ et en $N$ dans
$D_\QCoh(\mathcal{O}_Y)$, qui est un isomorphisme si
\begin{enumerate}
\item $N$ est parfait, ou
\item $K$ est parfait et $X \to S$ est propre, plat et de présentation finie.
\end{enumerate}
\end{lemma}

\begin{proof}
Le lemme \ref{lemma-fourier-Mukai-QCoh} donne un foncteur $\Phi$
comme dans l'énoncé. Observons que $a(\mathcal{O}_Y)$ appartient à
$D^+_\QCoh(\mathcal{O}_{X \times_S Y})$ d'après Dualité pour les schémas,
lemme \ref{duality-lemma-twisted-inverse-image-bounded-below}.
Ainsi, pour $K$ pseudo-cohérent, on a
$K' \in D_\QCoh(\mathcal{O}_{Y \times_S X})$
d'après Catégories dérivées des schémas, lemme
\ref{perfect-lemma-quasi-coherence-internal-hom},
et l'on obtient $\Phi'$ comme indiqué.

\medskip\noindent
Nous abrégeons
$\otimes^\mathbf{L} = \otimes_{\mathcal{O}_{X \times_S Y}}^\mathbf{L}$
et
$\SheafHom = R\SheafHom_{\mathcal{O}_{X \times_S Y}}$.
Soit $M$ dans $D_\QCoh(\mathcal{O}_X)$ et soit
$N$ dans $D_\QCoh(\mathcal{O}_Y)$. On a
\begin{align*}
\Hom_Y(\Phi(M), N)
& =
\Hom_Y(R\text{pr}_{2, *}(L\text{pr}_1^*M \otimes^\mathbf{L} K), N) \\
& =
\Hom_{X \times_S Y}(L\text{pr}_1^*M \otimes^\mathbf{L} K, a(N)) \\
& =
\Hom_{X \times_S Y}(L\text{pr}_1^*M,
R\SheafHom(K, a(N))) \\
& =
\Hom_X(M, R\text{pr}_{1, *}R\SheafHom(K, a(N)))
\end{align*}
où nous avons utilisé Cohomologie, lemmes \ref{cohomology-lemma-internal-hom}
et \ref{cohomology-lemma-adjoint}. Il existe des applications canoniques
$$
L\text{pr}_2^*N \otimes^\mathbf{L} R\SheafHom(K, a(\mathcal{O}_Y))
\xrightarrow{\alpha}
R\SheafHom(K, L\text{pr}_2^*N \otimes^\mathbf{L} a(\mathcal{O}_Y))
\xrightarrow{\beta}
R\SheafHom(K, a(N))
$$
Ici, $\alpha$ est l'application de
Cohomologie, lemme \ref{cohomology-lemma-internal-hom-diagonal-better},
et $\beta$ celle de Dualité pour les schémas, équation
(\ref{duality-equation-compare-with-pullback}).
En composant toutes ces flèches, on obtient la flèche fonctorielle affichée
dans l'énoncé du lemme.

\medskip\noindent
La flèche $\alpha$ est un isomorphisme d'après
Catégories dérivées des schémas, lemme
\ref{perfect-lemma-internal-hom-evaluate-tensor-isomorphism},
dès que $K$ ou $N$ est parfait.
La flèche $\beta$ est un isomorphisme si $N$ est parfait, d'après
Dualité pour les schémas, lemme \ref{duality-lemma-compare-with-pullback-perfect},
ou, en général, si $X \to S$ est
plat, propre et de présentation finie, d'après
Dualité pour les schémas, lemme
\ref{duality-lemma-compare-with-pullback-flat-proper}.
\end{proof}

\begin{lemma}
\label{lemma-fourier-mukai-left-adjoint}
\begin{reference}
Comparer avec la discussion de \cite{Rizzardo}.
\end{reference}
Soit $S$ un schéma noethérien. Soit $Y \to S$ un morphisme plat, propre
et de Gorenstein, et soit $X \to S$ un morphisme de type fini.
Notons $\omega^\bullet_{Y/S}$ le complexe dualisant relatif de
$Y$ sur $S$. Soit $\Phi : D_\QCoh(\mathcal{O}_X) \to D_\QCoh(\mathcal{O}_Y)$
un foncteur de Fourier-Mukai à noyau parfait
$K \in D_\QCoh(\mathcal{O}_{X \times_S Y})$. Notons
$$
K' = (Y \times_S X \to X \times_S Y)^*(K^\vee
\otimes_{\mathcal{O}_{X \times_S Y}}^\mathbf{L}
L\text{pr}_2^*\omega^\bullet_{Y/S})
\in
D_\QCoh(\mathcal{O}_{Y \times_S X})
$$
et notons $\Phi' : D_\QCoh(\mathcal{O}_Y) \to D_\QCoh(\mathcal{O}_X)$
la transformée de Fourier-Mukai correspondante. Il existe un isomorphisme
canonique
$$
\Hom_Y(N, \Phi(M)) \longrightarrow \Hom_X(\Phi'(N), M)
$$
fonctoriel en $M$ dans $D_\QCoh(\mathcal{O}_X)$ et en $N$ dans
$D_\QCoh(\mathcal{O}_Y)$.
\end{lemma}

\begin{proof}
Le lemme \ref{lemma-fourier-Mukai-QCoh} donne un foncteur $\Phi$
comme dans l'énoncé.

\medskip\noindent
Observons que la formation du complexe dualisant relatif commute
au changement de base dans notre situation, voir Dualité pour les schémas,
remarque \ref{duality-remark-relative-dualizing-complex}.
Ainsi, $L\text{pr}_2^*\omega^\bullet_{Y/S} = \omega^\bullet_{X \times_S Y/X}$.
De plus, $\omega^\bullet_{Y/S}$ est un objet
inversible de la catégorie dérivée, voir Dualité pour les schémas, lemme
\ref{duality-lemma-affine-flat-Noetherian-gorenstein}, et il est a fortiori
parfait.

\medskip\noindent
Pour démontrer effectivement le lemme, nous allons recourir à un artifice. Nous allons
montrer qu'en échangeant les rôles de $X$ et $Y$, ainsi que ceux de $K$ et $K'$,
on se trouve dans la situation du lemme \ref{lemma-fourier-mukai-right-adjoint},
ce qui donne le résultat. Il est clair que $K'$ est parfait comme
produit tensoriel d'objets parfaits, de sorte que la discussion du
lemme \ref{lemma-fourier-mukai-right-adjoint} s'y applique.
Pour montrer que la procédure du
lemme \ref{lemma-fourier-mukai-right-adjoint}, appliquée à $K'$
sur $Y \times_S X$, produit un complexe isomorphe à $K$, il suffit
(détails omis) de montrer que
$$
R\SheafHom(R\SheafHom(K, \omega^\bullet_{X \times_S Y/X}),
\omega^\bullet_{X \times_S Y/X}) = K
$$
C'est clair, car $K$ est parfait et $\omega^\bullet_{X \times_S Y/X}$
est inversible ; nous omettons les détails. Ainsi, le
lemme \ref{lemma-fourier-mukai-right-adjoint} fournit une application
$$
\Hom_Y(N, \Phi(M)) \longrightarrow \Hom_X(\Phi'(N), M)
$$
fonctorielle en $M$ dans $D_\QCoh(\mathcal{O}_X)$ et en $N$ dans
$D_\QCoh(\mathcal{O}_Y)$, qui est un isomorphisme puisque
$K'$ est parfait. Cela achève la démonstration.
\end{proof}

\begin{lemma}
\label{lemma-fourier-mukai-flat-proper-over-noetherian}
Soit $S$ un schéma noethérien.
\begin{enumerate}
\item Pour $X$ et $Y$ propres et plats sur $S$ et $K$ dans
$D_{perf}(\mathcal{O}_{X \times_S Y})$, on obtient un foncteur de Fourier-Mukai
$\Phi_K : D_{perf}(\mathcal{O}_X) \to D_{perf}(\mathcal{O}_Y)$.
\item Pour $X$, $Y$ et $Z$ propres et plats sur $S$, $K \in
D_{perf}(\mathcal{O}_{X \times_S Y})$, $K' \in
D_{perf}(\mathcal{O}_{Y \times_S Z})$, le composé
$\Phi_{K'} \circ \Phi_K : D_{perf}(\mathcal{O}_X) \to D_{perf}(\mathcal{O}_Z)$
est égal à $\Phi_{K''}$ avec $K'' \in D_{perf}(\mathcal{O}_{X \times_S Z})$,
calculé comme dans le lemme \ref{lemma-compose-fourier-mukai},
\item Pour $X$, $Y$, $K$ et $\Phi_K$ comme en (1), si $X \to S$ est de Gorenstein, alors
$\Phi_{K'} : D_{perf}(\mathcal{O}_Y) \to D_{perf}(\mathcal{O}_X)$ est un adjoint
à droite de $\Phi_K$, où $K' \in D_{perf}(\mathcal{O}_{Y \times_S X})$
est l'image inverse de $L\text{pr}_1^*\omega_{X/S}^\bullet
\otimes_{\mathcal{O}_{X \times_S Y}}^\mathbf{L} K^\vee$ par
$Y \times_S X \to X \times_S Y$.
\item Pour $X$, $Y$, $K$ et $\Phi_K$ comme en (1), si $Y \to S$ est de Gorenstein, alors
$\Phi_{K''} : D_{perf}(\mathcal{O}_Y) \to D_{perf}(\mathcal{O}_X)$ est un adjoint
à gauche de $\Phi_K$, où $K'' \in D_{perf}(\mathcal{O}_{Y \times_S X})$
est l'image inverse de $L\text{pr}_2^*\omega_{Y/S}^\bullet
\otimes_{\mathcal{O}_{X \times_S Y}}^\mathbf{L} K^\vee$ par
$Y \times_S X \to X \times_S Y$.
\end{enumerate}
\end{lemma}

\begin{proof}
L'assertion (1) est immédiate d'après le lemme \ref{lemma-fourier-mukai}, assertion (4).

\medskip\noindent
L'assertion (2) résulte du lemme \ref{lemma-compose-fourier-mukai} et du
fait que
$K'' = R\text{pr}_{13, *}(
L\text{pr}_{12}^*K
\otimes_{\mathcal{O}_{X \times_S Y \times_S Z}}^\mathbf{L}
L\text{pr}_{23}^*K')$ est parfait, par exemple d'après
Catégories dérivées des schémas, lemme
\ref{perfect-lemma-flat-proper-perfect-direct-image}.

\medskip\noindent
L'adjonction de l'assertion (3), sur tous les complexes à faisceaux de cohomologie
quasi-cohérents, résulte du lemme \ref{lemma-fourier-mukai-right-adjoint}, avec
$K'$ égal à l'image inverse de
$R\SheafHom_{\mathcal{O}_{X \times_S Y}}(K, a(\mathcal{O}_Y))$
par $Y \times_S X \to X \times_S Y$, où $a$ est l'adjoint à droite
de $R\text{pr}_{2, *} : D_\QCoh(\mathcal{O}_{X \times_S Y}) \to
D_\QCoh(\mathcal{O}_Y)$. Notons $f : X \to S$ le morphisme structural de $X$.
Puisque $f$ est propre, le foncteur
$f^! : D_\QCoh^+(\mathcal{O}_S) \to D_\QCoh^+(\mathcal{O}_X)$
est la restriction à $D_\QCoh^+(\mathcal{O}_S)$
de l'adjoint à droite de
$Rf_* : D_\QCoh(\mathcal{O}_X) \to D_\QCoh(\mathcal{O}_S)$, voir
Dualité pour les schémas, section \ref{duality-section-upper-shriek}.
Ainsi, le complexe dualisant relatif $\omega_{X/S}^\bullet$, tel qu'il est défini dans
Dualité pour les schémas, remarque
\ref{duality-remark-relative-dualizing-complex},
est égal à $\omega_{X/S}^\bullet = f^!\mathcal{O}_S$.
Comme la formation du complexe dualisant relatif
commute au changement de base (voir Dualité pour les schémas, remarque
\ref{duality-remark-relative-dualizing-complex}), on voit que
$a(\mathcal{O}_Y) = L\text{pr}_1^*\omega_{X/S}^\bullet$.
Ainsi
$$
R\SheafHom_{\mathcal{O}_{X \times_S Y}}(K, a(\mathcal{O}_Y))
\cong
L\text{pr}_1^*\omega_{X/S}^\bullet
\otimes_{\mathcal{O}_{X \times_S Y}}^\mathbf{L} K^\vee
$$
d'après Cohomologie, lemme \ref{cohomology-lemma-dual-perfect-complex}.
Enfin, comme $X \to S$ est supposé de Gorenstein, le complexe dualisant relatif
est inversible : cela résulte de Dualité pour les schémas, lemme
\ref{duality-lemma-affine-flat-Noetherian-gorenstein}.
On en déduit que $\omega_{X/S}^\bullet$ est parfait
(Cohomologie, lemme \ref{cohomology-lemma-invertible-derived})
et que $K'$ est donc parfait.
Par conséquent, $\Phi_{K'}$ envoie bien $D_{perf}(\mathcal{O}_Y)$
dans $D_{perf}(\mathcal{O}_X)$, ce qui achève la démonstration de (3).

\medskip\noindent
La démonstration de (4) est la même que celle de (3), si ce n'est qu'on utilise
le lemme \ref{lemma-fourier-mukai-left-adjoint} au lieu du
lemme \ref{lemma-fourier-mukai-right-adjoint}.
\end{proof}














\section{Résolutions et bornes}
\label{section-tricks-smooth}

\noindent
La diagonale d'un schéma lisse et propre possède une résolution commode.

\begin{lemma}
\label{lemma-on-product}
Soit $R$ un anneau noethérien. Soient $X$ et $Y$ des schémas de type fini sur $R$
qui possèdent la propriété de résolution. Pour tout
$\mathcal{O}_{X \times_R Y}$-module cohérent $\mathcal{F}$, il existe
une surjection $\mathcal{E} \boxtimes \mathcal{G} \to \mathcal{F}$,
où $\mathcal{E}$ est un $\mathcal{O}_X$-module localement libre de type fini
et $\mathcal{G}$ un $\mathcal{O}_Y$-module localement libre de type fini.
\end{lemma}

\begin{proof}
Soient $U \subset X$ et $V \subset Y$ des sous-schémas ouverts affines. Soit
$\mathcal{I} \subset \mathcal{O}_X$ le faisceau d'idéaux définissant la
structure réduite de sous-schéma fermé induite sur $X \setminus U$.
De même, soit $\mathcal{I}' \subset \mathcal{O}_Y$ le faisceau d'idéaux définissant la
structure réduite de sous-schéma fermé induite sur $Y \setminus V$.
Alors le faisceau d'idéaux
$$
\mathcal{J} = \Im(\text{pr}_1^*\mathcal{I} \otimes_{\mathcal{O}_{X \times_R Y}}
\text{pr}_2^*\mathcal{I}' \to \mathcal{O}_{X \times_R Y})
$$
vérifie $V(\mathcal{J}) = X \times_R Y \setminus U \times_R V$.
Pour toute section $s \in \mathcal{F}(U \times_R V)$, on peut trouver un entier
$n > 0$ et un morphisme $\mathcal{J}^n \to \mathcal{F}$ dont la restriction à
$U \times_R V$ donne $s$, voir
Cohomologie des schémas, lemme \ref{coherent-lemma-homs-over-open}.
Par hypothèse, on peut choisir des surjections
$\mathcal{E} \to \mathcal{I}$ et $\mathcal{G} \to \mathcal{I}'$.
Elles donnent les surjections correspondantes
$$
\mathcal{E} \boxtimes \mathcal{G} \to \mathcal{J}
\quad\text{et}\quad
\mathcal{E}^{\otimes n} \boxtimes \mathcal{G}^{\otimes n} \to \mathcal{J}^n
$$
et donc un morphisme
$\mathcal{E}^{\otimes n} \boxtimes \mathcal{G}^{\otimes n} \to \mathcal{F}$
dont l'image contient la section $s$ sur $U \times_R V$.
Comme on peut recouvrir $X \times_R Y$ par un nombre fini d'ouverts affines
de la forme $U \times_R V$ et que $\mathcal{F}|_{U \times_R V}$
est engendré par un nombre fini de sections (Propriétés, lemme
\ref{properties-lemma-finite-type-module}),
on en déduit qu'il existe une surjection
$$
\bigoplus\nolimits_{j = 1, \ldots, N}
\mathcal{E}_j^{\otimes n_j} \boxtimes \mathcal{G}_j^{\otimes n_j}
\to \mathcal{F}
$$
où $\mathcal{E}_j$ est localement libre de type fini sur $X$ et
$\mathcal{G}_j$ localement libre de type fini sur $Y$.
En posant $\mathcal{E} = \bigoplus \mathcal{E}_j^{\otimes n_j}$
et $\mathcal{G} = \bigoplus \mathcal{G}_j^{\otimes n_j}$,
on conclut que le lemme est vrai.
\end{proof}

\begin{lemma}
\label{lemma-on-product-general}
Soit $R$ un anneau. Soient $X$ et $Y$ des schémas quasi-compacts et quasi-séparés
sur $R$ qui possèdent la propriété de résolution. Pour tout
$\mathcal{O}_{X \times_R Y}$-module quasi-cohérent de type fini $\mathcal{F}$,
il existe une surjection $\mathcal{E} \boxtimes \mathcal{G} \to \mathcal{F}$,
où $\mathcal{E}$ est un $\mathcal{O}_X$-module localement libre de type fini
et $\mathcal{G}$ un $\mathcal{O}_Y$-module localement libre de type fini.
\end{lemma}

\begin{proof}
Cela résulte du lemme \ref{lemma-on-product} par un argument de passage à la limite.
Nous conseillons vivement au lecteur de ne pas lire la démonstration.
Comme $X \times_R Y$ est un sous-schéma fermé de $X \times_\mathbf{Z} Y$,
on peut sans inconvénient remplacer $R$ par $\mathbf{Z}$.
On peut écrire $\mathcal{F}$ comme quotient d'un
$\mathcal{O}_{X \times_R Y}$-module de présentation finie d'après
Propriétés, lemme
\ref{properties-lemma-finite-directed-colimit-surjective-maps}.
On peut donc supposer que $\mathcal{F}$ est de
présentation finie. Ensuite, on peut écrire $X = \lim X_i$,
où $X_i$ est de présentation finie sur $\mathbf{Z}$, et de même
$Y = \lim Y_j$, voir Limites, proposition \ref{limits-proposition-approximate}.
Alors $\mathcal{F}$ descend en $\mathcal{F}_{ij}$ sur un certain $X_i \times_R Y_j$
(Limites, lemme \ref{limits-lemma-descend-modules-finite-presentation}), tout comme
la propriété de résolution
(Catégories dérivées des schémas, lemme
\ref{perfect-lemma-resolution-property-descends}).
On applique alors le lemme \ref{lemma-on-product}
à $\mathcal{F}_{ij}$, puis on prend l'image inverse.
\end{proof}

\begin{lemma}
\label{lemma-diagonal-resolution}
Soit $R$ un anneau noethérien. Soit $X$ un schéma séparé de type fini
sur $R$ qui possède la propriété de résolution. Posons
$\mathcal{O}_\Delta = \Delta_*(\mathcal{O}_X)$, où
$\Delta : X \to X \times_R X$ est la diagonale de $X/k$.
Il existe une résolution
$$
\ldots \to
\mathcal{E}_2 \boxtimes \mathcal{G}_2 \to
\mathcal{E}_1 \boxtimes \mathcal{G}_1 \to
\mathcal{E}_0 \boxtimes \mathcal{G}_0 \to
\mathcal{O}_\Delta \to 0
$$
où chaque $\mathcal{E}_i$ et chaque $\mathcal{G}_i$ est un
$\mathcal{O}_X$-module localement libre de type fini.
\end{lemma}

\begin{proof}
Comme $X$ est séparé, le morphisme diagonal $\Delta$ est une immersion
fermée, et $\mathcal{O}_\Delta$ est donc un
$\mathcal{O}_{X \times_R X}$-module cohérent (Cohomologie des schémas, lemme
\ref{coherent-lemma-i-star-equivalence}).
Le lemme résulte donc immédiatement du lemme \ref{lemma-on-product}.
\end{proof}

\begin{lemma}
\label{lemma-Ext-0-regular}
Soit $X$ un schéma noethérien régulier de dimension $d < \infty$. Alors
\begin{enumerate}
\item pour $\mathcal{F}$, $\mathcal{G}$ qui sont des modules cohérents sur $\mathcal{O}_X$,
on a $\Ext^n_X(\mathcal{F}, \mathcal{G}) = 0$ pour $n > d$, et
\item pour $K, L \in D^b_{\textit{Coh}}(\mathcal{O}_X)$ et $a \in \mathbf{Z}$,
si $H^i(K) = 0$ pour $i < a + d$ et $H^i(L) = 0$ pour $i \geq a$, alors
$\Hom_X(K, L) = 0$.
\end{enumerate}
\end{lemma}

\begin{proof}
Pour démontrer (1), on utilise la suite spectrale
$$
H^p(X, \SheafExt^q(\mathcal{F}, \mathcal{G})) \Rightarrow
\Ext^{p + q}_X(\mathcal{F}, \mathcal{G})
$$
de Cohomologie, section \ref{cohomology-section-ext}. Soit $x \in X$.
On a
$$
\SheafExt^q(\mathcal{F}, \mathcal{G})_x =
\SheafExt^q_{\mathcal{O}_{X, x}}(\mathcal{F}_x, \mathcal{G}_x)
$$
voir Cohomologie, lemme \ref{cohomology-lemma-stalk-internal-hom}
(on utilise également ici le fait que $\mathcal{F}$ est pseudo-cohérent d'après
Catégories dérivées des schémas, lemme
\ref{perfect-lemma-identify-pseudo-coherent-noetherian}).
Posons $d_x = \dim(\mathcal{O}_{X, x})$.
Comme $\mathcal{O}_{X, x}$ est régulier, l'anneau
$\mathcal{O}_{X, x}$ est de dimension globale $d_x$, voir
Algèbre, proposition \ref{algebra-proposition-regular-finite-gl-dim}.
Ainsi, $\SheafExt^q_{\mathcal{O}_{X, x}}(\mathcal{F}_x, \mathcal{G}_x)$
est nul pour $q > d_x$. Il s'ensuit que les modules
$\SheafExt^q(\mathcal{F}, \mathcal{G})$ ont un support
de dimension au plus $d - q$. On a donc
$H^p(X, \SheafExt^q(\mathcal{F}, \mathcal{G})) = 0$ pour $p > d - q$
d'après Cohomologie, proposition \ref{cohomology-proposition-vanishing-Noetherian}.
Ceci prouve (1).

\medskip\noindent
Démonstration de (2).
On peut raisonner par récurrence sur le nombre de faisceaux de cohomologie non nuls
de $K$ et $L$. Le cas où ces nombres valent $0, 1$ résulte
de (1). Si le nombre de faisceaux de cohomologie non nuls de $K$
est $> 1$, soit $i \in \mathbf{Z}$ minimal tel que
$H^i(K)$ soit non nul. On obtient un triangle distingué
$$
H^i(K)[-i] \to K \to \tau_{\geq i + 1}K
$$
(Catégories dérivées, remarque
\ref{derived-remark-truncation-distinguished-triangle}),
et la nullité de $\Hom(K, L)$ résulte de celle
de $\Hom(H^i(K)[-i], L)$ et $\Hom(\tau_{\geq i + 1}K, L)$
d'après Catégories dérivées, lemme \ref{derived-lemma-representable-homological}.
De même si $L$ possède plus d'un faisceau de cohomologie non nul.
\end{proof}

\begin{lemma}
\label{lemma-split-complex-regular}
Soit $X$ un schéma noethérien régulier de dimension $d < \infty$.
Soient $K \in D^b_{\textit{Coh}}(\mathcal{O}_X)$ et $a \in \mathbf{Z}$.
Si $H^i(K) = 0$ pour $a < i < a + d$, alors
$K = \tau_{\leq a}K \oplus \tau_{\geq a + d}K$.
\end{lemma}

\begin{proof}
On a $\tau_{\leq a}K = \tau_{\leq a + d - 1}K$ par la nullité supposée
des faisceaux de cohomologie. D'après Catégories dérivées, remarque
\ref{derived-remark-truncation-distinguished-triangle},
on a un triangle distingué
$$
\tau_{\leq a}K \to K \to \tau_{\geq a + d}K \xrightarrow{\delta}
(\tau_{\leq a}K)[1]
$$
D'après Catégories dérivées, lemme \ref{derived-lemma-split}, il
suffit de montrer que le morphisme $\delta$ est nul.
Cela résulte du lemme \ref{lemma-Ext-0-regular}.
\end{proof}

\begin{lemma}
\label{lemma-diagonal-trick}
Soit $k$ un corps. Soit $X$ un schéma quasi-compact, séparé et lisse sur $k$.
Il existe des $\mathcal{O}_X$-modules localement libres de type fini
$\mathcal{E}$ et $\mathcal{G}$ tels que
$$
\mathcal{O}_\Delta \in \langle \mathcal{E} \boxtimes \mathcal{G} \rangle
$$
dans $D(\mathcal{O}_{X \times X})$, où la notation est celle de
Catégories dérivées, section \ref{derived-section-generators}.
\end{lemma}

\begin{proof}
Rappelons que $X$ est régulier d'après
Variétés, lemme \ref{varieties-lemma-smooth-regular}.
Ainsi, $X$ possède la propriété de résolution d'après
Catégories dérivées des schémas, lemme
\ref{perfect-lemma-regular-resolution-property}.
On peut donc choisir une résolution comme dans le lemme \ref{lemma-diagonal-resolution}.
Posons $\dim(X) = d$. Comme $X \times X$ est lisse sur $k$, il est régulier.
Ainsi, $X \times X$ est un schéma noethérien régulier et l'on a
$\dim(X \times X) = 2d$. L'objet
$$
K = (\mathcal{E}_{2d} \boxtimes \mathcal{G}_{2d} \to
\ldots \to
\mathcal{E}_0 \boxtimes \mathcal{G}_0)
$$
de $D_{perf}(\mathcal{O}_{X \times X})$
a pour faisceaux de cohomologie $\mathcal{O}_\Delta$
en degré $0$ et $\Ker(\mathcal{E}_{2d} \boxtimes \mathcal{G}_{2d} \to
\mathcal{E}_{2d-1} \boxtimes \mathcal{G}_{2d-1})$ en degré $-2d$, et zéro
à tous les autres degrés.
Le lemme \ref{lemma-split-complex-regular} montre donc que
$\mathcal{O}_\Delta$ est un facteur direct de $K$ dans
$D_{perf}(\mathcal{O}_{X \times X})$.
Manifestement, l'objet $K$ appartient à
$$
\left\langle
\bigoplus\nolimits_{i = 0, \ldots, 2d} \mathcal{E}_i \boxtimes \mathcal{G}_i
\right\rangle
\subset
\left\langle
\left(\bigoplus\nolimits_{i = 0, \ldots, 2d} \mathcal{E}_i\right)
\boxtimes
\left(\bigoplus\nolimits_{i = 0, \ldots, 2d} \mathcal{G}_i\right)
\right\rangle
$$
ce qui achève la démonstration. (Le lecteur peut consulter
Catégories dérivées, lemmes \ref{derived-lemma-generated-by-E-explicit} et
\ref{derived-lemma-in-cone-n}, pour voir que notre objet appartient à cette
catégorie.)
\end{proof}

\begin{lemma}
\label{lemma-smooth-proper-strong-generator}
Soit $k$ un corps. Soit $X$ un schéma propre et lisse sur $k$.
Alors $D_{perf}(\mathcal{O}_X)$
possède un générateur fort.
\end{lemma}

\begin{proof}
À l'aide du lemme \ref{lemma-diagonal-trick}, choisissons des
$\mathcal{O}_X$-modules localement libres de type fini $\mathcal{E}$ et $\mathcal{G}$ tels que
$\mathcal{O}_\Delta \in \langle \mathcal{E} \boxtimes \mathcal{G} \rangle$
dans $D(\mathcal{O}_{X \times X})$. Nous affirmons que $\mathcal{G}$
est un générateur fort de $D_{perf}(\mathcal{O}_X)$. Avec les notations de
Catégories dérivées, section \ref{derived-section-operate-on-full},
choisissons $m, n \geq 1$ tels que
$$
\mathcal{O}_\Delta \in
smd(add(\mathcal{E} \boxtimes \mathcal{G}[-m, m])^{\star n})
$$
C'est possible d'après Catégories dérivées, lemme
\ref{derived-lemma-find-smallest-containing-E}.
Soit $K$ un objet de $D_{perf}(\mathcal{O}_X)$. Comme
$L\text{pr}_1^*K \otimes_{\mathcal{O}_{X \times X}}^\mathbf{L} -$
est un foncteur exact et que
$$
L\text{pr}_1^*K \otimes_{\mathcal{O}_{X \times X}}^\mathbf{L}
(\mathcal{E} \boxtimes \mathcal{G}) =
(K \otimes_{\mathcal{O}_X}^\mathbf{L} \mathcal{E}) \boxtimes \mathcal{G}
$$
on déduit de
Catégories dérivées, remarque \ref{derived-remark-operations-functor}, que
$$
L\text{pr}_1^*K
\otimes_{\mathcal{O}_{X \times X}}^\mathbf{L}
\mathcal{O}_\Delta
\in
smd(add(
(K \otimes_{\mathcal{O}_X}^\mathbf{L} \mathcal{E})
\boxtimes \mathcal{G}[-m, m])^{\star n})
$$
En appliquant le foncteur exact $R\text{pr}_{2, *}$ et en observant que
$$
R\text{pr}_{2, *}
\left((K \otimes_{\mathcal{O}_X}^\mathbf{L} \mathcal{E}) \boxtimes
\mathcal{G}\right) =
R\Gamma(X, K \otimes_{\mathcal{O}_X}^\mathbf{L} \mathcal{E})
\otimes_k \mathcal{G}
$$
d'après Catégories dérivées des schémas, lemme
\ref{perfect-lemma-cohomology-base-change}, on en déduit que
$$
K = R\text{pr}_{2, *}(L\text{pr}_1^*K
\otimes_{\mathcal{O}_{X \times X}}^\mathbf{L} \mathcal{O}_\Delta)
\in
smd(add(R\Gamma(X, K \otimes_{\mathcal{O}_X}^\mathbf{L} \mathcal{E})
\otimes_k \mathcal{G}[-m, m])^{\star n})
$$
L'égalité résulte de la discussion de
l'exemple \ref{example-diagonal-fourier-mukai}.
Comme $K$ est parfait, il existe $a \leq b$ tels que
$H^i(X, K)$ ne soit non nul que pour $i \in [a, b]$. Comme $X$ est propre,
chaque $H^i(X, K)$ est de dimension finie. On en déduit que
le membre de droite est contenu dans
$smd(add(\mathcal{G}[-m + a, m + b])^{\star n})$, qui est
lui-même contenu dans $\langle \mathcal{G} \rangle_n$ d'après l'une des
références données ci-dessus. Cela achève la démonstration.
\end{proof}

\begin{lemma}
\label{lemma-diagonal-trick-proper}
Soit $k$ un corps. Soit $X$ un schéma propre et lisse sur $k$.
Il existe des entiers $m, n \geq 1$ et un
$\mathcal{O}_X$-module localement libre de type fini $\mathcal{G}$ tels que tout
$\mathcal{O}_X$-module cohérent appartienne à $smd(add(\mathcal{G}[-m, m])^{\star n})$,
avec les notations de Catégories dérivées, section
\ref{derived-section-operate-on-full}.
\end{lemma}

\begin{proof}
Dans la démonstration du lemme \ref{lemma-smooth-proper-strong-generator},
nous avons montré qu'il existe $m', n \geq 1$ tels que, pour tout
$\mathcal{O}_X$-module cohérent $\mathcal{F}$,
$$
\mathcal{F} \in smd(add(\mathcal{G}[-m' + a, m' + b])^{\star n})
$$
pour tous $a \leq b$ tels que $H^i(X, \mathcal{F})$ ne soit non nul que
pour $i \in [a, b]$. On peut donc prendre $a = 0$ et $b = \dim(X)$.
Prendre $m = \max(m', m' + b)$ achève la démonstration.
\end{proof}

\noindent
Le lemme suivant est le résultat de bornitude auquel renvoie
le titre de cette section.

\begin{lemma}
\label{lemma-boundedness}
Soit $k$ un corps. Soit $X$ un schéma lisse et propre sur $k$.
Soit $\mathcal{A}$ une catégorie abélienne. Soit
$H : D_{perf}(\mathcal{O}_X) \to \mathcal{A}$ un foncteur homologique
(Catégories dérivées, définition \ref{derived-definition-homological})
tel que, pour tout $K$ dans $D_{perf}(\mathcal{O}_X)$, l'objet
$H^i(K)$ ne soit non nul que pour un nombre fini de $i \in \mathbf{Z}$.
Alors il existe un entier $m \geq 1$ tel que
$H^i(\mathcal{F}) = 0$ pour tout $\mathcal{O}_X$-module cohérent
$\mathcal{F}$ et tout $i \not \in [-m, m]$.
De même pour les foncteurs cohomologiques.
\end{lemma}

\begin{proof}
Combiner le lemme \ref{lemma-diagonal-trick-proper} avec
Catégories dérivées, lemme \ref{derived-lemma-forward-cone-n}.
\end{proof}

\begin{lemma}
\label{lemma-bounded-fibres}
Soit $k$ un corps. Soient $X$ et $Y$ des schémas de type fini sur $k$.
Soit $K_0 \to K_1 \to K_2 \to \ldots$ un système d'objets
de $D_{perf}(\mathcal{O}_{X \times Y})$, et soit $m \geq 0$ un entier tel que
\begin{enumerate}
\item $H^q(K_i)$ ne soit non nul que pour $q \leq m$,
\item pour tout $\mathcal{O}_X$-module cohérent $\mathcal{F}$ tel que
$\dim(\text{Supp}(\mathcal{F})) = 0$, l'objet
$$
R\text{pr}_{2, *}(
\text{pr}_1^*\mathcal{F} \otimes_{\mathcal{O}_{X \times Y}}^\mathbf{L}
K_n)
$$
n'ait de faisceaux de cohomologie non nuls qu'aux degrés appartenant à
$[-m, m] \cup [-m - n, m - n]$ et, pour $n > 2m$, les morphismes de transition
induisent des isomorphismes sur les faisceaux de cohomologie aux degrés appartenant à $[-m, m]$.
\end{enumerate}
Alors $K_n$ n'a de faisceaux de cohomologie non nuls qu'aux degrés appartenant à
$[-m, m] \cup [-m - n, m - n]$ et, pour $n > 2m$, les
morphismes de transition induisent des isomorphismes sur les faisceaux de cohomologie aux degrés appartenant à
$[-m, m]$. De plus, si $X$ et $Y$ sont lisses sur $k$, alors, pour $n$
assez grand, on a $K_n = K \oplus C_n$ dans
$D_{perf}(\mathcal{O}_{X \times Y})$,
où $K$ n'a de cohomologie qu'aux degrés $[-m, m]$ et $C_n$ qu'aux
degrés $[-m - n, m - n]$, et les morphismes de transition
définissent des isomorphismes entre les diverses copies de $K$.
\end{lemma}

\begin{proof}
Soit $Z$ le support schématique d'un module $\mathcal{F}$ comme en (2).
Alors $Z \to \Spec(k)$ est fini, donc $Z \times Y \to Y$ est fini.
Il s'ensuit que, pour un objet $M$ de $D_\QCoh(\mathcal{O}_{X \times Y})$
dont les faisceaux de cohomologie sont supportés par $Z \times Y$, on a
$H^i(R\text{pr}_{2, *}(M)) = \text{pr}_{2, *}H^i(M)$, et le foncteur
$\text{pr}_{2, *}$ est fidèle sur les modules quasi-cohérents supportés
par $Z \times Y$ ; nous omettons les détails. On voit donc que les objets
$$
\text{pr}_1^*\mathcal{F} \otimes_{\mathcal{O}_{X \times Y}}^\mathbf{L} K_n
$$
de $D_{perf}(\mathcal{O}_{X \times Y})$ n'ont de faisceaux de cohomologie non nuls
qu'aux degrés appartenant à $[-m, m] \cup [-m - n, m - n]$ et que, pour $n > 2m$, les morphismes de transition
induisent des isomorphismes sur les faisceaux de cohomologie aux degrés appartenant à $[-m, m]$.
Soit $z \in X \times Y$ un point fermé dont l'image est le point fermé
$x \in X$. Alors on sait que
$$
K_{n, z} \otimes_{\mathcal{O}_{X \times Y, z}}^\mathbf{L}
\mathcal{O}_{X \times Y, z}/\mathfrak m_x^t\mathcal{O}_{X \times Y, z}
$$
n'a de cohomologie non nulle que dans les intervalles
$[-m, m] \cup [-m - n, m - n]$.
On déduit de Compléments sur l'algèbre, lemme
\ref{more-algebra-lemma-kollar-kovacs-pseudo-coherent},
que $K_{n, z}$ n'a de cohomologie non nulle
qu'aux degrés $[-m, m] \cup [-m - n, m - n]$. Comme cela vaut en tout
point fermé de $X \times Y$, on en déduit que $K_n$ n'a de faisceaux de cohomologie non nuls
qu'aux degrés $[-m, m] \cup [-m - n, m - n]$.
De même, on voit que les morphismes $K_n \to K_{n + 1}$
sont des isomorphismes sur les faisceaux de cohomologie aux degrés appartenant à $[-m, m]$
pour $n > 2m$.

\medskip\noindent
Si $X$ et $Y$ sont lisses sur $k$, alors $X \times Y$ est lisse
sur $k$, donc régulier d'après
Variétés, lemme \ref{varieties-lemma-smooth-regular}.
Nous obtiendrons donc la décomposition en somme directe de $K_n$
dès que $n > 2m + \dim(X \times Y)$, d'après
le lemme \ref{lemma-split-complex-regular}. La dernière assertion
en résulte clairement.
\end{proof}





\section{Foncteurs frères}
\label{section-sibling}

\noindent
Dans cette section, nous démontrons un résultat catégorique portant sur la notion suivante.

\begin{definition}
\label{definition-siblings}
Soit $\mathcal{A}$ une catégorie abélienne. Soit $\mathcal{D}$ une
catégorie triangulée. Nous disons que deux foncteurs exacts de catégories triangulées
$$
F, F' : D^b(\mathcal{A}) \longrightarrow \mathcal{D}
$$
sont des {\it frères}, ou que $F'$ est un {\it frère} de $F$,
si les deux conditions suivantes sont satisfaites
\begin{enumerate}
\item les foncteurs $F \circ i$ et $F' \circ i$ sont isomorphes,
où $i : \mathcal{A} \to D^b(\mathcal{A})$ est le foncteur d'inclusion, et
\item $F(K) \cong F'(K)$ pour tout $K$ de $D^b(\mathcal{A})$.
\end{enumerate}
\end{definition}

\noindent
La seconde condition est parfois une conséquence de la première.

\begin{lemma}
\label{lemma-sibling-fully-faithful}
Soit $\mathcal{A}$ une catégorie abélienne. Soit $\mathcal{D}$ une
catégorie triangulée. Soient
$F, F' : D^b(\mathcal{A}) \longrightarrow \mathcal{D}$
des foncteurs exacts de catégories triangulées. Supposons que
\begin{enumerate}
\item les foncteurs $F \circ i$ et $F' \circ i$ sont isomorphes,
où $i : \mathcal{A} \to D^b(\mathcal{A})$ est le foncteur d'inclusion, et
\item pour tous $X, Y \in \Ob(\mathcal{A})$, nous avons
$\Ext^q_\mathcal{D}(F(X), F(Y)) = 0$ pour $q < 0$ (par exemple
si $F$ est pleinement fidèle).
\end{enumerate}
Alors $F$ et $F'$ sont frères.
\end{lemma}

\begin{proof}
Soit $K \in D^b(\mathcal{A})$. Nous allons montrer que $F(K)$ est isomorphe à $F'(K)$.
On peut représenter $K$ par un complexe borné $A^\bullet$ d'objets de
$\mathcal{A}$. Après avoir remplacé $K$ par un décalé, on peut
supposer que $A^i = 0$ pour $i > 0$. Choisissons $n \geq 0$ tel que $A^{-i} = 0$
pour $i > n$. Les objets
$$
M_i = (A^{-i} \to \ldots \to A^0)[-i],\quad i = 0, \ldots, n
$$
forment un système de Postnikov dans $D^b(\mathcal{A})$ pour le complexe
$A^\bullet = A^{-n} \to \ldots \to A^0$ de $D^b(\mathcal{A})$.
Voir Catégories dérivées, exemple \ref{derived-example-key-postnikov}.
Comme $F$ et $F'$ sont tous deux des foncteurs exacts de catégories triangulées,
$$
F(M_i)
\quad\text{et}\quad
F'(M_i)
$$
forment tous deux un système de Postnikov dans $\mathcal{D}$ pour le complexe
$$
F(A^{-n}) \to \ldots \to F(A^0) =
F'(A^{-n}) \to \ldots \to F'(A^0)
$$
Puisque tous les $\Ext$ négatifs entre ces objets sont nuls par hypothèse,
l'unicité des systèmes de Postnikov permet de conclure
(Catégories dérivées, lemme \ref{derived-lemma-existence-postnikov-system})
que $F(K) = F(M_n[n]) \cong F'(M_n[n]) = F'(K)$.
\end{proof}

\begin{lemma}
\label{lemma-sibling-faithful}
Soient $F$ et $F'$ des frères comme dans la définition \ref{definition-siblings}.
Alors
\begin{enumerate}
\item si $F$ est essentiellement surjectif, alors $F'$ est essentiellement
surjectif,
\item si $F$ est pleinement fidèle, alors $F'$ est pleinement fidèle.
\end{enumerate}
\end{lemma}

\begin{proof}
L'assertion (1) est immédiate d'après la propriété (2) des foncteurs frères.

\medskip\noindent
Supposons $F$ pleinement fidèle. Notons $\mathcal{D}' \subset \mathcal{D}$
l'image essentielle de $F$, de sorte que $F : D^b(\mathcal{A}) \to \mathcal{D}'$
soit une équivalence. Puisque le foncteur $F'$ se factorise par $\mathcal{D}'$
d'après la propriété (2) des foncteurs frères, on peut considérer le foncteur
$H = F^{-1} \circ F' : D^b(\mathcal{A}) \to D^b(\mathcal{A})$.
Observons que $H$ est un frère du foncteur identité.
Comme il suffit de prouver que $H$ est pleinement fidèle,
on se ramène au problème examiné dans le paragraphe suivant.

\medskip\noindent
Posons $\mathcal{D} = D^b(\mathcal{A})$. Il faut montrer qu'un frère
$F : \mathcal{D} \to \mathcal{D}$ du foncteur identité est pleinement fidèle.
Notons $a_X : X \to F(X)$ l'isomorphisme fonctoriel pour
$X \in \Ob(\mathcal{A})$ fourni par la définition \ref{definition-siblings}.
Pour tout $K$ de $\mathcal{D}$ et tout triangle distingué
$K_1 \to K_2 \to K_3$ de $\mathcal{D}$,
si les applications
$$
F : \Hom(K, K_i[n]) \to \Hom(F(K), F(K_i[n]))
$$
sont des isomorphismes pour tout $n \in \mathbf{Z}$ et $i = 1, 3$, alors il en
est de même pour $i = 2$ et tout $n \in \mathbf{Z}$. On utilise ici
le $5$-lemme, Homologie, lemme \ref{homology-lemma-five-lemma}, et
Catégories dérivées, lemme \ref{derived-lemma-representable-homological} ;
nous omettons les détails. De même, si les applications
$$
F : \Hom(K_i[n], K) \to \Hom(F(K_i[n]), F(K))
$$
sont des isomorphismes pour tout $n \in \mathbf{Z}$ et $i = 1, 3$, alors il en
est de même pour $i = 2$ et tout $n \in \mathbf{Z}$. En utilisant les troncations
canoniques et une récurrence sur le nombre d'objets de cohomologie non nuls,
on voit qu'il suffit de montrer que
$$
F : \Ext^q(X, Y) \to \Ext^q(F(X), F(Y))
$$
est bijective pour tous $X, Y \in \Ob(\mathcal{A})$ et tout $q \in \mathbf{Z}$.
Comme $F$ est un frère de $\text{id}$, on a $F(X) \cong X$ et
$F(Y) \cong Y$ ; le membre de droite est donc nul pour $q < 0$.
Le cas $q = 0$ résulte de l'hypothèse que $F$ est un frère du
foncteur identité. Il reste à traiter les cas $q > 0$.

\medskip\noindent
Le cas $q = 1$ : injectivité. Un élément $\xi$ de $\Ext^1(X, Y)$
donne lieu à un triangle distingué
$$
Y \to E \to X \xrightarrow{\xi} Y[1]
$$
Observons que $E \in \Ob(\mathcal{A})$. Comme $F$ est un frère du
foncteur identité, on obtient un diagramme commutatif
$$
\xymatrix{
E \ar[d] \ar[r] & X \ar[d] \\
F(E) \ar[r] & F(X)
}
$$
dont les flèches verticales sont les isomorphismes $a_E$ et $a_X$.
Par TR3, le triangle distingué associé à $\xi$ dont on est parti
est isomorphe au triangle distingué
$$
F(Y) \to F(E) \to F(X) \xrightarrow{F(\xi)} F(Y[1]) = F(Y)[1]
$$
Ainsi, $\xi = 0$ si et seulement si $F(\xi)$ est nul, c'est-à-dire que
$F : \Ext^1(X, Y) \to \Ext^1(F(X), F(Y))$ est injective.

\medskip\noindent
Le cas $q = 1$ : surjectivité. Soit $\theta$ un élément de
$\Ext^1(F(X), F(Y))$. Il définit une extension de $F(X)$ par $F(Y)$
dans $\mathcal{A}$, que l'on peut écrire sous la forme $F(E)$
puisque $F$ est un frère du foncteur identité. On obtient ainsi un triangle
distingué
$$
F(Y) \xrightarrow{F(\alpha)} F(E)
\xrightarrow{F(\beta)} F(X)
\xrightarrow{\theta} F(Y[1]) = F(Y)[1]
$$
pour certains morphismes $\alpha : Y \to E$ et $\beta : E \to X$.
Comme $F$ est un frère du foncteur identité, la suite
$0 \to Y \to E \to X \to 0$
est une suite exacte courte de $\mathcal{A}$ ! On obtient donc un
triangle distingué
$$
Y \xrightarrow{\alpha} E \xrightarrow{\beta} X \xrightarrow{\delta} Y[1]
$$
pour un certain morphisme $\delta : X \to Y[1]$. En appliquant le foncteur exact
$F$, on obtient le triangle distingué
$$
F(Y) \xrightarrow{F(\alpha)} F(E) \xrightarrow{F(\beta)} F(X)
\xrightarrow{F(\delta)} F(Y)[1]
$$
En raisonnant comme ci-dessus, on voit que ces triangles sont isomorphes.
Il existe donc un diagramme commutatif
$$
\xymatrix{
F(X) \ar[d]^\gamma \ar[r]_{F(\delta)} & F(Y[1]) \ar[d]_\epsilon \\
F(X) \ar[r]^\theta & F(Y[1])
}
$$
pour certains isomorphismes $\gamma$, $\epsilon$ (on pourrait être plus précis, mais cela
ne sera pas nécessaire). On peut écrire $\gamma = F(\gamma')$ et
$\epsilon = F(\epsilon')$. On a alors
$\theta = F(\epsilon' \circ \delta \circ (\gamma')^{-1})$,
ce qui prouve la surjectivité.

\medskip\noindent
Le cas $q > 1$ : surjectivité. En utilisant les extensions de Yoneda, voir
Catégories dérivées, section \ref{derived-section-ext}, on voit que, pour tout
élément $\xi$ de $\Ext^q(F(X), F(Y))$, il existe
$F(X) = B_0, B_1, \ldots, B_{q - 1}, B_q = F(Y) \in \Ob(\mathcal{A})$ et des
éléments
$$
\xi_i \in \Ext^1(B_{i - 1}, B_i)
$$
tels que $\xi$ soit le composé $\xi_q \circ \ldots \circ \xi_1$.
Écrivons $B_i = F(A_i)$ (bien entendu, on a $A_i = B_i$, mais il n'est pas
nécessaire de l'utiliser), de sorte que
$$
\xi_i = F(\eta_i) \in \Ext^1(F(A_{i - 1}), F(A_i))
\quad\text{avec}\quad
\eta_i \in \Ext^1(A_{i - 1}, A_i)
$$
par surjectivité pour $q = 1$. Alors $\eta = \eta_q \circ \ldots \circ \eta_1$
est un élément de $\Ext^q(X, Y)$ tel que $F(\eta) = \xi$.

\medskip\noindent
Le cas $q > 1$ : injectivité. Un élément $\xi$ de $\Ext^q(X, Y)$
donne lieu à un triangle distingué
$$
Y[q - 1] \to E \to X \xrightarrow{\xi} Y[q]
$$
En appliquant $F$, on obtient un triangle distingué
$$
F(Y)[q - 1] \to F(E) \to F(X) \xrightarrow{F(\xi)} F(Y)[q]
$$
Si $F(\xi) = 0$, alors $F(E) \cong F(Y)[q - 1] \oplus F(X)$
dans $\mathcal{D}$, voir
Catégories dérivées, lemme \ref{derived-lemma-split}.
Comme $F$ est un frère du foncteur identité, on a
$E \cong F(E)$, et donc
$$
E \cong F(E) \cong F(Y)[q - 1] \oplus F(X) \cong Y[q - 1] \oplus X
$$
Autrement dit, $E$ est isomorphe à la
somme directe de ses objets de cohomologie. Cela implique que le
triangle distingué initial est scindé, c'est-à-dire que $\xi = 0$.
\end{proof}

\noindent
Donnons une définition non standard. Soit $\mathcal{A}$ une catégorie
abélienne. Disons que $\mathcal{A}$ {\it possède assez d'objets négatifs}
si, pour tout $X \in \Ob(\mathcal{A})$, il existe un objet $N$ tel que
\begin{enumerate}
\item il existe une surjection $N \to X$, et
\item $\Hom(X, N) = 0$.
\end{enumerate}
Démontrons deux lemmes sur cette notion afin de faciliter
la démonstration de la proposition \ref{proposition-siblings-isomorphic}.

\begin{lemma}
\label{lemma-good-map}
Soit $\mathcal{A}$ une catégorie abélienne possédant assez d'objets négatifs.
Soit $X \in D^b(\mathcal{A})$. Soit $b \in \mathbf{Z}$ tel que
$H^i(X) = 0$ pour $i > b$. Alors
il existe un morphisme $N[-b] \to X$ tel que le morphisme induit
$N \to H^b(X)$ soit surjectif et que $\Hom(H^b(X), N) = 0$.
\end{lemma}

\begin{proof}
À l'aide des foncteurs de troncation, on peut représenter $X$ par un complexe
$A^a \to A^{a + 1} \to \ldots \to A^b$ d'objets de $\mathcal{A}$.
Choisissons $N$ dans $\mathcal{A}$ de sorte qu'il existe une surjection
$t : N \to A^b$ et que $\Hom(A^b, N) = 0$. Alors la surjection $t$
définit un morphisme $N[-b] \to X$ comme voulu.
\end{proof}

\begin{lemma}
\label{lemma-good-map-zero}
Soit $\mathcal{A}$ une catégorie abélienne possédant assez d'objets négatifs.
Soit $f : X \to X'$ un morphisme de $D^b(\mathcal{A})$. Soit $b \in \mathbf{Z}$
tel que $H^i(X) = 0$ pour $i > b$ et $H^i(X') = 0$ pour $i \geq b$.
Alors il existe un morphisme $N[-b] \to X$ tel que le morphisme induit
$N \to H^b(X)$ soit surjectif, que $\Hom(H^b(X), N) = 0$, et
que le composé $N[-b] \to X \to X'$ soit nul.
\end{lemma}

\begin{proof}
On peut représenter $f$ par un morphisme $f^\bullet : A^\bullet \to B^\bullet$
de complexes bornés d'objets de $\mathcal{A}$, voir par exemple
Catégories dérivées, lemme \ref{derived-lemma-bounded-derived}.
Considérons l'objet
$$
C = \Ker(A^b \to A^{b + 1}) \times_{\Ker(B^b \to B^{b + 1})} B^{b - 1}
$$
de $\mathcal{A}$. Comme $H^b(B^\bullet) = 0$, on voit que
$C \to H^b(A^\bullet)$ est surjectif. D'autre part, le morphisme
$C \to A^b \to B^b$ est le même que le morphisme $C \to B^{b - 1} \to B^b$,
et le composé $C[-b] \to X \to X'$ est donc nul.
Puisque $\mathcal{A}$ possède assez d'objets négatifs, on peut trouver un objet $N$
muni d'une surjection $N \to C \oplus H^b(X)$ tel que
$\Hom(C \oplus H^b(X), N) = 0$. Alors $N$, avec le morphisme
$N[-b] \to X$, est une solution au problème posé par le lemme.
\end{proof}

\noindent
Nous encourageons le lecteur à lire l'original
\cite[Proposition 2.16]{Orlov-K3} pour les idées remarquables
qui interviennent dans la démonstration de la proposition suivante.

\begin{proposition}
\label{proposition-siblings-isomorphic}
\begin{reference}
\cite[Proposition 2.16]{Orlov-K3} ; le fait qu'il ne soit pas nécessaire
de supposer la nullité de $\Ext^q(N, X)$ pour $q > 0$ dans la définition
des objets négatifs ci-dessus est dû à \cite{Canonaco-Stellari}.
\end{reference}
Soient $F$ et $F'$ des frères comme dans la définition \ref{definition-siblings}.
Supposons que $F$ soit pleinement fidèle et que $\mathcal{A}$ possède assez
d'objets négatifs (voir ci-dessus). Alors $F$ et $F'$ sont des foncteurs isomorphes.
\end{proposition}

\begin{proof}
D'après l'assertion (2) de la définition \ref{definition-siblings}, l'image du foncteur
$F'$ est contenue dans l'image essentielle du foncteur $F$. Ainsi,
le foncteur $H = F^{-1} \circ F'$ est un frère du foncteur identité.
Cela nous ramène au cas décrit dans le paragraphe suivant.

\medskip\noindent
Soit $\mathcal{D} = D^b(\mathcal{A})$. Il faut montrer qu'un frère
$F : \mathcal{D} \to \mathcal{D}$ du foncteur identité est
isomorphe au foncteur identité. Pour un objet $X$ de $\mathcal{D}$,
disons que $X$ est de {\it largeur} $w = w(X)$ si $w \geq 0$ est minimal
parmi les entiers tels qu'il existe $a \in \mathbf{Z}$ avec $H^i(X) = 0$
pour $i \not \in [a, a + w - 1]$. Puisque $F$ est un frère du foncteur identité
et que $F \circ [n] = [n] \circ F$, on dispose déjà d'isomorphismes
$$
c_X : X \to F(X)
$$
pour $w(X) \leq 1$, compatibles aux décalages. De plus, si $X = A[-a]$ et
$X' = A'[-a]$ pour certains $A, A' \in \Ob(\mathcal{A})$, alors, pour tout morphisme
$f : X \to X'$, le diagramme
\begin{equation}
\label{equation-to-show}
\vcenter{
\xymatrix{
X \ar[d]_{c_X} \ar[r]_f &
X' \ar[d]^{c_{X'}} \\
F(X) \ar[r]^{F(f)} &
F(X')
}
}
\end{equation}
est commutatif.

\medskip\noindent
Montrons ensuite que, pour tout morphisme $f : X \to X'$ tel que
$w(X), w(X') \leq 1$, le diagramme (\ref{equation-to-show}) est commutatif.
Si $X$ ou $X'$ est nul, c'est clair. Sinon, on peut écrire
$X = A[-a]$ et $X' = A'[-a']$ pour d'uniques $A, A'$ dans $\mathcal{A}$
et $a, a' \in \mathbf{Z}$. Le cas $a = a'$ a été examiné ci-dessus.
Si $a' > a$, alors $f = 0$ (Catégories dérivées, lemme
\ref{derived-lemma-negative-exts}), et le résultat est clair.
Si $a' < a$, alors $f$ correspond à un élément
$\xi \in \Ext^q(A, A')$ avec $q = a - a'$. En utilisant les extensions de Yoneda, voir
Catégories dérivées, section \ref{derived-section-ext}, on peut trouver
$A = A_0, A_1, \ldots, A_{q - 1}, A_q = A' \in \Ob(\mathcal{A})$ et des
éléments
$$
\xi_i \in \Ext^1(A_{i - 1}, A_i)
$$
tels que $\xi$ soit le composé $\xi_q \circ \ldots \circ \xi_1$.
Autrement dit, en posant $X_i = A_i[-a + i]$,
on obtient des morphismes
$$
X = X_0 \xrightarrow{f_1} X_1 \to \ldots \to X_{q - 1}
\xrightarrow{f_q} X_q = X'
$$
dont le composé est $f$. Comme la commutativité de (\ref{equation-to-show})
pour $f_1, \ldots, f_q$ implique celle pour $f$, on se ramène au cas $q = 1$.
Dans ce cas, après décalage, on peut supposer que l'on a un triangle distingué
$$
A' \to E \to A \xrightarrow{f} A'[1]
$$
Observons que $E$ est un objet de $\mathcal{A}$. Considérons le
diagramme suivant
$$
\xymatrix{
E \ar[d]_{c_E} \ar[r] &
A \ar[d]_{c_A} \ar[r]_f &
A'[1] \ar[d]^{c_{A'}[1]}
\ar@{..>}@<-1ex>[d]_\gamma \ar@{..>}[ld]^\epsilon \ar[r] &
E[1] \ar[d]^{c_E[1]} \\
F(E) \ar[r] &
F(A) \ar[r]^{F(f)} &
F(A')[1] \ar[r] &
F(E)[1]
}
$$
dont les lignes sont des triangles distingués.
Le carré de droite est déjà commutatif, mais on ne sait pas encore que
le carré central l'est. D'après les axiomes d'une catégorie triangulée,
on peut trouver un morphisme $\gamma$ qui rende le diagramme commutatif.
Alors le composé de $\gamma - c_{A'}[1]$ avec
$F(A')[1] \to F(E)[1]$ est nul ; on peut donc
trouver $\epsilon : A'[1] \to F(A)$ tel que
$\gamma - c_{A'}[1] = F(f) \circ \epsilon$. Cependant, toute flèche
$A'[1] \to F(A)$ est nulle, car c'est une classe d'extension négative
entre objets de $\mathcal{A}$. Ainsi, $\gamma = c_{A'}[1]$,
et l'on conclut que le carré central est lui aussi commutatif, ce que l'on
voulait montrer.

\medskip\noindent
Pour achever la démonstration, nous allons raisonner par récurrence sur $w$
et construire des isomorphismes $c_X : X \to F(X)$ pour tout
$X$ tel que $w(X) \leq w$, compatibles avec tous les morphismes entre
de tels objets. Le cas initial $w = 1$ a été démontré ci-dessus. Supposons
le résultat connu pour un certain $w \geq 1$.

\medskip\noindent
Soit $X$ un objet tel que $w(X) = w + 1$. Choisissons $a \in \mathbf{Z}$ avec
$H^i(X) = 0$ pour $i \not \in [a, a + w]$. Posons $b = a + w$, de sorte que
$H^b(X)$ soit non nul. Choisissons $N[-b] \to X$ comme dans le lemme \ref{lemma-good-map}.
Choisissons un triangle distingué
$$
N[-b] \to X \to Y \to N[-b + 1]
$$
Le calcul de la suite exacte longue de cohomologie donne
$w(Y) \leq w$. Par récurrence, on obtient donc les flèches pleines
du diagramme suivant
$$
\xymatrix{
N[-b] \ar[r] \ar[d]_{c_N[-b]} &
X \ar[r] \ar@{..>}[d]_{c_{N[-b] \to X}} &
Y \ar[r] \ar[d]^{c_Y} &
N[-b + 1] \ar[d]^{c_N[-b + 1]} \\
F(N)[-b] \ar[r] &
F(X) \ar[r] &
F(Y) \ar[r] &
F(N)[-b + 1]
}
$$
On obtient la flèche pointillée $c_{N[-b] \to X}$.
D'après Catégories dérivées, lemme \ref{derived-lemma-uniqueness-third-arrow},
la flèche pointillée est unique, car $\Hom(X, F(N)[-b]) \cong \Hom(X, N[-b]) = 0$
par le choix de $N$. En fait, $c_{N[-b] \to X}$ est l'unique flèche
pointillée qui rende commutatif le carré de sommets $X, Y, F(X), F(Y)$.

\medskip\noindent
Soit $N'[-b] \to X$ un autre morphisme comme dans le lemme \ref{lemma-good-map},
et prouvons que $c_{N[-b] \to X} = c_{N'[-b] \to X}$.
Observons que le morphisme $(N \oplus N')[-b] \to X$ satisfait lui aussi les
conditions du lemme \ref{lemma-good-map}.
On peut donc supposer que $N'[-b] \to X$ se factorise
en $N'[-b] \to N[-b] \to X$ pour un certain morphisme $N' \to N$.
Choisissons des triangles distingués $N[-b] \to X \to Y \to N[-b + 1]$ et
$N'[-b] \to X \to Y' \to N'[-b + 1]$. D'après l'axiome TR3, il existe
un morphisme $g : Y' \to Y$ qui, avec $\text{id}_X$ et $N' \to N$,
forme un morphisme de triangles. Comme on dispose de
(\ref{equation-to-show}) pour $g$, on en déduit que
$$
(F(X) \to F(Y)) \circ c_{N'[-b] \to X} = (F(X) \to F(Y)) \circ c_{N[-b] \to X}
$$
L'unicité de $c_{N[-b] \to X}$ relevée dans la construction
ci-dessus montre alors que $c_{N'[-b] \to X} = c_{N[-b] \to X}$.

\medskip\noindent
On peut donc maintenant définir, pour $X$ de largeur $w + 1$, l'isomorphisme
$c_X : X \to F(X)$ comme la valeur commune des morphismes $c_{N[-b] \to X}$,
où $N[-b] \to X$ est comme dans le lemme \ref{lemma-good-map}. Pour achever
la démonstration, il faut montrer que les diagrammes (\ref{equation-to-show})
sont commutatifs pour tout morphisme $f : X \to X'$ entre objets tels que $w(X) \leq w + 1$
et $w(X') \leq w + 1$. Choisissons $a \leq b \leq a + w$ tel que
$H^i(X) = 0$ pour $i \not \in [a, b]$ et
$a' \leq b' \leq a' + w$ tel que $H^i(X') = 0$ pour
$i \not \in [a', b']$. Nous allons raisonner par récurrence sur
$(b' - a') + (b - a)$ pour démontrer l'assertion. (Le cas initial
est celui où ce nombre est nul, ce qui convient puisque $w \geq 1$.)
Distinguons deux cas.

\medskip\noindent
Cas I : $b' < b$. Dans ce cas, d'après le lemme \ref{lemma-good-map-zero},
on peut choisir $N[-b] \to X$ comme dans le lemme \ref{lemma-good-map},
de sorte que le composé $N[-b] \to X \to X'$ soit nul.
Choisissons un triangle distingué $N[-b] \to X \to Y \to N[-b + 1]$. Comme
$N[-b] \to X'$ est nul, on voit que $f$ se factorise
en $X \to Y \to X'$. Comme $H^i(Y)$ n'est non nul que pour $i \in [a, b - 1]$,
l'hypothèse de récurrence montre que (\ref{equation-to-show}) est commutatif pour
$Y \to X'$. Le diagramme (\ref{equation-to-show}) est commutatif pour
$X \to Y$ par construction si $w(X) = w + 1$, et par notre première
hypothèse de récurrence si $w(X) \leq w$.
Ainsi, (\ref{equation-to-show}) est commutatif pour $f$.

\medskip\noindent
Cas II : $b' \geq b$. Dans ce cas, choisissons $N'[-b'] \to X'$
comme dans le lemme \ref{lemma-good-map}.
On peut aussi supposer que $\Hom(H^{b'}(X), N') = 0$ (cela n'est
pertinent que si $b' = b$), car on peut, par exemple,
remplacer $N'$ par un objet $N''$ muni d'une surjection sur $N' \oplus H^{b'}(X)$
et tel que $\Hom(N' \oplus H^{b'}(X), N'') = 0$.
Choisissons un triangle distingué
$N'[-b'] \to X' \to Y' \to N'[-b' + 1]$. Comme
$\Hom(X, X') \to \Hom(X, Y')$ est injectif par notre choix de $N'$
(nous omettons les détails), il en va de même pour
$\Hom(X, F(X')) \to \Hom(X, F(Y'))$.
Il suffit donc, dans ce cas, de vérifier que
(\ref{equation-to-show}) est commutatif pour le composé $X \to Y'$
des morphismes $X \to X' \to Y'$.
Comme $H^i(Y')$ n'est non nul que pour $i \in [a', b' - 1]$,
on conclut par l'hypothèse de récurrence.
\end{proof}










\section{Établissement de la pleine fidélité}
\label{section-get-fully-faithful}

\noindent
Il nous sera utile de savoir quand un foncteur est pleinement fidèle ;
nous proposons la variante suivante de \cite[Lemma 2.15]{Orlov-K3}.

\begin{lemma}
\label{lemma-get-fully-faithful}
\begin{reference}
Variante de \cite[Lemma 2.15]{Orlov-K3}
\end{reference}
Soit $F : \mathcal{D} \to \mathcal{D}'$ un foncteur exact de
catégories triangulées. Soit $S \subset \Ob(\mathcal{D})$ un
ensemble d'objets. Supposons que
\begin{enumerate}
\item $F$ possède des adjoints à droite et à gauche,
\item pour $K \in \mathcal{D}$, si $\Hom(E, K[i]) = 0$ pour tous
$E \in S$ et $i \in \mathbf{Z}$, alors $K = 0$,
\item pour $K \in \mathcal{D}$, si $\Hom(K, E[i]) = 0$ pour tous
$E \in S$ et $i \in \mathbf{Z}$, alors $K = 0$,
\item l'application $\Hom(E, E'[i]) \to \Hom(F(E), F(E')[i])$ induite par $F$
est bijective pour tous $E, E' \in S$ et $i \in \mathbf{Z}$.
\end{enumerate}
Alors $F$ est pleinement fidèle.
\end{lemma}

\begin{proof}
Notons $F_r$ et $F_l$ les adjoints à droite et à gauche de $F$. Pour
$E \in S$, choisissons un triangle distingué
$$
E \to F_r(F(E)) \to C \to E[1]
$$
où la première flèche est l'unité de l'adjonction. Pour $E' \in S$, on a
$$
\Hom(E', F_r(F(E))[i]) = \Hom(F(E'), F(E)[i]) = \Hom(E', E[i])
$$
La dernière égalité vaut par l'hypothèse (4).
En appliquant donc le foncteur homologique $\Hom(E', -)$
(Catégories dérivées, lemme \ref{derived-lemma-representable-homological})
au triangle distingué ci-dessus, on en déduit que $\Hom(E', C[i]) = 0$
pour tout $i \in \mathbf{Z}$ et $E' \in S$. L'hypothèse (2) donne alors
$C = 0$ et $E = F_r(F(E))$.

\medskip\noindent
Pour $K \in \Ob(\mathcal{D})$, choisissons un triangle distingué
$$
F_l(F(K)) \to K \to C \to F_l(F(K))[1]
$$
où la première flèche est la counité de l'adjonction. Pour $E \in S$,
on a
$$
\Hom(F_l(F(K)), E[i]) = \Hom(F(K), F(E)[i]) =
\Hom(K, F_r(F(E))[i]) = \Hom(K, E[i])
$$
où la dernière égalité résulte du premier paragraphe.
On conclut donc comme précédemment que $\Hom(C, E[i]) = 0$ pour tout $E \in S$
et $i \in \mathbf{Z}$. Ainsi, $C = 0$ d'après l'hypothèse (3).
Le foncteur $F$ est donc pleinement fidèle d'après Catégories, lemme
\ref{categories-lemma-adjoint-fully-faithful}.
\end{proof}

\begin{lemma}
\label{lemma-duality-at-point}
Soit $k$ un corps. Soit $X$ un schéma de type fini sur $k$ qui
est régulier. Soit $x \in X$ un point fermé. Pour un
$\mathcal{O}_X$-module cohérent $\mathcal{F}$ à support en $x$, choisissons
un $\mathcal{O}_X$-module cohérent $\mathcal{F}'$ à support en $x$
tel que $\mathcal{F}_x$ et $\mathcal{F}'_x$ soient duaux de Matlis.
Alors il existe un isomorphisme
$$
\Hom_X(\mathcal{F}, M) =
H^0(X, M \otimes_{\mathcal{O}_X}^\mathbf{L} \mathcal{F}'[-d_x])
$$
où $d_x = \dim(\mathcal{O}_{X, x})$,
fonctoriel en $M$ dans $D_{perf}(\mathcal{O}_X)$.
\end{lemma}

\begin{proof}
Puisque $\mathcal{F}$ est à support en $x$, on a
$$
\Hom_X(\mathcal{F}, M) =
\Hom_{\mathcal{O}_{X, x}}(\mathcal{F}_x, M_x)
$$
et, de même, on a
$$
H^0(X, M \otimes_{\mathcal{O}_X}^\mathbf{L} \mathcal{F}'[-d_x]) =
\text{Tor}^{\mathcal{O}_{X, x}}_{d_x}(M_x, \mathcal{F}'_x)
$$
Il suffit donc de montrer que, pour un anneau local noethérien régulier $A$
de dimension $d$ et un $A$-module $N$ de longueur finie, si
$N'$ est le dual de Matlis de $N$, alors il existe un isomorphisme fonctoriel
$$
\Hom_A(N, K) = \text{Tor}^A_d(K, N')
$$
pour $K$ dans $D_{perf}(A)$. On peut écrire le membre de gauche sous la forme
$H^0(R\Hom_A(N, A) \otimes_A^\mathbf{L} K)$ d'après
Compléments sur l'algèbre, lemme \ref{more-algebra-lemma-dual-perfect-complex},
et le fait que $N$ détermine un objet parfait de $D(A)$.
La formule vaut donc parce que
$$
R\Hom_A(N, A) = R\Hom_A(N, A[d])[-d] = N'[-d]
$$
d'après Complexes dualisants, lemme \ref{dualizing-lemma-dualizing-finite-length},
et parce que $A[d]$ est un complexe dualisant normalisé sur $A$
($A$ est un anneau de Gorenstein d'après
Complexes dualisants, lemme \ref{dualizing-lemma-regular-gorenstein}).
\end{proof}

\begin{lemma}
\label{lemma-orthogonal-point-sheaf}
Soit $k$ un corps. Soit $X$ un schéma de type fini sur $k$ qui
est régulier. Soit $x \in X$ un point fermé et notons $\mathcal{O}_x$
le faisceau gratte-ciel en $x$ de valeur $\kappa(x)$. Soit $K$ dans
$D_{perf}(\mathcal{O}_X)$.
\begin{enumerate}
\item Si $\Ext^i_X(\mathcal{O}_x, K) = 0$, alors il existe un voisinage
ouvert $U$ de $x$ tel que $H^{i - d_x}(K)|_U = 0$, où
$d_x = \dim(\mathcal{O}_{X, x})$.
\item Si $\Hom_X(\mathcal{O}_x, K[i]) = 0$ pour tout
$i \in \mathbf{Z}$, alors $K$ est nul dans un voisinage ouvert de $x$.
\item Si $\Ext^i_X(K, \mathcal{O}_x) = 0$, alors il existe un voisinage
ouvert $U$ de $x$ tel que $H^i(K^\vee)|_U = 0$.
\item Si $\Hom_X(K, \mathcal{O}_x[i]) = 0$ pour tout
$i \in \mathbf{Z}$, alors $K$ est nul dans un voisinage ouvert de $x$.
\item Si $H^i(X, K \otimes_{\mathcal{O}_X}^\mathbf{L} \mathcal{O}_x) = 0$,
alors il existe un voisinage ouvert $U$ de $x$ tel que
$H^i(K)|_U = 0$.
\item Si $H^i(X, K \otimes_{\mathcal{O}_X}^\mathbf{L} \mathcal{O}_x) = 0$
pour $i \in \mathbf{Z}$, alors $K$ est nul dans un
voisinage ouvert de $x$.
\end{enumerate}
\end{lemma}

\begin{proof}
Observons que $H^i(X, K \otimes_{\mathcal{O}_X}^\mathbf{L} \mathcal{O}_x)$
est égal à $K_x  \otimes_{\mathcal{O}_{X, x}}^\mathbf{L} \kappa(x)$.
L'assertion (5) résulte donc de Compléments sur l'algèbre, lemme
\ref{more-algebra-lemma-cut-complex-in-two}.
L'assertion (6) résulte de l'assertion (5).
L'assertion (1) résulte de l'assertion (5), du lemme \ref{lemma-duality-at-point} et du
fait que le dual de Matlis de $\kappa(x)$ est $\kappa(x)$.
L'assertion (2) résulte de l'assertion (1).
L'assertion (3) résulte de l'assertion (5) et du fait que
$\Ext^i(K, \mathcal{O}_x) =
H^i(X, K^\vee \otimes_{\mathcal{O}_X}^\mathbf{L} \mathcal{O}_x)$ d'après
Cohomologie, lemme \ref{cohomology-lemma-dual-perfect-complex}.
L'assertion (4) résulte de l'assertion (3) et du fait que $K \cong (K^\vee)^\vee$
d'après le lemme qui vient d'être cité.
\end{proof}

\begin{lemma}
\label{lemma-hom-into-point-sheaf}
Soit $X$ un schéma noethérien. Soit $x \in X$ un point fermé et
notons $\mathcal{O}_x$ le faisceau gratte-ciel en $x$ de valeur $\kappa(x)$.
Soit $K$ dans $D^b_{\textit{Coh}}(\mathcal{O}_X)$. Soit $b \in \mathbf{Z}$.
Les conditions suivantes sont équivalentes :
\begin{enumerate}
\item $H^i(K)_x = 0$ pour tout $i > b$, et
\item $\Hom_X(K, \mathcal{O}_x[-i]) = 0$ pour tout $i > b$.
\end{enumerate}
\end{lemma}

\begin{proof}
Considérons le complexe $K_x$ dans $D^b_{\textit{Coh}}(\mathcal{O}_{X, x})$.
Il existe un entier $b_x \in \mathbf{Z}$ tel que $K_x$
puisse être représenté par un complexe borné supérieurement
$$
\ldots \to
\mathcal{O}_{X, x}^{\oplus n_{b_x - 2}} \to
\mathcal{O}_{X, x}^{\oplus n_{b_x - 1}} \to
\mathcal{O}_{X, x}^{\oplus n_{b_x}} \to 0 \to \ldots
$$
où $\mathcal{O}_{X, x}^{\oplus n_i}$ est placé en degré $i$
et où tous les morphismes de transition sont donnés par des matrices dont les
coefficients appartiennent à $\mathfrak m_x$. Voir
Compléments sur l'algèbre, lemme
\ref{more-algebra-lemma-lift-pseudo-coherent-from-residue-field}.
Le résultat s'en déduit facilement (et les conditions équivalentes
sont satisfaites si et seulement si $b \geq b_x$).
\end{proof}

\begin{lemma}
\label{lemma-get-fully-faithful-geometric}
Soit $k$ un corps. Soient $X$ et $Y$ des schémas propres sur $k$.
Supposons que $X$ soit régulier. Alors un foncteur exact $k$-linéaire
$F : D_{perf}(\mathcal{O}_X) \to D_{perf}(\mathcal{O}_Y)$
est pleinement fidèle si et seulement si
pour tous les points fermés $x, x' \in X$, les applications
$$
F : \Ext^i_X(\mathcal{O}_x, \mathcal{O}_{x'})
\longrightarrow
\Ext^i_Y(F(\mathcal{O}_x), F(\mathcal{O}_{x'}))
$$
sont des isomorphismes pour tout $i \in \mathbf{Z}$.
Ici, $\mathcal{O}_x$ est le faisceau gratte-ciel en $x$ de valeur $\kappa(x)$.
\end{lemma}

\begin{proof}
D'après le lemme \ref{lemma-always-right-adjoints}, le foncteur $F$
possède un adjoint à gauche et un adjoint à droite. On peut donc appliquer le critère
du lemme \ref{lemma-get-fully-faithful},
car les hypothèses (2) et (3) de ce lemme
résultent du lemme \ref{lemma-orthogonal-point-sheaf}.
\end{proof}

\begin{lemma}
\label{lemma-noah-pre}
\begin{reference}
Courriel de Noah Olander du 9 juin 2020
\end{reference}
Soit $k$ un corps. Soit $X$ un schéma propre sur $k$ qui est régulier.
Soit $F : D_{perf}(\mathcal{O}_X) \to D_{perf}(\mathcal{O}_X)$
un foncteur exact $k$-linéaire. Supposons que, pour tout
$\mathcal{O}_X$-module cohérent $\mathcal{F}$ tel que $\dim(\text{Supp}(\mathcal{F})) = 0$,
il existe un isomorphisme $\mathcal{F} \cong F(\mathcal{F})$.
Alors $F$ est pleinement fidèle.
\end{lemma}

\begin{proof}
D'après le lemme \ref{lemma-get-fully-faithful-geometric}, il suffit de montrer
que les applications
$$
F : \Ext^i_X(\mathcal{O}_x, \mathcal{O}_{x'})
\longrightarrow
\Ext^i_X(F(\mathcal{O}_x), F(\mathcal{O}_{x'}))
$$
sont des isomorphismes pour tout $i \in \mathbf{Z}$ et tous les points fermés
$x, x' \in X$. Par hypothèse, la source et le but sont isomorphes.
Si $x \not = x'$, alors les deux membres sont nuls et le résultat est vrai.
Si $x = x'$, il suffit alors de prouver que l'application est injective
ou surjective. Pour $i < 0$, les deux membres sont nuls et le résultat est vrai.
Pour $i = 0$, tout morphisme non nul $\alpha : \mathcal{O}_x \to \mathcal{O}_x$ de
$\mathcal{O}_X$-modules est un isomorphisme. Ainsi, $F(\alpha)$ est lui aussi un
isomorphisme, donc $F(\alpha)$ est non nul. Cela prouve le résultat pour $i = 0$.
Pour $i = 1$, un élément non nul $\xi$ de $\Ext^1(\mathcal{O}_x, \mathcal{O}_x)$
correspond à une suite exacte courte non scindée
$$
0 \to \mathcal{O}_x \to \mathcal{F} \to \mathcal{O}_x \to 0
$$
Comme $F(\mathcal{F}) \cong \mathcal{F}$, on voit que $F(\mathcal{F})$
est également une extension non scindée de $\mathcal{O}_x$ par $\mathcal{O}_x$.
Puisque $\mathcal{O}_x \cong F(\mathcal{O}_x)$ est un
$\mathcal{O}_X$-module simple et que $\mathcal{F} \cong F(\mathcal{F})$ est de
longueur $2$, on voit que, dans le triangle distingué
$$
F(\mathcal{O}_x) \to F(\mathcal{F}) \to F(\mathcal{O}_x)
\xrightarrow{F(\xi)} F(\mathcal{O}_x)[1]
$$
les deux premières flèches doivent former une suite exacte courte qui doit être
isomorphe à la suite exacte courte ci-dessus, et qui est donc non scindée.
Il s'ensuit que $F(\xi)$ est non nul, ce qui permet de conclure pour $i = 1$.
Pour $i > 1$, la composition des classes d'extensions définit une surjection
$$
\Ext^1(F(\mathcal{O}_x), F(\mathcal{O}_x)) \otimes \ldots \otimes
\Ext^1(F(\mathcal{O}_x), F(\mathcal{O}_x))
\longrightarrow
\Ext^i(F(\mathcal{O}_x), F(\mathcal{O}_x))
$$
Voir Dualité pour les schémas, lemme \ref{duality-lemma-regular-ideal-ext}.
La surjectivité en degré $1$ implique donc la surjectivité pour $i > 0$.
Cela achève la démonstration.
\end{proof}










\section{Foncteurs particuliers}
\label{section-special-functors}

\noindent
Dans cette section, nous démontrons quelques résultats sur des foncteurs d'un type particulier
que nous utiliserons plus loin dans ce chapitre.

\begin{definition}
\label{definition-siblings-geometric}
Soit $k$ un corps. Soient $X$ et $Y$ des schémas de type fini sur $k$.
Rappelons que
$D^b_{\textit{Coh}}(\mathcal{O}_X) = D^b(\textit{Coh}(\mathcal{O}_X))$
d'après Catégories dérivées des schémas, proposition \ref{perfect-proposition-DCoh}.
Nous disons que deux foncteurs exacts $k$-linéaires
$$
F, F' :
D^b_{\textit{Coh}}(\mathcal{O}_X) = D^b(\textit{Coh}(\mathcal{O}_X))
\longrightarrow
D^b_{\textit{Coh}}(\mathcal{O}_Y)
$$
sont {\it frères}, ou que $F'$ est un {\it frère} de $F$, si $F$ et $F'$
sont frères au sens de la définition \ref{definition-siblings},
la catégorie abélienne étant $\textit{Coh}(\mathcal{O}_X)$.
Si $X$ est régulier, alors
$D_{perf}(\mathcal{O}_X) = D^b_{\textit{Coh}}(\mathcal{O}_X)$ d'après
Catégories dérivées des schémas, lemme \ref{perfect-lemma-perfect-on-noetherian},
et nous employons la même terminologie pour les foncteurs exacts $k$-linéaires
$F, F' : D_{perf}(\mathcal{O}_X) \to D_{perf}(\mathcal{O}_Y)$.
\end{definition}

\begin{lemma}
\label{lemma-exact-functor-preserving-Coh}
Soit $k$ un corps. Soient $X$ et $Y$ des schémas de type fini sur $k$, avec
$X$ séparé. Soit
$F : D^b_{\textit{Coh}}(\mathcal{O}_X) \to D^b_{\textit{Coh}}(\mathcal{O}_Y)$
un foncteur exact $k$-linéaire qui envoie
$\textit{Coh}(\mathcal{O}_X) \subset D^b_{\textit{Coh}}(\mathcal{O}_X)$
dans
$\textit{Coh}(\mathcal{O}_Y) \subset D^b_{\textit{Coh}}(\mathcal{O}_Y)$.
Alors il existe un foncteur de Fourier-Mukai
$F' : D^b_{\textit{Coh}}(\mathcal{O}_X) \to D^b_{\textit{Coh}}(\mathcal{O}_Y)$
dont le noyau est un $\mathcal{O}_{X \times Y}$-module cohérent $\mathcal{K}$,
plat sur $X$ et à support fini sur $Y$ ; le foncteur ainsi obtenu est un frère de $F$.
\end{lemma}

\begin{proof}
Notons $H : \textit{Coh}(\mathcal{O}_X) \to \textit{Coh}(\mathcal{O}_Y)$
la restriction de $F$. Comme $F$ est un foncteur exact de catégories
triangulées, $H$ est un foncteur exact de catégories abéliennes.
Bien entendu, $H$ est $k$-linéaire puisque $F$ l'est. D'après
Foncteurs et morphismes, lemme \ref{functors-lemma-functor-coherent-over-field},
on obtient un $\mathcal{O}_{X \times Y}$-module cohérent
$\mathcal{K}$ qui est plat sur $X$ et à support fini sur $Y$.
Soit $F'$ le foncteur de Fourier-Mukai défini à l'aide de $\mathcal{K}$,
de sorte que $F'$ se restreint à $H$ sur $ \textit{Coh}(\mathcal{O}_X)$.
Le foncteur $F'$ envoie $D^b_{\textit{Coh}}(\mathcal{O}_X)$
dans $D^b_{\textit{Coh}}(\mathcal{O}_Y)$ d'après le
lemme \ref{lemma-fourier-mukai-Coh}.
Observons que $F$ et $F'$ satisfont les première et deuxième
conditions du lemme \ref{lemma-sibling-fully-faithful} et sont donc frères.
\end{proof}

\begin{remark}
\label{remark-difficult}
Si $F, F' : D^b_{\textit{Coh}}(\mathcal{O}_X) \to \mathcal{D}$ sont frères, si $F$
est pleinement fidèle et si $X$ est réduit et projectif sur $k$, alors
$F \cong F'$ ; cela résulte de la
proposition \ref{proposition-siblings-isomorphic} par l'argument
donné dans la démonstration du théorème \ref{theorem-fully-faithful}.
Cependant, en général, nous ne savons pas si des frères sont isomorphes.
Même dans la situation du lemme \ref{lemma-exact-functor-preserving-Coh},
il semble difficile de prouver que les frères $F$ et $F'$
sont des foncteurs isomorphes. Si $X$ est lisse et propre sur $k$
et si $F$ est pleinement fidèle, alors $F \cong F'$, comme le montre
\cite{Noah}.
Si vous disposez d'une démonstration ou d'un contre-exemple dans des situations plus générales,
veuillez écrire à
\href{mailto:stacks.project@gmail.com}{stacks.project@gmail.com}.
\end{remark}

\begin{lemma}
\label{lemma-two-functors-pre}
Soit $k$ un corps. Soient $X$ et $Y$ des schémas propres sur $k$. Supposons
que $X$ soit régulier. Soient
$F, G : D_{perf}(\mathcal{O}_X) \to D_{perf}(\mathcal{O}_Y)$
des foncteurs exacts $k$-linéaires tels que
\begin{enumerate}
\item $F(\mathcal{F}) \cong G(\mathcal{F})$ pour tout
$\mathcal{O}_X$-module cohérent $\mathcal{F}$ tel que $\dim(\text{Supp}(\mathcal{F})) = 0$,
\item $F$ soit pleinement fidèle.
\end{enumerate}
Alors l'image essentielle de $G$ est contenue dans l'image essentielle
de $F$.
\end{lemma}

\begin{proof}
Rappelons que $F$ et $G$ possèdent chacun un adjoint à gauche et un adjoint à droite, voir le
lemme \ref{lemma-always-right-adjoints}. En particulier,
l'image essentielle $\mathcal{A} \subset D_{perf}(\mathcal{O}_Y)$ de $F$
satisfait les conditions équivalentes de
Catégories dérivées, lemme \ref{derived-lemma-right-adjoint}.
Nous affirmons que $G$ se factorise par $\mathcal{A}$.
Comme $\mathcal{A} = {}^\perp(\mathcal{A}^\perp)$ d'après
Catégories dérivées, lemme \ref{derived-lemma-right-adjoint},
il suffit de montrer que $\Hom_Y(G(M), N) = 0$ pour
tout $M$ dans $D_{perf}(\mathcal{O}_X)$ et $N \in \mathcal{A}^\perp$.
On a
$$
\Hom_Y(G(M), N) = \Hom_X(M, G_r(N))
$$
où $G_r$ est l'adjoint à droite de $G$. Il suffit donc de prouver
que $G_r(N) = 0$. Comme
$G(\mathcal{F}) \cong F(\mathcal{F})$ pour $\mathcal{F}$ comme en (1),
on voit que
$$
\Hom_X(\mathcal{F}, G_r(N)) =
\Hom_Y(G(\mathcal{F}), N) =
\Hom_Y(F(\mathcal{F}), N) = 0
$$
puisque $N$ appartient à l'orthogonal à droite de l'image essentielle $\mathcal{A}$ de $F$.
Bien entendu, le groupe $\Hom_X(\mathcal{F}, G_r(N)[i])$ s'annule encore
pour tout $i \in \mathbf{Z}$. Ainsi, $G_r(N) = 0$ d'après le
lemme \ref{lemma-orthogonal-point-sheaf}, ce qui conclut.
\end{proof}

\begin{lemma}
\label{lemma-noah}
\begin{reference}
Courriel de Noah Olander du 8 juin 2020
\end{reference}
Soit $k$ un corps. Soit $X$ un schéma propre sur $k$ qui est régulier.
Soit $F : D_{perf}(\mathcal{O}_X) \to D_{perf}(\mathcal{O}_X)$
un foncteur exact $k$-linéaire. Supposons que, pour tout
$\mathcal{O}_X$-module cohérent $\mathcal{F}$ tel que $\dim(\text{Supp}(\mathcal{F})) = 0$,
il existe un isomorphisme $\mathcal{F} \cong F(\mathcal{F})$.
Alors il existe un automorphisme $f : X \to X$ sur $k$
qui induit l'identité
sur l'espace topologique sous-jacent\footnote{Cela force souvent $f$
à être l'identité, voir Variétés, lemme \ref{varieties-lemma-automorphism}.},
et un $\mathcal{O}_X$-module inversible $\mathcal{L}$
tels que $F$ et $F'(M) = f^*M \otimes_{\mathcal{O}_X}^\mathbf{L} \mathcal{L}$
soient frères.
\end{lemma}

\begin{proof}
D'après le lemme \ref{lemma-noah-pre}, le foncteur $F$ est pleinement fidèle.
D'après le lemme \ref{lemma-two-functors-pre}, l'image essentielle du
foncteur identité est contenue dans l'image essentielle de $F$, c'est-à-dire
que $F$ est essentiellement surjectif. Ainsi, $F$ est une équivalence.
Observons que le quasi-inverse $F^{-1}$ satisfait les mêmes hypothèses
que $F$.

\medskip\noindent
Soit $M \in D_{perf}(\mathcal{O}_X)$, et supposons $H^i(M) = 0$ pour $i > b$.
Comme $F$ est pleinement fidèle, on voit que
$$
\Hom_X(M, \mathcal{O}_x[-i]) =
\Hom_X(F(M), F(\mathcal{O}_x)[-i]) \cong
\Hom_X(F(M), \mathcal{O}_x[-i])
$$
pour tout $i \in \mathbf{Z}$ et tout point fermé $x$ de $X$.
Le lemme \ref{lemma-hom-into-point-sheaf} montre donc que les faisceaux de cohomologie de $F(M)$
sont nuls aux degrés $> b$.

\medskip\noindent
Soit $\mathcal{F}$ un $\mathcal{O}_X$-module cohérent. D'après
ce qui précède, les faisceaux de cohomologie non nuls de $F(\mathcal{F})$
ne se trouvent qu'aux degrés $\leq 0$.
Posons $\mathcal{G} = H^0(F(\mathcal{F}))$. Choisissons un triangle
distingué
$$
K \to F(\mathcal{F}) \to \mathcal{G} \to K[1]
$$
Alors les faisceaux de cohomologie non nuls de $K$ ne se trouvent qu'aux
degrés $\leq -1$.
En appliquant $F^{-1}$, on obtient un triangle distingué
$$
F^{-1}(K) \to \mathcal{F} \to F^{-1}(\mathcal{G}) \to F^{-1}(K')[1]
$$
Comme les faisceaux de cohomologie non nuls de $F^{-1}(K)$ ne se trouvent qu'aux
degrés $\leq -1$ (par le paragraphe précédent appliqué à $F^{-1}$),
on voit que la flèche $F^{-1}(K) \to \mathcal{F}$ est nulle
(Catégories dérivées, lemme \ref{derived-lemma-negative-exts}).
Ainsi, $K \to F(\mathcal{F})$ est nul, ce qui implique
que $F(\mathcal{F}) = \mathcal{G}$ par le choix du
premier triangle distingué.

\medskip\noindent
Le paragraphe précédent montre que $F$ préserve
$\textit{Coh}(\mathcal{O}_X)$ et définit même une équivalence
$H : \textit{Coh}(\mathcal{O}_X) \to \textit{Coh}(\mathcal{O}_X)$.
D'après Foncteurs et morphismes, lemme
\ref{functors-lemma-equivalence-coherent-over-field},
on obtient un automorphisme $f : X \to X$ sur $k$
et un $\mathcal{O}_X$-module inversible $\mathcal{L}$
tels que $H(\mathcal{F}) = f^*\mathcal{F} \otimes \mathcal{L}$.
Posons $F'(M) = f^*M \otimes_{\mathcal{O}_X}^\mathbf{L} \mathcal{L}$.
Le lemme \ref{lemma-sibling-fully-faithful}
montre que $F$ et $F'$ sont frères.
Pour voir que $f$ induit l'identité sur l'espace topologique
sous-jacent de $X$, on utilise $F(\mathcal{O}_x) \cong \mathcal{O}_x$
et le fait que le support de $\mathcal{O}_x$ est $\{x\}$.
Cela achève la démonstration.
\end{proof}

\begin{lemma}
\label{lemma-two-functors}
Soit $k$ un corps. Soient $X$ et $Y$ des schémas propres sur $k$.
Supposons que $X$ soit régulier.
Soient $F, G : D_{perf}(\mathcal{O}_X) \to D_{perf}(\mathcal{O}_Y)$
des foncteurs exacts $k$-linéaires tels que
\begin{enumerate}
\item $F(\mathcal{F}) \cong G(\mathcal{F})$ pour tout
$\mathcal{O}_X$-module cohérent $\mathcal{F}$ tel que $\dim(\text{Supp}(\mathcal{F})) = 0$,
\item $F$ soit pleinement fidèle, et
\item $G$ soit un foncteur de Fourier-Mukai dont le noyau appartient à
$D_{perf}(\mathcal{O}_{X \times Y})$.
\end{enumerate}
Alors il existe un foncteur de Fourier-Mukai
$F' : D_{perf}(\mathcal{O}_X) \to D_{perf}(\mathcal{O}_Y)$
dont le noyau appartient à $D_{perf}(\mathcal{O}_{X \times Y})$
et pour lequel $F$ et $F'$ sont frères.
\end{lemma}

\begin{proof}
L'image essentielle de $G$ est contenue dans l'image essentielle
de $F$ d'après le lemme \ref{lemma-two-functors-pre}.
Considérons le foncteur $H = F^{-1} \circ G$,
qui a un sens puisque $F$ est pleinement fidèle.
Le lemme \ref{lemma-noah} fournit un automorphisme $f : X \to X$
et un $\mathcal{O}_X$-module inversible $\mathcal{L}$ tels que
le foncteur $H' : K \mapsto f^*K \otimes \mathcal{L}$
soit un frère de $H$. En particulier,
$H$ est une auto-équivalence d'après le lemme \ref{lemma-sibling-faithful},
et $H$ induit une auto-équivalence de
$\textit{Coh}(\mathcal{O}_X)$ (puisque cela vaut pour son frère $H'$).
Ainsi, les quasi-inverses $H^{-1}$ et $(H')^{-1}$ existent et sont frères
(nous omettons un petit détail), et $(H')^{-1}$ envoie $M$ sur
$(f^{-1})^*(M \otimes_{\mathcal{O}_X}^\mathbf{L} \mathcal{L}^{\otimes -1})$,
qui est un foncteur de Fourier-Mukai (nous omettons les détails).
Alors, bien entendu, $F = G \circ H^{-1}$ est un frère de
$G \circ (H')^{-1}$. Comme les composés de foncteurs de Fourier-Mukai
sont de Fourier-Mukai d'après le
lemme \ref{lemma-compose-fourier-mukai},
on conclut.
\end{proof}





\section{Foncteurs pleinement fidèles}
\label{section-fully-faithful}

\noindent
Notre objectif est de prouver que les foncteurs pleinement fidèles entre catégories dérivées
sont frères de foncteurs de Fourier-Mukai, en suivant
\cite{Orlov-K3} et \cite{Ballard}.

\begin{situation}
\label{situation-fully-faithful}
Ici, $k$ est un corps. Nous avons des schémas propres et lisses $X$ et $Y$ sur $k$.
Nous avons un foncteur exact, $k$-linéaire et pleinement fidèle
$F : D_{perf}(\mathcal{O}_X) \to D_{perf}(\mathcal{O}_Y)$.
\end{situation}

\noindent
Avant de poursuivre, il est utile de lire au moins une partie de
Catégories dérivées, section \ref{derived-section-postnikov}.

\medskip\noindent
Rappelons que $X$ est régulier et possède donc la propriété de résolution
(Variétés, lemme \ref{varieties-lemma-smooth-regular} et
Catégories dérivées de schémas, lemme
\ref{perfect-lemma-regular-resolution-property}). Ainsi,
sur $X \times X$, nous pouvons choisir une résolution
$$
\ldots \to
\mathcal{E}_2 \boxtimes \mathcal{G}_2 \to
\mathcal{E}_1 \boxtimes \mathcal{G}_1 \to
\mathcal{E}_0 \boxtimes \mathcal{G}_0 \to
\mathcal{O}_\Delta \to 0
$$
où chaque $\mathcal{E}_i$ et chaque $\mathcal{G}_i$ est un
$\mathcal{O}_X$-module localement libre de type fini, voir le
lemme \ref{lemma-diagonal-resolution}.
En utilisant le complexe
\begin{equation}
\label{equation-original-complex}
\ldots \to
\mathcal{E}_2 \boxtimes \mathcal{G}_2 \to
\mathcal{E}_1 \boxtimes \mathcal{G}_1 \to
\mathcal{E}_0 \boxtimes \mathcal{G}_0
\end{equation}
dans $D_{perf}(\mathcal{O}_{X \times X})$, comme dans
Catégories dérivées, exemple \ref{derived-example-key-postnikov},
si, pour tout $n$, nous posons
$$
M_n = (\mathcal{E}_n \boxtimes \mathcal{G}_n \to \ldots \to
\mathcal{E}_0 \boxtimes \mathcal{G}_0)[-n]
$$
nous obtenons un système de Postnikov infini pour le complexe
(\ref{equation-original-complex}). Cela signifie que
les morphismes $M_0 \to M_1[1] \to M_2[2] \to \ldots$, ainsi que
$M_n \to \mathcal{E}_n \boxtimes \mathcal{G}_n$ et
$\mathcal{E}_n \boxtimes \mathcal{G}_n \to M_{n - 1}$,
satisfont à certaines conditions consignées dans
Catégories dérivées, définition \ref{derived-definition-postnikov-system}.
Posons
$$
\mathcal{F}_n = \Ker(\mathcal{E}_n \boxtimes \mathcal{G}_n \to
\mathcal{E}_{n - 1} \boxtimes \mathcal{G}_{n - 1})
$$
Observons que, puisque $\mathcal{O}_\Delta$ est plat sur $X$ via $\text{pr}_1$,
il en va de même de $\mathcal{F}_n$ pour tout $n$ (c'est une observation
commode, mais non essentielle). Nous avons
$$
H^q(M_n[n]) = \left\{
\begin{matrix}
\mathcal{O}_\Delta & \text{si} & q = 0 \\
\mathcal{F}_n & \text{si} & q = -n \\
0 & \text{si} & q \not = 0, -n
\end{matrix}
\right.
$$
Ainsi, pour $n \geq \dim(X \times X)$, nous avons
$$
M_n[n] \cong \mathcal{O}_\Delta \oplus \mathcal{F}_n[n]
$$
dans $D_{perf}(\mathcal{O}_{X \times X})$, d'après le

lemme \ref{lemma-split-complex-regular}.
\medskip\noindent
Nous nous intéressons au complexe
\begin{equation}
\label{equation-complex}
\ldots \to
\mathcal{E}_2 \boxtimes F(\mathcal{G}_2) \to
\mathcal{E}_1 \boxtimes F(\mathcal{G}_1) \to
\mathcal{E}_0 \boxtimes F(\mathcal{G}_0)
\end{equation}
dans $D_{perf}(\mathcal{O}_{X \times Y})$,
car la « totalisation » de ce complexe devrait
nous donner le noyau du foncteur de Fourier-Mukai que nous cherchons à construire.
Pour tous $i, j \geq 0$, nous avons
\begin{align*}
\Ext^q_{X \times Y}(\mathcal{E}_i \boxtimes F(\mathcal{G}_i),
\mathcal{E}_j \boxtimes F(\mathcal{G}_j))
& =
\bigoplus\nolimits_p
\Ext^{q + p}_X(\mathcal{E}_i, \mathcal{E}_j) \otimes_k
\Ext^{-p}_Y(F(\mathcal{G}_i), F(\mathcal{G}_j)) \\
& =
\bigoplus\nolimits_p
\Ext^{q + p}_X(\mathcal{E}_i, \mathcal{E}_j) \otimes_k
\Ext^{-p}_X(\mathcal{G}_i, \mathcal{G}_j)
\end{align*}
La seconde égalité vaut parce que $F$ est pleinement fidèle, et la première
résulte de Catégories dérivées de schémas, lemme
\ref{perfect-lemma-kunneth-Ext}.
On obtient que ces $\Ext^q$ sont nuls pour $q < 0$.
Ainsi, d'après Catégories dérivées, lemme
\ref{derived-lemma-existence-postnikov-system},
nous pouvons construire un système de Postnikov infini $K_0, K_1, K_2, \ldots$
dans $D_{perf}(\mathcal{O}_{X \times Y})$ pour le complexe
(\ref{equation-complex}).
Parallèlement à ce qui se passe pour $M_0, M_1, M_2, \ldots$, cela signifie
que nous obtenons des morphismes
$K_0 \to K_1[1] \to K_2[2] \to \ldots$, ainsi que
$K_n \to \mathcal{E}_n \boxtimes F(\mathcal{G}_n)$ et
$\mathcal{E}_n \boxtimes F(\mathcal{G}_n) \to K_{n - 1}$
dans $D_{perf}(\mathcal{O}_{X \times Y})$,
qui satisfont à certaines conditions consignées dans

Catégories dérivées, définition \ref{derived-definition-postnikov-system}.
\medskip\noindent
Soit $\mathcal{F}$ un $\mathcal{O}_X$-module cohérent dont le support
est un ensemble fini de points, c'est-à-dire tel que $\dim(\text{Supp}(\mathcal{F})) = 0$.
Considérons le foncteur exact entre catégories triangulées
$$
D_{perf}(\mathcal{O}_{X \times Y})
\longrightarrow
D_{perf}(\mathcal{O}_Y),\quad
N \longmapsto R\text{pr}_{2, *}(\text{pr}_1^*\mathcal{F}
\otimes^\mathbf{L}_{\mathcal{O}_{X \times Y}} N)
$$
Il s'ensuit que les objets $R\text{pr}_{2, *}(\text{pr}_1^*\mathcal{F}
\otimes^\mathbf{L}_{\mathcal{O}_{X \times Y}} K_i)$
forment un système de Postnikov pour le complexe de
$D_{perf}(\mathcal{O}_Y)$ dont les termes sont
$$
R\text{pr}_{2, *}(
(\mathcal{F} \otimes \mathcal{E}_i) \boxtimes F(\mathcal{G}_i)) =
\Gamma(X, \mathcal{F} \otimes \mathcal{E}_i) \otimes_k F(\mathcal{G}_i) =
F(\Gamma(X, \mathcal{F} \otimes \mathcal{E}_i) \otimes_k \mathcal{G}_i)
$$
Nous avons utilisé ici que $\mathcal{F} \otimes \mathcal{E}_i$ a une
cohomologie supérieure nulle, puisque son support est de dimension $0$.
D'autre part, en appliquant le foncteur exact
$$
D_{perf}(\mathcal{O}_{X \times X})
\longrightarrow
D_{perf}(\mathcal{O}_Y),\quad
N \longmapsto F(R\text{pr}_{2, *}(\text{pr}_1^*\mathcal{F}
\otimes^\mathbf{L}_{\mathcal{O}_{X \times X}} N))
$$
nous constatons que les objets
$F(R\text{pr}_{2, *}(\text{pr}_1^*\mathcal{F}
\otimes^\mathbf{L}_{\mathcal{O}_{X \times X}} M_n))$
forment un second système de Postnikov infini
pour le complexe de $D_{perf}(\mathcal{O}_Y)$ dont les termes sont
$$
F(R\text{pr}_{2, *}(
(\mathcal{F} \otimes \mathcal{E}_i) \boxtimes \mathcal{G}_i)) =
F(\Gamma(X, \mathcal{F} \otimes \mathcal{E}_i) \otimes_k \mathcal{G}_i)
$$
C'est le même complexe que précédemment ! Par unicité des systèmes de Postnikov
(Catégories dérivées, lemme \ref{derived-lemma-existence-postnikov-system}),
laquelle s'applique puisque
$$
\Ext^q_Y(
F(\Gamma(X, \mathcal{F} \otimes \mathcal{E}_i) \otimes_k \mathcal{G}_i),
F(\Gamma(X, \mathcal{F} \otimes \mathcal{E}_j) \otimes_k \mathcal{G}_j)) = 0,
\quad q < 0
$$
car $F$ est pleinement fidèle, nous obtenons un système d'isomorphismes
$$
F(R\text{pr}_{2, *}(\text{pr}_1^*\mathcal{F}
\otimes^\mathbf{L}_{\mathcal{O}_{X \times X}} M_n[n]))
\cong
R\text{pr}_{2, *}(\text{pr}_1^*\mathcal{F}
\otimes^\mathbf{L}_{\mathcal{O}_{X \times Y}} K_n[n])
$$
dans $D_{perf}(\mathcal{O}_Y)$, compatible avec les morphismes de
$D_{perf}(\mathcal{O}_Y)$ induits par les morphismes
$$
M_{n - 1}[n - 1] \to M_n[n]
\quad\text{et}\quad
K_{n - 1}[n - 1] \to K_n[n]
$$
$$
M_n \to \mathcal{E}_n \boxtimes \mathcal{G}_n
\quad\text{et}\quad
K_n \to \mathcal{E}_n \boxtimes F(\mathcal{G}_n)
$$
$$
\mathcal{E}_n \boxtimes \mathcal{G}_n \to M_{n - 1}
\quad\text{et}\quad
\mathcal{E}_n \boxtimes F(\mathcal{G}_n) \to K_{n - 1}
$$
qui font partie de la structure des systèmes de Postnikov.
Pour $n$ suffisamment grand, nous obtenons une décomposition en somme directe
$$
F(R\text{pr}_{2, *}(\text{pr}_1^*\mathcal{F}
\otimes^\mathbf{L}_{\mathcal{O}_{X \times X}} M_n[n]))
=
F(\mathcal{F}) \oplus
F(R\text{pr}_{2, *}(
\text{pr}_1^*\mathcal{F} \otimes_{\mathcal{O}_{X \times Y}} \mathcal{F}_n
))[n]
$$
qui correspond à la décomposition en somme directe de $M_n$ construite ci-dessus
(nous utilisons la platitude de $\mathcal{F}_n$ sur $X$ via $\text{pr}_1$
pour écrire un produit tensoriel ordinaire dans la formule ci-dessus, mais cela
n'est pas essentiel à l'argument).
D'après le lemme \ref{lemma-boundedness}, il existe un entier $m \geq 0$
tel que le premier facteur de cette décomposition en somme directe n'ait de
faisceaux de cohomologie non nuls que dans l'intervalle $[-m, m]$, tandis que le
second facteur n'a de faisceaux de cohomologie non nuls que dans l'intervalle
$[-m - n, m + \dim(X) - n]$.
On en conclut que le système $K_0 \to K_1[1] \to K_2[2] \to \ldots$
de $D_{perf}(\mathcal{O}_{X \times Y})$ satisfait aux hypothèses du
lemme \ref{lemma-bounded-fibres}, quitte à remplacer $m$ par un entier plus grand. On peut donc écrire
$$
K_n[n] = K \oplus C_n
$$
pour $n \gg 0$, de manière compatible avec les morphismes de transition, et
$C_n$ n'a de faisceaux de cohomologie non nuls que dans l'intervalle $[-m - n, m - n]$.
Notons $G$ le foncteur de Fourier-Mukai associé à $K$.
En regroupant ce qui précède, nous obtenons
$$
\begin{matrix}
G(\mathcal{F}) \oplus
R\text{pr}_{2, *}(
\text{pr}_1^*\mathcal{F} \otimes_{\mathcal{O}_{X \times Y}}^\mathbf{L} C_n)
\cong \\
R\text{pr}_{2, *}(\text{pr}_1^*\mathcal{F}
\otimes^\mathbf{L}_{\mathcal{O}_{X \times Y}} K_n[n]) \cong \\
F(R\text{pr}_{2, *}(\text{pr}_1^*\mathcal{F}
\otimes^\mathbf{L}_{\mathcal{O}_{X \times X}} M_n[n]))
\cong \\
F(\mathcal{F}) \oplus
F(R\text{pr}_{2, *}(
\text{pr}_1^*\mathcal{F} \otimes_{\mathcal{O}_{X \times Y}} \mathcal{F}_n
))[n]
\end{matrix}
$$
En examinant les degrés où vivent ces objets, on conclut que, pour $n \gg m$,
nous obtenons un isomorphisme
$$
F(\mathcal{F}) \cong G(\mathcal{F})
$$
De plus, cela vaut pour tout faisceau cohérent $\mathcal{F}$ sur $X$
dont le support est de dimension $0$.

\begin{lemma}
\label{lemma-fully-faithful}
Soit $k$ un corps. Soient $X$ et $Y$ des schémas propres et lisses sur $k$.
Étant donné un foncteur exact, $k$-linéaire et pleinement fidèle
$F : D_{perf}(\mathcal{O}_X) \to D_{perf}(\mathcal{O}_Y)$,
il existe un foncteur de Fourier-Mukai
$F' : D_{perf}(\mathcal{O}_X) \to D_{perf}(\mathcal{O}_Y)$ dont le noyau
appartient à $D_{perf}(\mathcal{O}_{X \times Y})$ et qui est frère de $F$.
\end{lemma}

\begin{proof}
Appliquons le lemme \ref{lemma-two-functors} à $F$ et au foncteur
$G$ construit ci-dessus.
\end{proof}

\noindent
Le théorème suivant reste vrai sans supposer que $X$ soit projectif,
voir \cite{Noah}.

\begin{theorem}[Orlov]
\label{theorem-fully-faithful}
\begin{reference}
\cite[Theorem 2.2]{Orlov-K3} ; cela est démontré dans \cite{Noah}
sans l'hypothèse que $X$ soit projectif
\end{reference}
Soit $k$ un corps. Soient $X$ et $Y$ des schémas propres et lisses sur $k$,
avec $X$ projectif sur $k$. Tout foncteur exact, $k$-linéaire et pleinement fidèle
$F : D_{perf}(\mathcal{O}_X) \to D_{perf}(\mathcal{O}_Y)$
est un foncteur de Fourier-Mukai associé à un noyau de
$D_{perf}(\mathcal{O}_{X \times Y})$.
\end{theorem}

\begin{proof}
Soit $F'$ le foncteur de Fourier-Mukai qui est frère de $F$
comme dans le lemme \ref{lemma-fully-faithful}.
D'après la proposition \ref{proposition-siblings-isomorphic}, nous avons $F \cong F'$
pourvu que nous montrions que $\textit{Coh}(\mathcal{O}_X)$ possède assez
d'objets négatifs. Toutefois, si $X = \Spec(k)$, par exemple,
ce n'est pas le cas. Nous décomposons donc d'abord $X = \coprod X_i$
en ses composantes connexes (et irréductibles), puis nous
montrons qu'il suffit de prouver le résultat pour chacun des
foncteurs composés (pleinement fidèles)
$$
F_i :
D_{perf}(\mathcal{O}_{X_i}) \to
D_{perf}(\mathcal{O}_X) \to
D_{perf}(\mathcal{O}_Y)
$$
Nous omettons les détails. Nous pouvons donc supposer $X$ irréductible.

\medskip\noindent
Le cas $\dim(X) = 0$. Ici, $X$ est le spectre d'une extension finie (séparable)
$k'/k$, et $D_{perf}(\mathcal{O}_X)$ est donc équivalente à la catégorie
des espaces vectoriels gradués sur $k'$ dans laquelle
$\mathcal{O}_X$ correspond à
l'espace vectoriel trivial de dimension $1$, placé en degré $0$.
Il est immédiat que deux foncteurs frères
$F, F' : D_{perf}(\mathcal{O}_X) \to D_{perf}(\mathcal{O}_Y)$
sont isomorphes. En effet, nous disposons d'un isomorphisme
$F(\mathcal{O}_X) \cong F'(\mathcal{O}_X)$
compatible avec l'action de la $k$-algèbre
$k' = \text{End}_{D_{perf}(\mathcal{O}_X)}(\mathcal{O}_X)$,
qui se prolonge canoniquement en un isomorphisme sur tout espace vectoriel gradué sur $k'$.

\medskip\noindent
Le cas $\dim(X) > 0$. Ici, $X$ est une variété projective lisse
de dimension $> 1$. Soit $\mathcal{F}$ un
$\mathcal{O}_X$-module cohérent.
Nous devons montrer qu'il existe un module cohérent $\mathcal{N}$ tel que
\begin{enumerate}
\item il existe une surjection $\mathcal{N} \to \mathcal{F}$, et
\item $\Hom(\mathcal{F}, \mathcal{N}) = 0$.
\end{enumerate}
Choisissons un $\mathcal{O}_X$-module inversible ample $\mathcal{L}$.
Nous affirmons que $\mathcal{N} = (\mathcal{L}^{\otimes n})^{\oplus r}$
convient pour $n \ll 0$ et $r$ assez grand.
La condition (1) résulte de
Propriétés, proposition \ref{properties-proposition-characterize-ample}.
Enfin, nous avons
$$
\Hom(\mathcal{F}, \mathcal{L}^{\otimes n}) =
H^0(X, \SheafHom(\mathcal{F}, \mathcal{L}^{\otimes n})) =
H^0(X, \SheafHom(\mathcal{F}, \mathcal{O}_X) \otimes \mathcal{L}^{\otimes n})
$$
Comme le dual $\SheafHom(\mathcal{F}, \mathcal{O}_X)$ est sans torsion, ce groupe
est nul pour $n \ll 0$ d'après Variétés, lemme
\ref{varieties-lemma-vanishin-h0-negative}. Cela achève la démonstration.
\end{proof}

\begin{proposition}
\label{proposition-equivalence}
Soit $k$ un corps. Soient $X$ et $Y$ des schémas propres et lisses sur $k$.
Si $F : D_{perf}(\mathcal{O}_X) \to D_{perf}(\mathcal{O}_Y)$
est une équivalence exacte $k$-linéaire de catégories triangulées, alors
il existe un foncteur de Fourier-Mukai
$F' : D_{perf}(\mathcal{O}_X) \to D_{perf}(\mathcal{O}_Y)$ dont le
noyau appartient à $D_{perf}(\mathcal{O}_{X \times Y})$,
qui est une équivalence et un frère de $F$.
\end{proposition}

\begin{proof}
Le foncteur $F'$ du lemme \ref{lemma-fully-faithful}
est une équivalence d'après le lemme \ref{lemma-sibling-faithful}.
\end{proof}

\begin{lemma}
\label{lemma-uniqueness}
Soit $k$ un corps. Soit $X$ un schéma propre et lisse sur $k$.
Soit $K \in D_{perf}(\mathcal{O}_{X \times X})$. Si le foncteur de Fourier-Mukai
$\Phi_K : D_{perf}(\mathcal{O}_X) \to D_{perf}(\mathcal{O}_X)$
est isomorphe au foncteur identité, alors
$K \cong \Delta_*\mathcal{O}_X$ dans $_{perf}(\mathcal{O}_{X \times X})$.
\end{lemma}

\begin{proof}
Soit $i$ le plus petit entier tel que le faisceau de cohomologie $H^i(K)$ soit
non nul. Soient $\mathcal{E}$ et $\mathcal{G}$ des $\mathcal{O}_X$-modules
localement libres de type fini. Alors
\begin{align*}
H^i(X \times X, K \otimes_{\mathcal{O}_{X \times X}}^\mathbf{L}
(\mathcal{E} \boxtimes \mathcal{G}))
& =
H^i(X, R\text{pr}_{2, *}(K \otimes_{\mathcal{O}_{X \times X}}^\mathbf{L}
(\mathcal{E} \boxtimes \mathcal{G}))) \\
& =
H^i(X, \Phi_K(\mathcal{E}) \otimes_{\mathcal{O}_X}^\mathbf{L} \mathcal{G}) \\
& \cong
H^i(X, \mathcal{E} \otimes \mathcal{G})
\end{align*}
ce qui est nul si $i < 0$. D'autre part, nous pouvons choisir
$\mathcal{E}$ et $\mathcal{G}$ de sorte qu'il existe une surjection
$\mathcal{E}^\vee \boxtimes \mathcal{G}^\vee \to H^i(K)$
d'après le lemme \ref{lemma-on-product}.
Dans ce cas, le membre de gauche des égalités n'est pas nul.
On conclut donc que $H^i(K) = 0$ pour $i < 0$.

\medskip\noindent
Soit $i$ le plus grand entier tel que $H^i(K)$ soit non nul.
Le même argument, avec $\mathcal{E}$ et $\mathcal{G}$ à support de dimension $0$,
montre que $i \leq 0$.
On conclut donc que $K$ est donné par un unique
$\mathcal{O}_{X \times X}$-module cohérent $\mathcal{K}$ placé en degré $0$.

\medskip\noindent
Puisque $R\text{pr}_{2, *}(\text{pr}_1^*\mathcal{F} \otimes \mathcal{K})$
est $\mathcal{F}$, en prenant $\mathcal{F}$ à support en des points fermés,
nous voyons que le support de $\mathcal{K}$ est fini sur $X$ via
$\text{pr}_2$. Comme $R\text{pr}_{2, *}(\mathcal{K}) \cong \mathcal{O}_X$,
on conclut, d'après Foncteurs et morphismes, lemme
\ref{functors-lemma-pushforward-invertible-pre},
que $\mathcal{K} = s_*\mathcal{O}_X$ pour une section $s : X \to X \times X$
de la seconde projection. Alors $\Phi_K(M) = f^*M$, où
$f = \text{pr}_1 \circ s$, et cela ne peut se produire que si $s$
est le morphisme diagonal, comme souhaité.
\end{proof}








\section{Une catégorie de noyaux de Fourier-Mukai}
\label{section-category-Fourier-Mukai-kernels}

\noindent
Soit $S$ un schéma. Nous affirmons qu'il existe une catégorie
dont
\begin{enumerate}
\item les objets sont les schémas propres et lisses sur $S$ ;
\item les morphismes de $X$ vers $Y$ sont les classes d'isomorphisme
des objets de $D_{perf}(\mathcal{O}_{X \times_S Y})$ ;
\item la composée de la classe d'isomorphisme de
$K \in D_{perf}(\mathcal{O}_{X \times_S Y})$
et de la classe d'isomorphisme de $K'$ dans $D_{perf}(\mathcal{O}_{Y \times_S Z})$
est la classe d'isomorphisme de
$$
R\text{pr}_{13, *}(
L\text{pr}_{12}^*K
\otimes_{\mathcal{O}_{X \times_S Y \times_S Z}}^\mathbf{L}
L\text{pr}_{23}^*K')
$$
qui appartient à $D_{perf}(\mathcal{O}_{X \times_S Z})$ d'après
Catégories dérivées des schémas, lemme
\ref{perfect-lemma-flat-proper-perfect-direct-image-general} ;
\item le morphisme identité de $X$ vers $X$ est la
classe d'isomorphisme de $\Delta_{X/S, *}\mathcal{O}_X$,
qui appartient à $D_{perf}(\mathcal{O}_{X \times_S X})$
d'après Compléments sur les morphismes, lemme
\ref{more-morphisms-lemma-perfect-closed-immersion-perfect-direct-image},
et le fait que $\Delta_{X/S}$ est un morphisme parfait d'après
Diviseurs, lemme
\ref{divisors-lemma-immersion-smooth-into-smooth-regular-immersion} et
Compléments sur les morphismes, lemme \ref{more-morphisms-lemma-regular-immersion-perfect}.
\end{enumerate}
Vérifions l'associativité de la composition
des morphismes ; nous omettons de vérifier que les morphismes identité
sont bien des identités. Pour cela, supposons donnés
$X, Y, Z, W$ et
$c \in D_{perf}(\mathcal{O}_{X \times_S Y})$,
$c' \in D_{perf}(\mathcal{O}_{Y \times_S Z})$, et
$c'' \in D_{perf}(\mathcal{O}_{Z \times_S W})$. Alors
\begin{align*}
c'' \circ (c' \circ c)
& \cong
\text{pr}^{134}_{14, *}(
\text{pr}^{134, *}_{13}
\text{pr}^{123}_{13, *}(\text{pr}^{123, *}_{12}c \otimes
\text{pr}^{123, *}_{23}c')
\otimes \text{pr}^{134, *}_{34}c'') \\
& \cong
\text{pr}^{134}_{14, *}(
\text{pr}^{1234}_{134, *}
\text{pr}^{1234, *}_{123}(\text{pr}^{123, *}_{12}c \otimes
\text{pr}^{123, *}_{23}c')
\otimes \text{pr}^{134, *}_{34}c'') \\
& \cong
\text{pr}^{134}_{14, *}(
\text{pr}^{1234}_{134, *}
(\text{pr}^{1234, *}_{12}c \otimes
\text{pr}^{1234, *}_{23}c')
\otimes \text{pr}^{134, *}_{34}c'') \\
& \cong
\text{pr}^{134}_{14, *}
\text{pr}^{1234}_{134, *}
((\text{pr}^{1234, *}_{12}c \otimes
\text{pr}^{1234, *}_{23}c')
\otimes \text{pr}^{1234, *}_{34}c'') \\
& \cong
\text{pr}^{1234}_{14, *}(
(\text{pr}^{1234, *}_{12}c \otimes
\text{pr}^{1234, *}_{23}c') \otimes
\text{pr}^{1234, *}_{34}c'')
\end{align*}
Nous employons ici la notation
$$
p^{1234}_{134} : X \times_S Y \times_S Z \times_S W
\to X \times_S Z \times_S W
\quad\text{et}\quad
p^{134}_{14} : X \times_S Z \times_S W \to X \times_S W
$$
pour les projections, et de même pour les autres indices.
Nous écrivons aussi $\text{pr}_*$ au lieu de $R\text{pr}_*$ et
$\text{pr}^*$ au lieu de $L\text{pr}^*$, et nous omettons
tous les indices supérieurs et inférieurs sur $\otimes$.
La première égalité est la définition de la composition.
La deuxième égalité vaut parce que
$\text{pr}^{134, *}_{13} \text{pr}^{123}_{13, *} =
\text{pr}^{1234}_{134, *} \text{pr}^{1234, *}_{123}$
par changement de base (Catégories dérivées des schémas, lemme
\ref{perfect-lemma-compare-base-change}).
La troisième égalité vaut parce que les images inverses se composent
de la manière attendue et commutent aux produits tensoriels, voir
Cohomologie, lemmes \ref{cohomology-lemma-derived-pullback-composition} et
\ref{cohomology-lemma-pullback-tensor-product}.
La quatrième égalité résulte de la « formule de projection » pour
$p^{1234}_{134}$, voir Catégories dérivées des schémas, lemme
\ref{perfect-lemma-cohomology-base-change}.
La cinquième égalité exprime la compatibilité de l'image directe propre
avec la composition, voir
Cohomologie, lemme \ref{cohomology-lemma-derived-pushforward-composition}.
Comme le produit tensoriel est associatif,
cela achève la démonstration de l'associativité de la composition.

\begin{lemma}
\label{lemma-base-change-is-functor}
Soit $S' \to S$ un morphisme de schémas.
La règle qui associe
\begin{enumerate}
\item à tout schéma propre et lisse $X$ sur $S$ le schéma $X' =  S' \times_S X$, et
\item à la classe d'isomorphisme d'un objet $K$
de $D_{perf}(\mathcal{O}_{X \times_S Y})$ la classe d'isomorphisme de
$L(X' \times_{S'} Y' \to X \times_S Y)^*K$
dans $D_{perf}(\mathcal{O}_{X' \times_{S'} Y'})$
\end{enumerate}
est un foncteur de la catégorie définie pour $S$ vers la catégorie
définie pour $S'$.
\end{lemma}

\begin{proof}
Pour le voir, supposons donnés $X, Y, Z$ et
$K \in D_{perf}(\mathcal{O}_{X \times_S Y})$ et
$M \in D_{perf}(\mathcal{O}_{Y \times_S Z})$.
Notons
$K' \in D_{perf}(\mathcal{O}_{X' \times_{S'} Y'})$ et
$M' \in D_{perf}(\mathcal{O}_{Y' \times_{S'} Z'})$
leurs images inverses comme dans l'énoncé du lemme.
Le carré
$$
\xymatrix{
X' \times_{S'} Y' \times_{S'} Z' \ar[r] \ar[d]_{\text{pr}'_{13}} &
X \times_S Y \times_S Z \ar[d]^{\text{pr}_{13}} \\
X' \times_{S'} Z' \ar[r] &
X \times_S Z
}
$$
est cartésien et $\text{pr}_{13}$ est propre et lisse.
D'après Catégories dérivées des schémas, lemme
\ref{perfect-lemma-flat-proper-perfect-direct-image-general},
l'image inverse dérivée par la flèche horizontale inférieure
du composé
$$
R\text{pr}_{13, *}(
L\text{pr}_{12}^*K
\otimes_{\mathcal{O}_{X \times_S Y \times_S Z}}^\mathbf{L}
L\text{pr}_{23}^*M)
$$
est bien (canoniquement) isomorphe à
$$
R\text{pr}'_{13, *}(
L(\text{pr}'_{12})^*K'
\otimes_{\mathcal{O}_{X' \times_{S'} Y' \times_{S'} Z'}}^\mathbf{L}
L(\text{pr}'_{23})^*M')
$$
comme voulu. Nous omettons quelques détails.
\end{proof}





\section{Équivalences relatives}
\label{section-relative-equivalences}

\noindent
Dans cette section, nous démontrons quelques lemmes sur la notion suivante.

\begin{definition}
\label{definition-relative-equivalence-kernel}
Soit $S$ un schéma. Soient $X \to S$ et $Y \to S$ des morphismes lisses et propres.
Un objet $K \in D_{perf}(\mathcal{O}_{X \times_S Y})$
est appelé {\it noyau de Fourier-Mukai d'une équivalence relative
de $X$ vers $Y$ au-dessus de $S$}
s'il existe un objet $K' \in D_{perf}(\mathcal{O}_{X \times_S Y})$
tel que
$$
\Delta_{X/S, *}\mathcal{O}_X \cong
R\text{pr}_{13, *}(L\text{pr}_{12}^*K
\otimes_{\mathcal{O}_{X \times_S Y \times_S X}}^\mathbf{L}
L\text{pr}_{23}^*K')
$$
dans $D(\mathcal{O}_{X \times_S X})$ et
$$
\Delta_{Y/S, *}\mathcal{O}_Y \cong
R\text{pr}_{13, *}(L\text{pr}_{12}^*K'
\otimes_{\mathcal{O}_{Y \times_S X \times_S Y}}^\mathbf{L}
L\text{pr}_{23}^*K)
$$
dans $D(\mathcal{O}_{Y \times_S Y})$. Autrement dit, la classe d'isomorphisme
de $K$ définit une flèche inversible dans la catégorie définie à la
section \ref{section-category-Fourier-Mukai-kernels}.
\end{definition}

\noindent
La terminologie est volontairement lourde.

\begin{lemma}
\label{lemma-equivalences-rek}
Avec les notations de la définition \ref{definition-relative-equivalence-kernel},
soit $K$ le noyau de Fourier-Mukai d'une équivalence relative de $X$
vers $Y$ au-dessus de $S$. Alors les foncteurs de Fourier-Mukai correspondants
$\Phi_K : D_\QCoh(\mathcal{O}_X) \to D_\QCoh(\mathcal{O}_Y)$
(lemme \ref{lemma-fourier-Mukai-QCoh})
et $\Phi_K : D_{perf}(\mathcal{O}_X) \to D_{perf}(\mathcal{O}_Y)$
(lemme \ref{lemma-fourier-mukai})
sont des équivalences.
\end{lemma}

\begin{proof}
Cela résulte immédiatement du lemme \ref{lemma-compose-fourier-mukai} et de
l'exemple \ref{example-diagonal-fourier-mukai}.
\end{proof}

\begin{lemma}
\label{lemma-base-change-rek}
Avec les notations de la définition \ref{definition-relative-equivalence-kernel},
soit $K$ le noyau de Fourier-Mukai d'une équivalence relative de $X$
vers $Y$ au-dessus de $S$. Soit $S_1 \to S$ un morphisme de schémas. Posons
$X_1 = S_1 \times_S X$ et $Y_1 = S_1 \times_S Y$. Alors l'image inverse
$K_1 = L(X_1 \times_{S_1} Y_1 \to X \times_S Y)^*K$ est
le noyau de Fourier-Mukai d'une équivalence relative de $X_1$
vers $Y_1$ au-dessus de $S_1$.
\end{lemma}

\begin{proof}
Soit $K' \in D_{perf}(\mathcal{O}_{Y \times_S X})$ l'objet dont l'existence est
supposée dans la définition \ref{definition-relative-equivalence-kernel}.
Notons $K'_1$ l'image inverse de $K'$ par
$Y_1 \times_{S_1} X_1 \to Y \times_S X$.
Il suffit alors de démontrer que
$$
\Delta_{X_1/S_1, *}\mathcal{O}_X \cong
R\text{pr}_{13, *}(L\text{pr}_{12}^*K_1
\otimes_{\mathcal{O}_{X_1 \times_{S_1} Y_1 \times_{S_1} X_1}}^\mathbf{L}
L\text{pr}_{23}^*K_1')
$$
dans $D(\mathcal{O}_{X_1 \times_{S_1} X_1})$, et de même pour l'autre
condition. Comme
$$
\xymatrix{
X_1 \times_{S_1} Y_1 \times_{S_1} X_1 \ar[r] \ar[d]_{\text{pr}_{13}} &
X \times_S Y \times_S X \ar[d]^{\text{pr}_{13}} \\
X_1 \times_{S_1} X_1 \ar[r] &
X \times_S X
}
$$
est cartésien, il suffit, d'après Catégories dérivées des schémas, lemme
\ref{perfect-lemma-flat-proper-perfect-direct-image-general},
de montrer que
$$
\Delta_{X_1/S_1, *}\mathcal{O}_{X_1}
\cong
L(X_1 \times_{S_1} X_1 \to X \times_S X)^*\Delta_{X/S, *}\mathcal{O}_X 
$$
Cela résultera à son tour de la tor-indépendance de $X$ et de
$X_1 \times_{S_1} X_1$ au-dessus de $X \times_S X$, voir
Catégories dérivées des schémas, lemme \ref{perfect-lemma-compare-base-change}.
Cette tor-indépendance se vérifie directement, mais résulte aussi
du résultat plus général de Compléments sur les morphismes, lemme
\ref{more-morphisms-lemma-case-of-tor-independence}, appliqué au carré
de sommets $X, X, X, S$ et à son changement de base par $S_1 \to S$.
\end{proof}

\begin{lemma}
\label{lemma-descend-rek}
Soit $S = \lim_{i \in I} S_i$ la limite d'un système inductif de schémas
dont les morphismes de transition affines sont $g_{i'i} : S_{i'} \to S_i$.
Supposons que $S_i$ soit quasi-compact et quasi-séparé pour tout $i \in I$.
Soit $0 \in I$. Soient $X_0 \to S_0$ et $Y_0 \to S_0$ des morphismes lisses et propres.
Posons $X_i = S_i \times_{S_0} X_0$ pour $i \geq 0$,
et $X = S \times_{S_0} X_0$ ; de même pour $Y_0$. Si $K$ est le
noyau de Fourier-Mukai d'une équivalence relative de $X$ vers $Y$ au-dessus de $S$,
alors, pour un certain $i \geq 0$, il existe un
noyau de Fourier-Mukai d'une équivalence relative de $X_i$ vers $Y_i$ au-dessus de $S_i$.
\end{lemma}

\begin{proof}
Soit $K' \in D_{perf}(\mathcal{O}_{Y \times_S X})$ l'objet dont l'existence est
supposée dans la définition \ref{definition-relative-equivalence-kernel}.
Puisque $X \times_S Y = \lim X_i \times_{S_i} Y_i$, il existe un
$i$ et des objets $K_i$ et $K'_i$ dans
$D_{perf}(\mathcal{O}_{Y_i \times_{S_i} X_i})$
dont les images inverses sur $Y \times_S X$ sont $K$ et $K'$.
Voir Catégories dérivées des schémas, lemme \ref{perfect-lemma-descend-perfect}.
D'après Catégories dérivées des schémas, lemme
\ref{perfect-lemma-flat-proper-perfect-direct-image-general},
l'objet
$$
R\text{pr}_{13, *}(L\text{pr}_{12}^*K_i
\otimes_{\mathcal{O}_{X_i \times_{S_i} Y_i \times_{S_i} X_i}}^\mathbf{L}
L\text{pr}_{23}^*K_i')
$$
est parfait, et son image inverse sur $X \times_S X$ est égale à
$$
R\text{pr}_{13, *}(L\text{pr}_{12}^*K
\otimes_{\mathcal{O}_{X \times_S Y \times_S X}}^\mathbf{L}
L\text{pr}_{23}^*K') \cong \Delta_{X/S, *}\mathcal{O}_X
$$
Voir la démonstration du lemme \ref{lemma-base-change-rek}.
D'autre part, puisque $X_i \to S$ est lisse et séparé,
l'objet
$$
\Delta_{i, *}\mathcal{O}_{X_i}
$$
de $D(\mathcal{O}_{X_i \times_{S_i} X_i})$ est lui aussi parfait
(d'après Compléments sur les morphismes, lemmes
\ref{more-morphisms-lemma-smooth-diagonal-perfect} et
\ref{more-morphisms-lemma-perfect-proper-perfect-direct-image}), et
son image inverse sur $X \times_S X$ est égale à
$$
\Delta_{X/S, *}\mathcal{O}_X
$$
Voir la démonstration du lemme \ref{lemma-base-change-rek}. Ainsi, d'après
Catégories dérivées des schémas, lemme \ref{perfect-lemma-descend-perfect},
quitte à augmenter $i$, on peut supposer que
$$
\Delta_{i, *}\mathcal{O}_{X_i} \cong
R\text{pr}_{13, *}(L\text{pr}_{12}^*K_i
\otimes_{\mathcal{O}_{X_i \times_{S_i} Y_i \times_{S_i} X_i}}^\mathbf{L}
L\text{pr}_{23}^*K_i')
$$
comme voulu. Le même raisonnement s'applique après interversion des rôles de $K$ et $K'$.
\end{proof}








\section{Absence de déformations}
\label{section-no-deformations}

\noindent
Le titre de cette section renvoie au lemme \ref{lemma-no-deformations}.

\begin{lemma}
\label{lemma-deform-koszul}
Soit $(R, \mathfrak m, \kappa) \to (A, \mathfrak n, \lambda)$
un homomorphisme local plat d'anneaux locaux
essentiellement de présentation finie.
Soit $\overline{f}_1, \ldots, \overline{f}_r \in \mathfrak n/\mathfrak m A
\subset A/\mathfrak m A$ une suite régulière. Soit $K \in D(A)$. Supposons que
\begin{enumerate}
\item $K$ soit parfait,
\item $K \otimes_A^\mathbf{L} A/\mathfrak m A$ soit isomorphe dans
$D(A/\mathfrak m A)$ au
complexe de Koszul associé à $\overline{f}_1, \ldots, \overline{f}_r$.
\end{enumerate}
Alors $K$ est isomorphe dans $D(A)$ à un complexe de Koszul associé à une suite
régulière $f_1, \ldots, f_r \in A$ qui relève les éléments donnés
$\overline{f}_1, \ldots, \overline{f}_r$. De plus, $A/(f_1, \ldots, f_r)$
est plat sur $R$.
\end{lemma}

\begin{proof}
Utilisons des complexes de chaînes dans la démonstration de ce lemme.
Le complexe de Koszul $K_\bullet(\overline{f}_1, \ldots, \overline{f}_r)$
est défini dans Compléments d'algèbre, définition
\ref{more-algebra-definition-koszul-complex}.
D'après Compléments d'algèbre, lemme \ref{more-algebra-lemma-lift-complex-stably-frees},
on peut représenter $K$ par un complexe
$$
K_\bullet :
A \to A^{\oplus r} \to \ldots \to A^{\oplus r} \to A
$$
dont le produit tensoriel avec $A/\mathfrak mA$ est égal (!)
à $K_\bullet(\overline{f}_1, \ldots, \overline{f}_r)$.
Notons $f_1, \ldots, f_r \in A$ les composantes de la
flèche $A^{\oplus r} \to A$. Ces $f_i$ relèvent les
$\overline{f}_i$. D'après Algèbre, lemme
\ref{algebra-lemma-grothendieck-regular-sequence-general}
$f_1, \ldots, f_r$ forment une suite régulière dans $A$ et $A/(f_1, \ldots, f_r)$
est plat sur $R$. Posons $J = (f_1, \ldots, f_r) \subset A$.
Considérons le diagramme
$$
\xymatrix{
K_\bullet \ar[rd] \ar@{..>}[rr]_{\varphi_\bullet} & &
K_\bullet(f_1, \ldots, f_r) \ar[ld] \\
& A/J
}
$$
Comme $f_1, \ldots, f_r$ forment une suite régulière, la flèche sud-ouest
est un quasi-isomorphisme (voir
Compléments d'algèbre, lemme \ref{more-algebra-lemma-regular-koszul-regular}).
On peut donc trouver la flèche pointillée qui fait commuter le
diagramme, par exemple grâce à
Algèbre, lemme \ref{algebra-lemma-compare-resolutions}.
En réduisant modulo $\mathfrak m$, on obtient un diagramme commutatif
$$
\xymatrix{
K_\bullet(\overline{f}_1, \ldots, \overline{f}_r)
\ar[rd] \ar[rr]_{\overline{\varphi}_\bullet} & &
K_\bullet(\overline{f}_1, \ldots, \overline{f}_r) \ar[ld] \\
& (A/\mathfrak m A)/(\overline{f}_1, \ldots, \overline{f}_r)
}
$$
par notre choix de $K_\bullet$. Ainsi, $\overline{\varphi}$ est un isomorphisme
dans la catégorie dérivée $D(A/\mathfrak m A)$. Il s'ensuit que
$\overline{\varphi} \otimes_{A/\mathfrak m A}^\mathbf{L} \lambda$
est un isomorphisme. Puisque $\overline{f}_i \in \mathfrak n / \mathfrak m A$,
on voit que
$$
\text{Tor}_i^{A/\mathfrak m A}(
K_\bullet(\overline{f}_1, \ldots, \overline{f}_r), \lambda)
=
K_i(\overline{f}_1, \ldots, \overline{f}_r) \otimes_{A/\mathfrak m A} \lambda
$$
Par conséquent, $\varphi_i \bmod \mathfrak n$ est inversible.
Comme $A$ est local, cela signifie que $\varphi_i$ est un
isomorphisme, ce qui achève la démonstration.
\end{proof}

\begin{lemma}
\label{lemma-limit-arguments}
Soit $R \to S$ un morphisme plat de type fini d'anneaux noethériens.
Soit $\mathfrak q \subset S$ un idéal premier au-dessus de
$\mathfrak p \subset R$. Soit $K \in D(S)$ parfait.
Soit $f_1, \ldots, f_r \in \mathfrak q S_\mathfrak q$
une suite régulière telle que $S_\mathfrak q/(f_1, \ldots, f_r)$
soit plat sur $R$ et telle que
$K \otimes_S^\mathbf{L} S_\mathfrak q$ soit isomorphe au
complexe de Koszul associé à $f_1, \ldots, f_r$. Alors il existe
$g \in S$, $g \not \in \mathfrak q$, tel que
\begin{enumerate}
\item $f_1, \ldots, f_r$ soient les images de
$f'_1, \ldots, f'_r \in S_g$,
\item $f'_1, \ldots, f'_r$ forment une suite régulière dans $S_g$,
\item $S_g/(f'_1, \ldots, f'_r)$ soit plat sur $R$,
\item $K \otimes_S^\mathbf{L} S_g$ soit isomorphe au
complexe de Koszul associé à $f_1, \ldots, f_r$.
\end{enumerate}
\end{lemma}

\begin{proof}
La définition des localisés permet de trouver $g \in S$, $g \not \in \mathfrak q$,
qui vérifie (1). Après avoir remplacé $g$ par
$gg'$ pour un certain $g' \in S$, $g' \not \in \mathfrak q$,
on peut supposer que (2) est vérifiée ; voir
Algèbre, lemme \ref{algebra-lemma-regular-sequence-in-neighbourhood}.
D'après Algèbre, théorème \ref{algebra-theorem-openness-flatness},
$S_g/(f'_1, \ldots, f'_r)$ est plat sur $R$
dans un voisinage ouvert de $\mathfrak q$.
Après avoir de nouveau remplacé $g$ par $gg'$ pour un certain
$g' \in S$, $g' \not \in \mathfrak q$, on peut donc supposer que (3) est aussi vérifiée.
Enfin, on obtient (4) après un remplacement supplémentaire grâce à
Compléments d'algèbre, lemme \ref{more-algebra-lemma-colimit-perfect-complexes}.
\end{proof}

\noindent
Pour une généralisation du lemme suivant, voir
Compléments sur les morphismes d'espaces, lemme
\ref{spaces-more-morphisms-lemma-where-isomorphism}.

\begin{lemma}
\label{lemma-isomorphism-in-neighbourhood}
Soit $S$ un schéma noethérien. Soit $s \in S$.
Soit $p : X \to Y$ un morphisme de schémas sur $S$.
Supposons que
\begin{enumerate}
\item $Y \to S$ et $X \to S$ soient propres,
\item $X$ soit plat sur $S$,
\item $X_s \to Y_s$ soit un isomorphisme.
\end{enumerate}
Alors il existe un voisinage ouvert $U \subset S$ de $s$
tel que le changement de base $X_U \to Y_U$ soit un isomorphisme.
\end{lemma}

\begin{proof}
Le morphisme $p$ est propre d'après Morphismes, lemme
\ref{morphisms-lemma-closed-immersion-proper}.
D'après Cohomologie des schémas, lemme
\ref{coherent-lemma-proper-finite-fibre-finite-in-neighbourhood}
il existe un ouvert $Y_s \subset V \subset Y$ tel que
$p|_{p^{-1}(V)} : p^{-1}(V) \to V$ soit fini.
D'après Compléments sur les morphismes, théorème
\ref{more-morphisms-theorem-criterion-flatness-fibre-Noetherian}
il existe un ouvert $X_s \subset U \subset X$ tel que
$p|_U : U \to Y$ soit plat. Après avoir retranché les images de
$X \setminus U$ et de $Y \setminus V$ (qui sont des parties fermées
ne contenant pas $s$), on peut supposer que $p$ est plat et fini.
Alors $p$ est ouvert (Morphismes, lemme \ref{morphisms-lemma-fppf-open})
et $Y_s \subset p(X) \subset Y$ ; après avoir rétréci $S$,
on peut donc supposer que $p$ est surjectif.
Comme $p_s : X_s \to Y_s$ est un isomorphisme, l'application
$$
p^\sharp : \mathcal{O}_Y \longrightarrow p_*\mathcal{O}_X
$$
de $\mathcal{O}_Y$-modules cohérents ($p$ est fini)
devient un isomorphisme après image inverse par $i : Y_s \to Y$
(par exemple d'après Cohomologie des schémas, lemme
\ref{coherent-lemma-affine-base-change}).
Le lemme de Nakayama implique alors que
$\mathcal{O}_{Y, y} \to (p_*\mathcal{O}_X)_y$ est surjective
pour tout $y \in Y_s$. Il existe donc un ouvert $Y_s \subset V \subset Y$
tel que $p^\sharp|_V$ soit surjective
(Modules, lemme \ref{modules-lemma-finite-type-surjective-on-stalk}).
Après avoir de nouveau rétréci $S$, on peut donc supposer que
$p^\sharp$ est surjective, ce qui signifie que $p$ est une immersion
fermée (puisque $p$ est déjà fini).
Ainsi, $p$ est une immersion fermée plate et surjective
de schémas noethériens, donc un isomorphisme ; voir
Morphismes, section \ref{morphisms-section-flat-closed-immersions}.
\end{proof}

\begin{lemma}
\label{lemma-no-deformations}
Soit $k$ un corps. Soit $S$ un schéma de type fini sur $k$
muni d'un point $k$-rationnel $s$. Soit $Y \to S$ un morphisme lisse et propre.
Soit $X = Y_s \times S \to S$ la famille constante de fibre
$Y_s$. Soit $K$ le noyau de Fourier-Mukai d'une équivalence relative
de $X$ vers $Y$ au-dessus de $S$. Supposons que l'on ait
$$
L(Y_s \times_S Y_s \to X \times_S Y)^*K \cong 
\Delta_{Y_s/k, *} \mathcal{O}_{Y_s}
$$
dans $D(\mathcal{O}_{Y_s \times Y_s})$. Alors il existe un voisinage ouvert
$s \in U \subset S$ tel que $Y|_U$ soit isomorphe à $Y_s \times U$ sur $U$.
\end{lemma}

\begin{proof}
Notons $i : Y_s \times Y_s = X_s \times Y_s \to X \times_S Y$
l'immersion fermée naturelle. (Désormais, nous noterons $Y_s$, et non $X_s$,
la fibre de $X$ au-dessus de $s$.) Soit
$z \in Y_s \times Y_s = (X \times_S Y)_s \subset X \times_S Y$
un point fermé. Comme indiqué, nous considérons $z$ à la fois comme un point fermé
de $Y_s \times Y_s$ et comme un point fermé de $X \times_S Y$.

\medskip\noindent
Cas I : $z \not \in \Delta_{Y_s/k}(Y_s)$. Notons $\mathcal{O}_z$
le $\mathcal{O}_{Y_s \times Y_s}$-module cohérent supporté en $z$
dont la fibre est $\kappa(z)$. Alors $i_*\mathcal{O}_z$ est le
$\mathcal{O}_{X \times_S Y}$-module cohérent supporté en $z$
dont la fibre est $\kappa(z)$. Notre hypothèse signifie que
$$
K \otimes_{\mathcal{O}_{X \times_S Y}}^\mathbf{L} i_*\mathcal{O}_z =
Li^*K \otimes_{\mathcal{O}_{Y_s \times Y_s}}^\mathbf{L} \mathcal{O}_z = 0
$$
D'après le lemme \ref{lemma-orthogonal-point-sheaf}, il existe donc
un voisinage ouvert $U(z) \subset X \times_S Y$ de $z$
tel que $K|_{U(z)} = 0$. Dans ce cas, on pose $Z(z) = \emptyset$
comme sous-schéma fermé de $U(z)$.

\medskip\noindent
Cas II : $z \in \Delta_{Y_s/k}(Y_s)$. Comme $Y_s$ est lisse sur $k$,
on sait que $\Delta_{Y_s/k} : Y_s \to Y_s \times Y_s$ est une
immersion régulière ; voir Compléments sur les morphismes, lemme
\ref{more-morphisms-lemma-smooth-diagonal-perfect}.
Choisissons une suite régulière $\overline{f}_1, \ldots, \overline{f}_r \in
\mathcal{O}_{Y_s \times Y_s, z}$ engendrant l'idéal de
$\Delta_{Y_s/k}(Y_s)$. Comme une suite régulière est régulière au sens de Koszul
(Compléments d'algèbre, lemme \ref{more-algebra-lemma-regular-koszul-regular}),
notre hypothèse signifie que
$$
K_z \otimes_{\mathcal{O}_{X \times_S Y, z}}^\mathbf{L}
\mathcal{O}_{Y_s \times Y_s, z}
\in D(\mathcal{O}_{Y_s \times Y_s, z})
$$
est représenté par le complexe de Koszul associé à
$\overline{f}_1, \ldots, \overline{f}_r$ sur
$\mathcal{O}_{Y_s \times Y_s, z}$.
En appliquant le lemme \ref{lemma-deform-koszul} à
$\mathcal{O}_{S, s} \to \mathcal{O}_{X \times_S Y, z}$,
on conclut que $K_z \in D(\mathcal{O}_{X \times_S Y, z})$ est
représenté par le complexe de Koszul associé à une suite régulière
$f_1, \ldots, f_r \in \mathcal{O}_{X \times_S Y, z}$
qui relève la suite régulière
$\overline{f}_1, \ldots, \overline{f}_r$
et que, de plus, $\mathcal{O}_{X \times_S Y}/(f_1, \ldots, f_r)$
est plat sur $\mathcal{O}_{S, s}$.
Par des arguments de passage à la limite (lemme \ref{lemma-limit-arguments}),
on conclut qu'il existe un voisinage ouvert affine
$U(z) \subset X \times_S Y$ de $z$ et un sous-schéma fermé
$Z(z) \subset U(z)$ tels que
\begin{enumerate}
\item $Z(z) \to U(z)$ soit une immersion fermée régulière,
\item $K|_{U(z)}$ soit quasi-isomorphe à $\mathcal{O}_{Z(z)}$,
\item $Z(z) \to S$ soit plat,
\item $Z(z)_s = \Delta_{Y_s/k}(Y_s) \cap U(z)_s$
comme sous-schémas fermés de $U(z)_s$.
\end{enumerate}

\noindent
D'après la propriété (2), pour $z, z' \in Y_s \times Y_s$, on
a $Z(z) \cap U(z') = Z(z') \cap U(z)$ en tant que sous-schémas fermés.
On obtient donc un voisinage ouvert
$$
U = \bigcup\nolimits_{z \in Y_s \times Y_s\text{ fermé}} U(z)
$$
de $Y_s \times Y_s$ dans $X \times_S Y$ et un sous-schéma fermé $Z \subset U$
tels que (1) $Z \to U$ soit une immersion fermée régulière,
(2) $Z \to S$ soit plat et (3) $Z_s = \Delta_{Y_s/k}(Y_s)$.
Comme $X \times_S Y \to S$ est propre, après avoir remplacé $S$
par un voisinage ouvert de $s$, on peut supposer que $U = X \times_S Y$.
Comme les projections $Z_s \to Y_s$ et $Z_s \to X_s$
sont des isomorphismes, on conclut qu'après avoir rétréci $S$,
on peut supposer que $Z \to Y$ et $Z \to X$ sont des isomorphismes ; voir
le lemme \ref{lemma-isomorphism-in-neighbourhood}.
Ceci achève la démonstration.
\end{proof}

\begin{lemma}
\label{lemma-no-deformations-better}
Soit $k$ un corps algébriquement clos. Soit $X$
un schéma lisse et propre sur $k$.
Soit $f : Y \to S$ un morphisme lisse et propre, où $S$ est de type fini sur $k$.
Soit $K$ le noyau de Fourier-Mukai d'une équivalence relative
de $X \times S$ vers $Y$ au-dessus de $S$. Alors $S$ peut être recouvert par des
sous-schémas ouverts $U$ tels qu'il existe un $U$-isomorphisme
$f^{-1}(U) \cong Y_0 \times U$ pour un certain $Y_0$ propre et lisse sur $k$.
\end{lemma}

\begin{proof}
Choisissons un point fermé $s \in S$. Comme $k$ est algébriquement clos,
c'est un point $k$-rationnel. Posons $Y_0 = Y_s$. La restriction
$K_0$ de $K$ à $X \times Y_0$ est le noyau de Fourier-Mukai d'une
équivalence relative de $X$ vers $Y_0$ au-dessus de $\Spec(k)$ d'après le
lemme \ref{lemma-base-change-rek}. Soit $K'_0$,
dans $D_{perf}(\mathcal{O}_{Y_0 \times X})$,
l'objet dont l'existence est postulée dans la
définition \ref{definition-relative-equivalence-kernel}.
Alors $K'_0$ est le noyau de Fourier-Mukai d'une
équivalence relative de $Y_0$ vers $X$ au-dessus de $\Spec(k)$,
par la symétrie inhérente à la
définition \ref{definition-relative-equivalence-kernel}.
Ainsi, d'après le
lemme \ref{lemma-base-change-rek},
on voit que l'image inverse
$$
M = (Y_0 \times X \times S \to Y_0 \times X)^*K'_0
$$
sur $(Y_0 \times S) \times_S (X \times S) = Y_0 \times X \times S$
est le noyau de Fourier-Mukai d'une
équivalence relative de $Y_0 \times S$ vers $X \times S$ au-dessus de $S$.
Considérons maintenant le noyau
$$
K_{new} =
R\text{pr}_{13, *}(L\text{pr}_{12}^*M
\otimes_{\mathcal{O}_{(Y_0 \times S) \times_S (X \times S)
\times_S Y}}^\mathbf{L}
L\text{pr}_{23}^*K)
$$
sur $(Y_0 \times S) \times_S Y$. C'est le noyau de Fourier-Mukai d'une
équivalence relative de $Y_0 \times S$ vers $Y$ au-dessus de $S$, car il est
la composée de deux flèches inversibles dans
la catégorie construite à la
section \ref{section-category-Fourier-Mukai-kernels}.
De plus, cette composition commute au changement de base
(lemme \ref{lemma-base-change-is-functor}).
On voit donc que l'image inverse de $K_{new}$ sur
$((Y_0 \times S) \times_S Y)_s = Y_0 \times Y_0$
est égale à la composée de $K_0$ et de $K'_0$,
et donc à l'identité dans cette catégorie.
Autrement dit, on a
$$
L(Y_0 \times Y_0 \to (Y_0 \times S) \times_S Y)^*K_{new}
\cong
\Delta_{Y_0/k, *}\mathcal{O}_{Y_0}
$$
D'après le lemme \ref{lemma-no-deformations}, on conclut donc que $Y \to S$
est isomorphe à $Y_0 \times S$ dans un voisinage ouvert de $s$.
Ceci achève la démonstration.
\end{proof}






\section{Dénombrabilité}
\label{section-countability}

\noindent
Dans cette section, nous démontrons quelques lemmes élémentaires sur la dénombrabilité
de certains ensembles. Soit $\mathcal{C}$ une catégorie. Dans cette section,
nous dirons que $\mathcal{C}$ est {\it dénombrable} si
\begin{enumerate}
\item pour tous $X, Y \in \Ob(\mathcal{C})$, l'ensemble
$\Mor_\mathcal{C}(X, Y)$ est dénombrable, et
\item l'ensemble des classes d'isomorphisme d'objets de $\mathcal{C}$
est dénombrable.
\end{enumerate}

\begin{lemma}
\label{lemma-countable-finite-type}
Soit $R$ un anneau noethérien dénombrable. Alors la catégorie des schémas de type
fini sur $R$ est dénombrable.
\end{lemma}

\begin{proof}
Omis.
\end{proof}

\begin{lemma}
\label{lemma-countable-abelian}
Soit $\mathcal{A}$ une catégorie abélienne dénombrable.
Alors $D^b(\mathcal{A})$ est dénombrable.
\end{lemma}

\begin{proof}
Il suffit de démontrer l'assertion pour $D(\mathcal{A})$, puisque les autres
en sont des sous-catégories pleines. Comme tout objet de $D(\mathcal{A})$
est un complexe d'objets de $\mathcal{A}$, il est immédiat que l'ensemble des
classes d'isomorphisme d'objets de $D^b(\mathcal{A})$ est dénombrable.
De plus, pour des complexes bornés $A^\bullet$ et $B^\bullet$ de $\mathcal{A}$,
il est clair que $\Hom_{K^b(\mathcal{A})}(A^\bullet, B^\bullet)$ est dénombrable.
On a
$$
\Hom_{D^b(\mathcal{A})}(A^\bullet, B^\bullet) =
\colim_{s : (A')^\bullet \to A^\bullet
\text{ qis et }(A')^\bullet\text{ borné}}
\Hom_{K^b(\mathcal{A})}((A')^\bullet, B^\bullet)
$$
d'après Catégories dérivées, lemme \ref{derived-lemma-bounded-derived}.
Il s'agit donc d'un ensemble dénombrable, en tant que colimite dénombrable d'ensembles dénombrables.
\end{proof}

\begin{lemma}
\label{lemma-countable-perfect}
Soit $X$ un schéma de type fini sur un anneau noethérien dénombrable.
Alors les catégories $D_{perf}(\mathcal{O}_X)$ et
$D^b_{\textit{Coh}}(\mathcal{O}_X)$ sont dénombrables.
\end{lemma}

\begin{proof}
Observons que $X$ est noethérien d'après
Morphismes, lemme \ref{morphisms-lemma-finite-type-noetherian}.
Ainsi, $D_{perf}(\mathcal{O}_X)$ est une sous-catégorie pleine de
$D^b_{\textit{Coh}}(\mathcal{O}_X)$ d'après
Catégories dérivées des schémas, lemme \ref{perfect-lemma-perfect-on-noetherian}.
Il suffit donc de démontrer
le résultat pour $D^b_{\textit{Coh}}(\mathcal{O}_X)$.
Rappelons que
$D^b_{\textit{Coh}}(\mathcal{O}_X) = D^b(\textit{Coh}(\mathcal{O}_X))$
d'après
Catégories dérivées des schémas, proposition \ref{perfect-proposition-DCoh}.
Par le lemme \ref{lemma-countable-abelian},
il suffit de démontrer que $\textit{Coh}(\mathcal{O}_X)$ est
dénombrable. Nous omettons cette vérification.
\end{proof}

\begin{lemma}
\label{lemma-countable-isos}
Soit $K$ un corps algébriquement clos.
Soit $S$ un schéma de type fini sur $K$.
Soient $X \to S$ et $Y \to S$ des morphismes de type fini.
Il existe un ensemble dénombrable $I$ et, pour tout $i \in I$, un couple
$(S_i \to S, h_i)$ possédant les propriétés suivantes :
\begin{enumerate}
\item $S_i \to S$ est un morphisme de type fini, et l'on pose
$X_i = X \times_S S_i$ et $Y_i = Y \times_S S_i$ ;
\item $h_i : X_i \to Y_i$ est un isomorphisme sur $S_i$ ;
\item pour tout point fermé $s \in S(K)$, si $X_s \cong Y_s$
sur $K = \kappa(s)$, alors $s$ appartient à l'image de $S_i \to S$
pour un certain $i$.
\end{enumerate}
\end{lemma}

\begin{proof}
Le corps $K$ est la réunion filtrante de ses sous-corps dénombrables.
Dualement, $\Spec(K)$ est la limite cofiltrante des spectres
des sous-corps dénombrables de $K$.
Le lemme \ref{limits-lemma-descend-finite-presentation} de Limites
assure donc que l'on peut trouver un sous-corps dénombrable
$k$ et des morphismes $X_0 \to S_0$ et $Y_0 \to S_0$
de schémas de type fini sur $k$ tels que
$X \to S$ et $Y \to S$ soient leurs changements de base.

\medskip\noindent
Par le lemme \ref{lemma-countable-finite-type}, il existe un ensemble dénombrable $I$ et
des couples $(S_{0, i} \to S_0, h_{0, i})$ tels que
\begin{enumerate}
\item $S_{0, i} \to S_0$ est un morphisme de type fini, et l'on pose
$X_{0, i} = X_0 \times_{S_0} S_{0, i}$ et
$Y_{0, i} = Y_0 \times_{S_0} S_{0, i}$ ;
\item $h_{0, i} : X_{0, i} \to Y_{0, i}$ est un isomorphisme sur $S_{0, i}$.
\end{enumerate}
de façon que tout couple $(T \to S_0, h_T)$, avec $T \to S_0$ de type fini
et $h_T : X_0 \times_{S_0} T \to Y_0 \times_{S_0} T$ un isomorphisme,
soit isomorphe à l'un d'eux.
Notons $(S_i \to S, h_i)$ le changement de base de $(S_{0, i} \to S_0, h_{0, i})$
par $\Spec(K) \to \Spec(k)$.
Nous affirmons que cela convient.

\medskip\noindent
Soit $s \in S(K)$ et soit $h_s : X_s \to Y_s$ un isomorphisme sur
$K = \kappa(s)$. On peut écrire $K$ comme réunion filtrante de ses
$k$-sous-algèbres de type fini. Par conséquent, d'après
Limites, proposition
\ref{limits-proposition-characterize-locally-finite-presentation}, et
le lemme \ref{limits-lemma-descend-finite-presentation},
on peut trouver une telle $k$-sous-algèbre de type fini
$K \supset A \supset k$ telle que
\begin{enumerate}
\item il existe un diagramme commutatif
$$
\xymatrix{
\Spec(K) \ar[d]_s \ar[r] &
\Spec(A) \ar[d]^{s'} \\
S \ar[r] &
S_0}
$$
pour un certain morphisme $s' : \Spec(A) \to S_0$ sur $k$ ;
\item $h_s$ est le changement de base d'un isomorphisme
$h_{s'} : X_0 \times_{S_0, s'} \Spec(A) \to
X_0 \times_{S_0, s'} \Spec(A)$ sur $A$.
\end{enumerate}
Alors $(s' : \Spec(A) \to S_0, h_{s'})$ est bien sûr isomorphe
au couple $(S_{0, i} \to S_0, h_{0, i})$ pour un certain $i \in I$.
Ceci achève la démonstration, car le diagramme commutatif
de (1) montre que $s$ appartient à l'image
du changement de base de $s'$ à $\Spec(K)$.
\end{proof}

\begin{lemma}
\label{lemma-countable-equivs}
Soit $K$ un corps algébriquement clos. Il existe un ensemble dénombrable $I$
et, pour tout $i \in I$, un système $(S_i/K, X_i \to S_i, Y_i \to S_i, M_i)$
possédant les propriétés suivantes :
\begin{enumerate}
\item $S_i$ est un schéma de type fini sur $K$ ;
\item $X_i \to S_i$ et $Y_i \to S_i$ sont des morphismes
propres et lisses ;
\item $M_i \in D_{perf}(\mathcal{O}_{X_i \times_{S_i} Y_i})$
est le noyau de Fourier-Mukai d'une équivalence relative de
$X_i$ vers $Y_i$ sur $S_i$, et
\item pour tous schémas lisses et propres $X$ et $Y$ sur $K$
tels qu'il existe une équivalence exacte et $K$-linéaire
$D_{perf}(\mathcal{O}_X) \to D_{perf}(\mathcal{O}_Y)$
il existe $i \in I$ et $s \in S_i(K)$
tels que $X \cong (X_i)_s$ et $Y \cong (Y_i)_s$.
\end{enumerate}
\end{lemma}

\begin{proof}
Choisissons un sous-corps dénombrable $k \subset K$, par exemple le sous-corps premier.
D'après les lemmes \ref{lemma-countable-finite-type} et \ref{lemma-countable-perfect},
il existe un ensemble dénombrable de classes d'isomorphisme de systèmes
sur $k$ satisfaisant aux assertions (1), (2) et (3) du lemme.
Nous pouvons donc choisir un ensemble dénombrable
$I$ et, pour tout $i \in I$, un tel système
$$
(S_{0, i}/k, X_{0, i} \to S_{0, i}, Y_{0, i} \to S_{0, i}, M_{0, i})
$$
sur $k$ tel que chaque classe d'isomorphisme soit représentée au moins une fois.
Notons $(S_i/K, X_i \to S_i, Y_i \to S_i, M_i)$ le changement de base
à $K$ du système affiché. Ce système possède les propriétés (1), (2) et (3),
voir le lemme \ref{lemma-base-change-rek}. Démontrons la propriété (4).

\medskip\noindent
Considérons des schémas lisses et propres $X$ et $Y$ sur $K$
tels qu'il existe une équivalence exacte et $K$-linéaire
$F : D_{perf}(\mathcal{O}_X) \to D_{perf}(\mathcal{O}_Y)$.
Par la proposition \ref{proposition-equivalence},
nous pouvons supposer qu'il existe un objet
$M \in D_{perf}(\mathcal{O}_{X \times Y})$
tel que $F = \Phi_M$ soit le foncteur de Fourier-Mukai correspondant.
Par le lemme \ref{lemma-fourier-mukai-flat-proper-over-noetherian},
il existe un objet $M'$ dans $D_{perf}(\mathcal{O}_{Y \times X})$
tel que $\Phi_{M'}$ soit l'adjoint à droite de $\Phi_M$.
Comme $\Phi_M$ est une équivalence, cela signifie que
$\Phi_{M'}$ est le quasi-inverse de $\Phi_M$.
Par le lemme \ref{lemma-fourier-mukai-flat-proper-over-noetherian},
les foncteurs de Fourier-Mukai définis par les objets
$$
A = R\text{pr}_{13, *}(
L\text{pr}_{12}^*M
\otimes_{\mathcal{O}_{X \times Y \times X}}^\mathbf{L}
L\text{pr}_{23}^*M')
$$
dans $D_{perf}(\mathcal{O}_{X \times X})$ et
$$
B = R\text{pr}_{13, *}(
L\text{pr}_{12}^*M'
\otimes_{\mathcal{O}_{Y \times X \times Y}}^\mathbf{L}
L\text{pr}_{23}^*M)
$$
dans $D_{perf}(\mathcal{O}_{Y \times Y})$
sont isomorphes à
$\text{id} : D_{perf}(\mathcal{O}_X) \to D_{perf}(\mathcal{O}_X)$
et à
$\text{id} : D_{perf}(\mathcal{O}_Y) \to D_{perf}(\mathcal{O}_Y)$
respectivement. Par conséquent,
$A \cong \Delta_{X/K, *}\mathcal{O}_X$ et
$B \cong \Delta_{Y/K, *}\mathcal{O}_Y$
d'après le lemme \ref{lemma-uniqueness}. Ainsi, $M$ est le
noyau de Fourier-Mukai d'une équivalence relative de $X$ vers $Y$
sur $K$, par définition.

\medskip\noindent
On peut écrire $K$ comme colimite filtrante de ses
$k$-sous-algèbres de type fini $A \subset K$. D'après
Limites, lemme \ref{limits-lemma-descend-finite-presentation},
on peut trouver $X_0, Y_0$ de type fini sur $A$ dont les
changements de base à $K$ donnent $X$ et $Y$.
D'après Limites, lemmes
\ref{limits-lemma-eventually-proper} et \ref{limits-lemma-descend-smooth},
quitte à agrandir $A$, on peut supposer $X_0$ et $Y_0$
lisses et propres sur $A$.
Par le lemme \ref{lemma-descend-rek},
quitte à agrandir $A$, on peut supposer que $M$ est l'image inverse d'un
objet $M_0  \in D_{perf}(\mathcal{O}_{X_0 \times_{\Spec(A)} Y_0})$
qui est le noyau de Fourier-Mukai d'une équivalence relative
de $X_0$ vers $Y_0$ sur $\Spec(A)$.
Ainsi, le système $(S_0/k, X_0 \to S_0, Y_0 \to S_0, M_0)$
est isomorphe à
$(S_{0, i}/k, X_{0, i} \to S_{0, i}, Y_{0, i} \to S_{0, i}, M_{0, i})$
pour un certain $i \in I$.
Comme $S_i = S_{0, i} \times_{\Spec(k)} \Spec(K)$,
on en déduit que (4) est vraie pour $s : \Spec(K) \to S_i$
induit par le morphisme $\Spec(K) \to \Spec(A) \cong S_{0, i}$
provenant de $A \subset K$.
\end{proof}








\section{Dénombrabilité des variétés équivalentes au sens dérivé}
\label{section-countable-derived-equivalent}

\noindent
Dans cette section, nous démontrons un résultat d'Anel et To\"en, voir \cite{AT}.

\begin{definition}
\label{definition-derived-equivalent}
Soit $k$ un corps. Soient $X$ et $Y$ des schémas projectifs lisses sur $k$.
On dit que $X$ et $Y$ sont {\it équivalents au sens dérivé} s'il existe une équivalence $k$-linéaire
exacte
$D_{perf}(\mathcal{O}_X) \to D_{perf}(\mathcal{O}_Y)$.
\end{definition}

\noindent
Voici le résultat.

\begin{theorem}
\label{theorem-countable}
\begin{reference}
Légère amélioration de \cite{AT}
\end{reference}
Soit $K$ un corps algébriquement clos. Soit $\mathbf{X}$ un schéma lisse et propre
sur $K$. Il existe au plus un nombre dénombrable de classes d'isomorphisme
de schémas lisses et propres $\mathbf{Y}$ sur $K$ qui sont
équivalents à $\mathbf{X}$ au sens dérivé.
\end{theorem}

\begin{proof}
Choisissons un ensemble dénombrable $I$ et, pour $i \in I$, des systèmes
$(S_i/K, X_i \to S_i, Y_i \to S_i, M_i)$ satisfaisant aux propriétés
(1), (2), (3) et (4) du lemme \ref{lemma-countable-equivs}.
Fixons $i \in I$ et posons $S = S_i$, $X = X_i$, $Y = Y_i$ et
$M = M_i$. Il suffit manifestement de montrer que
l'ensemble des classes d'isomorphisme des fibres $Y_s$ pour $s \in S(K)$
telles que $X_s \cong \mathbf{X}$ est dénombrable.
C'est ce que nous démontrons dans le paragraphe suivant.

\medskip\noindent
Soit $S$ un schéma de type fini sur $K$, et soient $X \to S$ et $Y \to S$
des morphismes propres et lisses ; soit $M \in D_{perf}(\mathcal{O}_{X \times_S Y})$
le noyau de Fourier-Mukai d'une équivalence relative de $X$
vers $Y$ sur $S$. Nous allons montrer que
l'ensemble des classes d'isomorphisme des fibres $Y_s$ pour $s \in S(K)$
telles que $X_s \cong \mathbf{X}$ est dénombrable.
Par le lemme \ref{lemma-countable-isos}, appliqué
aux familles $\mathbf{X} \times S \to S$ et $X \to S$,
il existe un ensemble dénombrable $I$ et, pour $i \in I$, un couple
$(S_i \to S, h_i)$ possédant les propriétés suivantes :
\begin{enumerate}
\item $S_i \to S$ est un morphisme de type fini, et l'on pose
$X_i = X \times_S S_i$ ;
\item $h_i : \mathbf{X} \times S_i \to X_i$
est un isomorphisme sur $S_i$ ;
\item pour tout point fermé $s \in S(K)$, si $\mathbf{X} \cong X_s$
sur $K = \kappa(s)$, alors $s$ appartient à l'image de $S_i \to S$
pour un certain $i$.
\end{enumerate}
Posons $Y_i = Y \times_S S_i$. Notons
$M_i \in D_{perf}(\mathcal{O}_{X_i \times_{S_i} Y_i})$
l'image inverse de $M$. Par le lemme \ref{lemma-base-change-rek},
$M_i$ est le noyau de Fourier-Mukai d'une équivalence relative de
$X_i$ vers $Y_i$ sur $S_i$. Comme $I$ est dénombrable, la
propriété (3) ramène la démonstration à prouver que
l'ensemble des classes d'isomorphisme des fibres $Y_{i, s}$ pour $s \in S_i(K)$
est dénombrable.
En fait, cet ensemble est fini d'après
le lemme \ref{lemma-no-deformations-better},
ce qui achève la démonstration.
\end{proof}







\begin{multicols}{2}[\section{Autres chapitres}]
\noindent
Préliminaires
\begin{enumerate}
\item \hyperref[introduction-section-phantom]{Introduction}
\item \hyperref[conventions-section-phantom]{Conventions}
\item \hyperref[sets-section-phantom]{Théorie des ensembles}
\item \hyperref[categories-section-phantom]{Catégories}
\item \hyperref[topology-section-phantom]{Topologie}
\item \hyperref[sheaves-section-phantom]{Faisceaux sur les espaces}
\item \hyperref[sites-section-phantom]{Sites et faisceaux}
\item \hyperref[stacks-section-phantom]{Champs}
\item \hyperref[fields-section-phantom]{Corps}
\item \hyperref[algebra-section-phantom]{Algèbre commutative}
\item \hyperref[brauer-section-phantom]{Groupes de Brauer}
\item \hyperref[homology-section-phantom]{Algèbre homologique}
\item \hyperref[derived-section-phantom]{Catégories dérivées}
\item \hyperref[simplicial-section-phantom]{Méthodes simpliciales}
\item \hyperref[more-algebra-section-phantom]{Compléments d'algèbre}
\item \hyperref[smoothing-section-phantom]{Lissage des morphismes d'anneaux}
\item \hyperref[modules-section-phantom]{Faisceaux de modules}
\item \hyperref[sites-modules-section-phantom]{Modules sur les sites}
\item \hyperref[injectives-section-phantom]{Objets injectifs}
\item \hyperref[cohomology-section-phantom]{Cohomologie des faisceaux}
\item \hyperref[sites-cohomology-section-phantom]{Cohomologie sur les sites}
\item \hyperref[dga-section-phantom]{Algèbre différentielle graduée}
\item \hyperref[dpa-section-phantom]{Algèbre à puissances divisées}
\item \hyperref[sdga-section-phantom]{Faisceaux différentiels gradués}
\item \hyperref[hypercovering-section-phantom]{Hyperrecouvrements}
\end{enumerate}
Schémas
\begin{enumerate}
\setcounter{enumi}{25}
\item \hyperref[schemes-section-phantom]{Schémas}
\item \hyperref[constructions-section-phantom]{Constructions de schémas}
\item \hyperref[properties-section-phantom]{Propriétés des schémas}
\item \hyperref[morphisms-section-phantom]{Morphismes de schémas}
\item \hyperref[coherent-section-phantom]{Cohomologie des schémas}
\item \hyperref[divisors-section-phantom]{Diviseurs}
\item \hyperref[limits-section-phantom]{Limites de schémas}
\item \hyperref[varieties-section-phantom]{Variétés}
\item \hyperref[topologies-section-phantom]{Topologies sur les schémas}
\item \hyperref[descent-section-phantom]{Descente}
\item \hyperref[perfect-section-phantom]{Catégories dérivées de schémas}
\item \hyperref[more-morphisms-section-phantom]{Compléments sur les morphismes}
\item \hyperref[flat-section-phantom]{Compléments sur la platitude}
\item \hyperref[groupoids-section-phantom]{Schémas en groupoïdes}
\item \hyperref[more-groupoids-section-phantom]{Compléments sur les schémas en groupoïdes}
\item \hyperref[etale-section-phantom]{Morphismes étales de schémas}
\end{enumerate}
Thèmes de théorie des schémas
\begin{enumerate}
\setcounter{enumi}{41}
\item \hyperref[chow-section-phantom]{Homologie de Chow}
\item \hyperref[intersection-section-phantom]{Théorie de l'intersection}
\item \hyperref[pic-section-phantom]{Schémas de Picard des courbes}
\item \hyperref[weil-section-phantom]{Théories cohomologiques de Weil}
\item \hyperref[adequate-section-phantom]{Modules adéquats}
\item \hyperref[dualizing-section-phantom]{Complexes dualisants}
\item \hyperref[duality-section-phantom]{Dualité pour les schémas}
\item \hyperref[discriminant-section-phantom]{Discriminants et différentes}
\item \hyperref[derham-section-phantom]{Cohomologie de de Rham}
\item \hyperref[local-cohomology-section-phantom]{Cohomologie locale}
\item \hyperref[algebraization-section-phantom]{Géométrie algébrique et formelle}
\item \hyperref[curves-section-phantom]{Courbes algébriques}
\item \hyperref[resolve-section-phantom]{Résolution des surfaces}
\item \hyperref[models-section-phantom]{Réduction semi-stable}
\item \hyperref[functors-section-phantom]{Foncteurs et morphismes}
\item \hyperref[equiv-section-phantom]{Catégories dérivées de variétés}
\item \hyperref[pione-section-phantom]{Groupes fondamentaux de schémas}
\item \hyperref[etale-cohomology-section-phantom]{Cohomologie étale}
\item \hyperref[crystalline-section-phantom]{Cohomologie cristalline}
\item \hyperref[proetale-section-phantom]{Cohomologie pro-étale}
\item \hyperref[relative-cycles-section-phantom]{Cycles relatifs}
\item \hyperref[more-etale-section-phantom]{Compléments de cohomologie étale}
\item \hyperref[trace-section-phantom]{La formule des traces}
\end{enumerate}
Espaces algébriques
\begin{enumerate}
\setcounter{enumi}{64}
\item \hyperref[spaces-section-phantom]{Espaces algébriques}
\item \hyperref[spaces-properties-section-phantom]{Propriétés des espaces algébriques}
\item \hyperref[spaces-morphisms-section-phantom]{Morphismes d'espaces algébriques}
\item \hyperref[decent-spaces-section-phantom]{Espaces algébriques décents}
\item \hyperref[spaces-cohomology-section-phantom]{Cohomologie des espaces algébriques}
\item \hyperref[spaces-limits-section-phantom]{Limites d'espaces algébriques}
\item \hyperref[spaces-divisors-section-phantom]{Diviseurs sur les espaces algébriques}
\item \hyperref[spaces-over-fields-section-phantom]{Espaces algébriques sur des corps}
\item \hyperref[spaces-topologies-section-phantom]{Topologies sur les espaces algébriques}
\item \hyperref[spaces-descent-section-phantom]{Descente et espaces algébriques}
\item \hyperref[spaces-perfect-section-phantom]{Catégories dérivées d'espaces}
\item \hyperref[spaces-more-morphisms-section-phantom]{Compléments sur les morphismes d'espaces}
\item \hyperref[spaces-flat-section-phantom]{Platitude sur les espaces algébriques}
\item \hyperref[spaces-groupoids-section-phantom]{Groupoïdes dans les espaces algébriques}
\item \hyperref[spaces-more-groupoids-section-phantom]{Compléments sur les groupoïdes dans les espaces}
\item \hyperref[bootstrap-section-phantom]{Amorçage}
\item \hyperref[spaces-pushouts-section-phantom]{Sommes amalgamées d'espaces algébriques}
\end{enumerate}
Thèmes de géométrie
\begin{enumerate}
\setcounter{enumi}{81}
\item \hyperref[spaces-chow-section-phantom]{Groupes de Chow des espaces}
\item \hyperref[groupoids-quotients-section-phantom]{Quotients de groupoïdes}
\item \hyperref[spaces-more-cohomology-section-phantom]{Compléments sur la cohomologie des espaces}
\item \hyperref[spaces-simplicial-section-phantom]{Espaces simpliciaux}
\item \hyperref[spaces-duality-section-phantom]{Dualité pour les espaces}
\item \hyperref[formal-spaces-section-phantom]{Espaces algébriques formels}
\item \hyperref[restricted-section-phantom]{Algébrisation des espaces formels}
\item \hyperref[spaces-resolve-section-phantom]{Résolution des surfaces revisitée}
\end{enumerate}
Théorie des déformations
\begin{enumerate}
\setcounter{enumi}{89}
\item \hyperref[formal-defos-section-phantom]{Théorie formelle des déformations}
\item \hyperref[defos-section-phantom]{Théorie des déformations}
\item \hyperref[cotangent-section-phantom]{Le complexe cotangent}
\item \hyperref[examples-defos-section-phantom]{Problèmes de déformation}
\end{enumerate}
Champs algébriques
\begin{enumerate}
\setcounter{enumi}{93}
\item \hyperref[algebraic-section-phantom]{Champs algébriques}
\item \hyperref[examples-stacks-section-phantom]{Exemples de champs}
\item \hyperref[stacks-sheaves-section-phantom]{Faisceaux sur les champs algébriques}
\item \hyperref[criteria-section-phantom]{Critères de représentabilité}
\item \hyperref[artin-section-phantom]{Axiomes d'Artin}
\item \hyperref[quot-section-phantom]{Espaces de Quot et de Hilbert}
\item \hyperref[stacks-properties-section-phantom]{Propriétés des champs algébriques}
\item \hyperref[stacks-morphisms-section-phantom]{Morphismes de champs algébriques}
\item \hyperref[stacks-limits-section-phantom]{Limites de champs algébriques}
\item \hyperref[stacks-cohomology-section-phantom]{Cohomologie des champs algébriques}
\item \hyperref[stacks-perfect-section-phantom]{Catégories dérivées de champs}
\item \hyperref[stacks-introduction-section-phantom]{Introduction aux champs algébriques}
\item \hyperref[stacks-more-morphisms-section-phantom]{Compléments sur les morphismes de champs}
\item \hyperref[stacks-geometry-section-phantom]{La géométrie des champs}
\end{enumerate}
Thèmes de théorie des espaces de modules
\begin{enumerate}
\setcounter{enumi}{107}
\item \hyperref[moduli-section-phantom]{Champs de modules}
\item \hyperref[moduli-curves-section-phantom]{Espaces de modules de courbes}
\end{enumerate}
Divers
\begin{enumerate}
\setcounter{enumi}{109}
\item \hyperref[examples-section-phantom]{Exemples}
\item \hyperref[exercises-section-phantom]{Exercices}
\item \hyperref[guide-section-phantom]{Guide bibliographique}
\item \hyperref[desirables-section-phantom]{Desiderata}
\item \hyperref[coding-section-phantom]{Style de programmation}
\item \hyperref[obsolete-section-phantom]{Obsolète}
\item \hyperref[fdl-section-phantom]{Licence de documentation libre GNU}
\item \hyperref[index-section-phantom]{Index généré automatiquement}
\end{enumerate}
\end{multicols}


\bibliography{my}
\bibliographystyle{amsalpha}

\end{document}

